% !TeX root = ../Book4.tex
% 纲要标题：# 第四编　导出范畴与现代同调语言
% 纲要位置：计划纲要.md，第 2890 行。

\BookPart{导出范畴与现代同调语言}{b4:part:04}
\BookPartGuide
前面用消解计算导出函子时，曾多次得到“更换消解不影响答案”的结论。本编进一步问：能否把消解和原对象放在同一个范畴中，使这些比较成为真正的同构，同时保留不同次数之间的映射？导出范畴回答这一问题。其态射不仅包含普通同态，也包含扩张类；好三角则把短正合列及其连接映射一起保留下来。

\ProseParagraphBreak
\chapref{b4:ch:16}先从复形、链同伦和平移出发建立同伦范畴。映射锥给出好三角，显式矩阵证明旋转、态射补全和八面体公理。随后，\chapref{b4:ch:17}把拟同构倒转，用屋顶演算构造导出范畴，并说明三角结构为何仍然成立。读这两章主要需要第二编的复形、锥和消解。学习\secref{b4:sec:1704}的有界投射与内射模型时，还需\subsecref{b4:subsec:1502:01}、\subsecref{b4:subsec:1502:02}的 Cartan--Eilenberg 消解定义、存在性与比较，以及\cref{b4:cor:double-complex-comparison}的双复形比较结论。

\ProseParagraphBreak
\chapref{b4:ch:18}把对象级导出函子提升为复形级函子。全右、全左导出的自然变换方向分别由内射和投射替换决定；导出张量与导出 Hom 的构造则把 Tor、Ext、伴随和基变换联系起来。前三节在有界方向下证明所需结论，\secref{b4:sec:1804}则构造无界半自由替换。

\ProseParagraphBreak
复合导出的问题在这里也有了复形层面的表述。\secref{b4:sec:1004}先考察前函子或后函子正合时，怎样简化 \(R^n(GF)(M)\)；\secref{b4:sec:1503}在两个函子都只有左正合性时，用 \(R^pG(R^qF(M))\) 组成谱序列，保留两次导出的信息。\cref{b4:thm:total-derived-composition}则在 \(D^+\) 上构造自然比较 \(R(GF)\to RG\circ RF\)。比较的存在不保证它可逆；当 \(F\) 把内射对象送到 \(G\)-零调对象时，比较成为同构，这也正是上述 Grothendieck 谱序列的附加条件。即便如此，谱序列的极限页给出的仍是目标上同调的滤过相邻商，不自动给出直和分解。

\ProseParagraphBreak
\chapref{b4:ch:19}讨论另一类问题：一个三角范畴能否重新识别出对象、短正合列和上同调？\secref{b4:sec:1901}至\secref{b4:sec:1903}从标准截断走到抽象 \(t\) 结构及其交换范畴的心。\secref{b4:sec:1904}至\secref{b4:sec:1906}进一步用扭对改变心，并在两顶点线性映射范畴中完整计算非标准心，比较心内高阶扩张与环境中的移位态射；这一计算不能由“存在另一颗心”的形式结论代替。本章使用\chapref{b4:ch:16}至\chapref{b4:ch:17}的三角与消解理论。

\ProseParagraphBreak
\chapref{b4:ch:20}介绍 DG 语言、完美复形和倾斜理论。\subsecref{b4:subsec:2002:01}、\subsecref{b4:subsec:2002:02}用有界有限生成投射复形建立完美模型与有限厚生成；三顶点代数中自身扩张、自同态及正则模生成的计算也由有界理论完成。判定倾斜对象还要用\subsecref{b4:subsec:2002:03}证明的完美复形厚性；该小节中任意余积和紧对象的反向刻画依赖\secref{b4:sec:1804}的无界半自由替换。紧对象刻画及一般导出 Morita 定理的后半论证给出证明概要，说明无界构造、有限因子分解及双模替换的作用。

\ProseParagraphBreak
本编始终采用上链微分和 \(C[1]^n=C^{n+1}\)，固定锥为 \(D^n\oplus C^{n+1}\)。阅读不同文献时，须一并核对平移方向、旋转负号以及左右模作用。主要参考为 Weibel\cite[Sections 10.1--10.8]{Weibel1994HomologicalAlgebra}、Beilinson--Bernstein--Deligne\cite[Section 1.3]{BBD1982FaisceauxPervers}及 Keller\cite[Sections 3.1 and 4]{Keller1998AbelianDerived}\cite[Sections 5.3, 8--9]{Keller1994DerivingDG}；各章还列出具体的参考章节或定理。

% !TeX root = ../../Book4.tex
% 纲要标题：## 第17章　导出范畴的构造
% 纲要位置：计划纲要.md，第 2921 行。

\chapter{导出范畴的构造}
\label{b4:ch:17}
\BookChapterNumber{b4:ch:17}

\begin{chapterabstract}
导出范畴把拟同构变为可逆态射，使消解与原对象真正同构。本章从局部化的泛性质出发，证明拟同构允许屋顶演算，并核对三角结构如何下降。有界方向的截断使有界导出范畴成为满子范畴；投射和内射复形进一步把导出态射化为链同伦类，得到 Ext 的平移态射解释。
\end{chapterabstract}

\begin{learninggoals}
\begin{itemize}
\item 用泛性质区别导出范畴与仅列出上同调对象的做法。
\item 通分、合成和比较屋顶，说明拟同构的消去条件。
\item 证明导出范畴的三角结构及短正合列所产生的好三角。
\item 区分项有界、上同调有界和特殊消解的有界方向。
\item 用投射或内射模型计算导出态射，并核对 Ext 的次数与符号。
\end{itemize}
\end{learninggoals}

本章使用\chapref{b4:ch:16}的三角公理和 Hom 正合性。\secref{b4:sec:1704}的模型另外调用\chapref{b4:ch:07}的投射、内射消解，\chapref{b4:ch:08}的 Ext，\cref{b4:thm:ce-resolution}中有下界复形的 Cartan--Eilenberg 消解，以及\cref{b4:cor:double-complex-comparison}的双复形比较。第十五章的消解构造独立于其谱序列应用。前三节建立导出范畴及其三角结构；第17.4节进一步给出计算态射所需的有界模型。

\ProseParagraphBreak
局部化与屋顶可对照 Weibel\cite[Section 10.3; Proposition 10.4.1]{Weibel1994HomologicalAlgebra}；有界消解模型见该书\cite[Corollary 10.4.7; Theorem 10.4.8]{Weibel1994HomologicalAlgebra}。本章同时说明一般交换范畴的集合大小约定，并在模范畴中给出局部小性的具体证明。无界消解的存在不作为这一章的隐含假设。

% !TeX root = ../../Book4.tex
% 纲要标题：### 17.1 局部化
% 纲要位置：计划纲要.md，第 2923 行。

\section{局部化}
\label{b4:sec:1701}

拟同构保存同调，却未必已有同伦逆。将全部拟同构正式变为可逆态射，能允许更换消解，同时仍保留态射的复合及其泛性质。


% 纲要内容（保留纲要核验项目）：
%
% * 把拟同构变成可逆态射
% * 局部化的泛性质
% * 集合大小与构造假设
%

\subsection{拟同构的可逆化}
\label{b4:subsec:1701:01}

链同伦等价必为拟同构，但反向不成立。消解 \(P\to M\) 的作用恰在于：虽然它可能不是同伦等价，上同调信息却与 \(M\) 完全相同。若希望既允许换消解又保持态射的复合，就应把所有拟同构都当作可逆态射。

\begin{definition}\label{b4:def:derived-category}
设 \(\mathcal A\) 为局部小交换范畴，\(S\) 为 \(K(\mathcal A)\) 中的拟同构，即在各次上同调上诱导同构的态射类。在适当的集合大小约定下，取局部化
\[
 D(\mathcal A)=K(\mathcal A)[S^{-1}].
\]
所得范畴称为 \(\mathcal A\) 的\term[daochu-fanchou]{导出范畴}[derived category]，其局部化函子记为 \(Q:K(\mathcal A)\to D(\mathcal A)\)。这里的逆是新范畴中的逆，不要求拟同构原本具有复形映射意义的逆。
\end{definition}
\symidx{\(D(\mathcal A)\)：复形的导出范畴}

一个零调复形在 \(D(\mathcal A)\) 中成为零，因为它到零复形的映射是拟同构。反过来，各次上同调函子能经 \(Q\) 分解，所以在 \(D(\mathcal A)\) 中为零的对象必零调。由此看出，导出范畴的零对象比同伦范畴的零对象多。

\subsection{局部化的泛性质}
\label{b4:subsec:1701:02}

\begin{definition}\label{b4:def:category-localization}
设 \(\mathcal C\) 为范畴，\(S\) 为一类态射，给定函子 \(Q:\mathcal C\to\mathcal C[S^{-1}]\)。若 \(Q(s)\) 对每个 \(s\in S\) 都可逆，并且对每个范畴 \(\mathcal E\)，任何将 \(S\) 送为同构的函子 \(F:\mathcal C\to\mathcal E\) 都唯一写成 \(F=\overline FQ\)，自然变换也唯一经此分解，则称 \(Q\) 为 \(\mathcal C\) 在 \(S\) 处的\term[fanchou-jubuhua]{范畴局部化}[localization of a category]。采用对象不变的标准构造时，这里的唯一性是严格的；更换等价模型后则为唯一到自然同构。
\end{definition}

小范畴的一个直接构造，是在原有箭头之外为每个 \(s\in S\) 添一个反向符号 \(s^{-1}\)，取有限可复合箭头串，再按原范畴的复合关系以及 \(ss^{-1}=1,s^{-1}s=1\) 取商。复合由连接箭头串给出。函子 \(F\) 在新符号上只能取 \(F(s)^{-1}\)，因而唯一延拓；自然变换在新符号上的自然性由在 \(s\) 上的自然性乘逆得到。这也证明了自然变换的分解性质。

\ProseParagraphBreak
这种构造证明存在，却不便计算。下一节证明：同伦范畴中的拟同构允许把任意箭头串缩成一个分母在左的屋顶。这里不能先假定任意范畴都有这样的分式形式；它需要专门的交换与消去条件。

\ProseParagraphBreak
还可直接从复形范畴把拟同构倒过来。两种做法等价：拟同构包含链同伦等价，故\cref{b4:prop:homotopy-universal}迫使任何这样的函子先经 \(K(\mathcal A)\) 分解；之后再用局部化的泛性质即可。

\subsection{集合大小与构造假设}
\label{b4:subsec:1701:03}

若 \(\mathcal A\) 的对象不是一个集合，前述箭头串也可能形成真类，不能仅凭原范畴局部小就断言局部化仍局部小。本书对一般交换范畴使用固定的较大集合范围来作局部化；讨论 Hom 群时，要么假定局部化在所用范围内局部小，要么使用下面模范畴的大小控制，或\secref{b4:sec:1704}的有界消解模型。这里不把无界的任意交换范畴情形一概当作已解决。

\begin{proposition}\label{b4:prop:module-derived-size}
设 \(R\) 为含幺环，\(X\) 为左 \(R\)-模复形。存在一集合的拟同构 \(X_i\to X\)，使每个拟同构 \(s:Y\to X\) 都能由某个 \(X_i\to X\) 经拟同构 \(X_i\to Y\) 细化。因此在允许屋顶演算后，\(D(R\text{-}\mathbf{Mod})\) 的每个态射类确实组成集合。
\end{proposition}
\begin{proof}
取无限基数 \(\kappa\) 至少为 \(|R|\)、\(\aleph_0\) 及所有 \(|X^n|\) 的上界。给定 \(s\)，为 \(X\) 每次上同调的每个类选一个 \(Y\) 中余圈，其像代表该类。所选元素生成一个各次大小至多 \(\kappa\) 的子复形 \(Y_0\subseteq Y\)。

归纳构造 \(Y_j\)：对 \(Y_j\) 中每个余圈 \(z\)，若 \(s(z)\) 在 \(X\) 中为余边，则 \(s\) 的上同调单射性保证 \(z=\mathrm{d}y\) 对某个 \(y\in Y\) 成立。把这些 \(y\) 加入并生成子复形，得 \(Y_{j+1}\)。每一步至多加入 \(\kappa\) 个元素；取可数并 \(Y'=\bigcup_jY_j\)，其大小仍至多 \(\kappa\)。起初的选择保证 \(H(Y')\to H(X)\) 满，每个落入核的余圈会在下一步成为余边，所以该映射也单。于是 \(Y'\to X\) 和 \(Y'\to Y\) 都是拟同构。

所有大小至多 \(\kappa\) 的模复形可以在固定集合上编码，它们到 \(X\) 的映射也组成集合。取其代表即得所需的一集合分母。每个屋顶 \(X\leftarrow Y\to Z\) 都可换成以某个 \(X_i\) 为中间对象的屋顶，而 \(\bigcup_i\operatorname{Hom}_{K}(X_i,Z)\) 为集合，其商仍为集合。
\end{proof}

这里用 \(\le\kappa\) 控制大小，避免了“可数次合并以后仍严格小于某个基数”的不当推断。相关大小问题可对照 Weibel\cite[Set-Theoretic Considerations 10.3.6; Proposition 10.4.4]{Weibel1994HomologicalAlgebra}。

% !TeX root = ../../Book4.tex
% 纲要标题：### 17.2 三角范畴公理
% 纲要标签：b4:sec:1602
% 纲要位置：计划纲要.md，第 3661 行。

\section{三角范畴公理}
\label{b4:sec:1602}

映射锥的旋转和复合规律可以独立表述为三角范畴公理。核对八面体时要保留各箭头的符号，才能将抽象相容性落实到复形公式。


% 纲要内容（保留纲要核验项目）：
%
% * 好三角与旋转
% * 态射公理与八面体公理
% * 三角范畴不是一般交换范畴
% * 同伦范畴的三角结构通过复形层面的显式锥构造验证；八面体公理的矩阵映射与符号单独核对，不把其成立仅列作形式约定
%

\subsection{好三角与旋转}
\label{b4:subsec:1602:01}

对可复合的 \(X\xrightarrow{u}Y\xrightarrow{v}Z\)，锥 \(C_u\) 记录第一步的差别，\(C_v\) 记录第二步的差别，而 \(C_{vu}\) 记录整个复合的差别。八面体公理把三者联系为一条好三角
\[
 C_u\longrightarrow C_{vu}\longrightarrow C_v\longrightarrow C_u[1].
\]
三角范畴将这种联系与锥的补全、旋转及态射补全一并作为公理，因此不必要求对象本身仍是复形。

\begin{definition}\label{b4:def:triangulated-category}
设 \(\mathcal T\) 为加性范畴，给定自等价 \([1]\) 及其逆，并固定它们的迭代。再指定一类称为好三角的图式 \(X\xrightarrow{u}Y\xrightarrow{v}Z\xrightarrow{w}X[1]\)。若这些数据满足下面四条公理，则称 \(\mathcal T\) 连同这些数据为\term[sanjiao-fanchou]{三角范畴}[triangulated category]。
\begin{enumerate}
\item[(TR1)] 好三角对同构封闭；\(X\xrightarrow{1}X\to0\to X[1]\) 是好三角；每个态射 \(u\) 都能补成好三角。
\item[(TR2)] 一个三角为好三角，当且仅当其旋转
\(Y\xrightarrow{v}Z\xrightarrow{w}X[1]\xrightarrow{-u[1]}Y[1]\) 为好三角。
\item[(TR3)] 两个好三角的前两项之间若已有交换方块，则可以补上第三项的态射，成为三角态射。
\item[(TR4)] 对可复合的 \(X\xrightarrow{u}Y\xrightarrow{v}Z\)，取分别补全 \(u,v,vu\) 的好三角，其后两条箭头记为 \(i_u,p_u\)、\(i_v,p_v\)、\(i_{vu},p_{vu}\)，第三对象记为 \(C_u,C_v,C_{vu}\)。存在 \(\alpha:C_u\to C_{vu}\)、\(\beta:C_{vu}\to C_v\)，满足
\[
\alpha i_u=i_{vu}v,\quad p_{vu}\alpha=p_u,\quad
\beta i_{vu}=i_v,\quad p_v\beta=u[1]p_{vu},
\]
并使 \(C_u\xrightarrow{\alpha}C_{vu}\xrightarrow{\beta}C_v\xrightarrow{i_u[1]p_v}C_u[1]\) 为好三角。
\end{enumerate}
\end{definition}

旋转中的负号不能任意删去。它与平移复形的负微分一起，使同一条长正合列在换起点后仍由三角产生。连续旋转三次所得是三个箭头都带负号的平移三角；在三个对象上分别取 \(1,-1,1\) 作同构，可将它改写为箭头 \(u[1],v[1],-w[1]\) 的好三角。这里最后一条箭头仍带负号，不能把三个符号一起删去。

\ProseParagraphBreak
以下先从前三条公理推导 Hom 正合性及其应用，最后用锥的具体计算验证 \(K(\mathcal A)\) 满足全部公理。这样，在验证八面体时调用的正合性只依赖已经核对的前三条公理。

\subsection{态射公理与八面体公理}
\label{b4:subsec:1602:02}

(TR3) 只要求补全存在，没有要求唯一。(TR4) 则规定这种补全对于复合态射还应满足什么相容关系。把三个三角写成
\[
 X\xrightarrow{u}Y\xrightarrow{i_u}C_u\xrightarrow{p_u}X[1],\quad
 Y\xrightarrow{v}Z\xrightarrow{i_v}C_v\xrightarrow{p_v}Y[1],
\]
\[
 X\xrightarrow{vu}Z\xrightarrow{i_{vu}}C_{vu}\xrightarrow{p_{vu}}X[1].
\]
所要求的是 \(\alpha:C_u\to C_{vu}\)、\(\beta:C_{vu}\to C_v\)，满足
\[
 \alpha i_u=i_{vu}v,\quad p_{vu}\alpha=p_u,\quad
 \beta i_{vu}=i_v,\quad p_v\beta=u[1]p_{vu},
\]
并使
\[
 C_u\xrightarrow{\alpha}C_{vu}\xrightarrow{\beta}C_v
 \xrightarrow{i_u[1]p_v}C_u[1]
\]
为好三角。这一形式称为\term[bamian-ti-gongli]{八面体公理}[octahedral axiom]，名称来自把上述四个三角画成八面体的四个面。把六个未平移对象放在同一张平面图中，并在底端补出连接箭头的目标，就能同时追踪四个三角：
\[
\begin{tikzcd}[column sep=1.2em,row sep=2.1em]
&&X\arrow[dl,"u"']\arrow[dr,"vu"]&&\\
&Y\arrow[rr,"v"]\arrow[dl,"i_u"']&&Z\arrow[dl,"i_{vu}"']\arrow[dr,"i_v"]&\\
C_u\arrow[rr,"\alpha"]\arrow[d,"p_u"']&&C_{vu}\arrow[rr,"\beta"]\arrow[dll,"p_{vu}"]&&C_v\arrow[d,"p_v"]\arrow[dll,"\gamma"]\\
X[1]\arrow[rrrr,bend right=25,"{u[1]}"']&&C_u[1]&&Y[1]\arrow[ll,"{i_u[1]}"] .
\end{tikzcd}
\]
这里 \(\gamma=i_u[1]p_v\)。顶部的两条路径比较 \(v\circ u\)；中间两条新箭头把 \(u,vu,v\) 的三个锥连成第四个三角。底端平移对象上的箭头记录各三角的连接态射，图中的交换关系正是前面四个等式。后面的锥矩阵计算逐一验证这些路径，以及第四个三角为何仍为好三角。

\begin{lemma}\label{b4:lem:pretriangle-hom}
设局部小加性范畴 \(\mathcal T\) 带平移及满足 (TR1)--(TR3) 的好三角。对好三角 \(X\xrightarrow{u}Y\xrightarrow{v}Z\xrightarrow{w}X[1]\)，相邻箭头复合为零，且对任意对象 \(W\in\mathcal T\)，两个 Hom 序列
\[
 \operatorname{Hom}(W,X)\longrightarrow\operatorname{Hom}(W,Y)
 \longrightarrow\operatorname{Hom}(W,Z),
\]
\[
 \operatorname{Hom}(Z,W)\longrightarrow\operatorname{Hom}(Y,W)
 \longrightarrow\operatorname{Hom}(X,W)
\]
均在中间正合。三角态射中任意两个分量为同构，则第三个也为同构。
\end{lemma}
\begin{proof}
从恒等三角到给定三角应用 (TR3)，前两个分量取 \(1_X,u\)，第三个分量由零对象出发，得到 \(vu=0\)。旋转后同理得到其余零复合。

若 \(a:W\to Y\) 满足 \(va=0\)，要证明的是存在 \(b_0:W\to X\)，使 \(ub_0=a\)。(TR3) 从前两项的交换方块补出第三项的态射；因此取第二项为零的好三角，使方块的交换条件恰好是 \(va=0\)，并让补出的态射经逆平移成为所需的 \(b_0\)。具体地，对上方好三角
\(W\to0\to W[1]\xrightarrow{-1}W[1]\)
和下方旋转三角
\(Y\xrightarrow{v}Z\xrightarrow{w}X[1]\xrightarrow{-u[1]}Y[1]\)
使用 (TR3)，前两个分量取 \(a,0\)。所得第三个分量 \(b:W[1]\to X[1]\) 满足 \(u[1]b=a[1]\)，所以 \(a=ub[-1]\)。这给出第一列的反向包含。

若 \(a:Y\to W\) 满足 \(au=0\)，恒等三角的逆旋转给出好三角 \(0\to W\xrightarrow{1}W\to0\)。从给定三角到这个三角应用 (TR3)，前两个分量取 \(0:X\to0\) 和 \(a:Y\to W\)；前方块由 \(au=0\) 交换，补出的第三个分量 \(b:Z\to W\) 满足 \(bv=a\)。结合 \(vu=0\)，这就证明了反变 Hom 序列在中间正合。

把两类正合列旋转延长。若一个三角态射的前两个分量 \(a:X\to X'\)、\(b:Y\to Y'\) 同构，对每个 \(W\) 比较两个三角产生的连续五项
\[
\begin{gathered}
 \operatorname{Hom}(W,X)\longrightarrow\operatorname{Hom}(W,Y)
 \longrightarrow\operatorname{Hom}(W,Z)\longrightarrow{}\\
 \operatorname{Hom}(W,X[1])\longrightarrow\operatorname{Hom}(W,Y[1]).
\end{gathered}
\]
它们在中间三项正合，五个比较映射依次由 \(a,b,c,a[1],b[1]\) 诱导。除中央由 \(c:Z\to Z'\) 诱导的映射外，其余四个都是同构，所以\cref{b4:thm:five-lemma}给出 \(\operatorname{Hom}(W,c)\) 为同构。取 \(W=Z'\)，由满射性得到 \(s:Z'\to Z\) 满足 \(cs=1_{Z'}\)；再取 \(W=Z\)，由单射性和 \(csc=c\) 得 \(sc=1_Z\)。故 \(c\) 本身同构。其他两种情形由旋转归约。
\end{proof}

因此，同一态射的两个好三角可以用 (TR3) 补出同构；但补出的同构是否唯一，是另一回事。

\subsection{三角范畴与交换范畴的区别}
\label{b4:subsec:1602:03}

交换范畴用核和余核描述一个态射，三角范畴则用第三个对象同时保留它前后次数的信息。不能把这个对象简单称作余核：即使 \(u=0\)，它也通常是 \(Y\oplus X[1]\)，而不是 \(Y\)。

\begin{proposition}\label{b4:prop:triangle-mono-splits}
设 \(\mathcal T\) 是三角范畴。\(\mathcal T\) 中每个范畴意义的单态射均分裂，每个范畴意义的满态射也均分裂。
\end{proposition}
\begin{proof}
若 \(u:X\to Y\) 单，把它补成好三角 \(X\to Y\to Z\xrightarrow{w}X[1]\)。由 \(u[1]w=0\) 和 \(u[1]\) 单得 \(w=0\)。旋转并应用\cref{b4:lem:pretriangle-hom}的反变正合列，得到
\[
 \operatorname{Hom}(Y,X)\longrightarrow\operatorname{Hom}(X,X)
 \xrightarrow{(-w[-1])^*}\operatorname{Hom}(Z[-1],X).
\]
最后的映射为零，故 \(1_X\) 提升为 \(r:Y\to X\)，即 \(ru=1_X\)。满态射的论证对偶。
\end{proof}

在交换群范畴中，单态射 \(\mathbb Z\xrightarrow{2}\mathbb Z\) 没有左逆。三角范畴却能将每个态射补成好三角，旋转和 Hom 正合性因而对单态射施加了额外的分裂约束。后面将用同一个乘二映射说明：逐次单射的复形映射，在 \(K(\mathbf{Ab})\) 中未必是单态射。

\subsection{显式锥构造与八面体的符号}
\label{b4:subsec:1602:04}

\begin{theorem}\label{b4:thm:homotopy-triangulated}
设 \(\mathcal A\) 为局部小加性范畴。以平移 \([1]\) 和\cref{b4:def:good-cone-triangle}的好三角，\(K(\mathcal A)\) 是三角范畴；\(K^b(\mathcal A),K^+(\mathcal A),K^-(\mathcal A)\) 也各自为三角范畴。
\end{theorem}
\begin{proof}
好三角按定义是与标准锥三角同构的三角，所以对同构封闭。\(\operatorname{Cone}(1_A)\) 可缩，在 \(K(\mathcal A)\) 中是零对象；因而恒等态射的标准锥三角与 \(A\xrightarrow{1}A\to0\to A[1]\) 同构。对任意 \(u\in\operatorname{Hom}_{K(\mathcal A)}(A,B)\)，选择代表它的复形映射 \(f:A\to B\)，其标准锥三角的首箭头就是 \(u=[f]\)，从而补全了 \(u\)。这验证了 (TR1) 的三个要求。

为核对 (TR2)，写 \(C_f=B\oplus A[1]\)。包含 \(i_f:B\to C_f\) 的锥在次数 \(n\) 为
\[
 B^n\oplus A^{n+1}\oplus B^{n+1},\qquad
 \mathrm{d}(b,a,b')=(\mathrm{d}b+fa+b',-\mathrm{d}a,-\mathrm{d}b').
\]
映射 \(r(b,a,b')=a\) 和 \(s(a)=(0,a,-fa)\) 给出它与 \(A[1]\) 的同伦等价。确实 \(rs=1\)，而 \(1-sr=\mathrm{d}h+h\mathrm{d}\)，其中 \(h(b,a,b')=(0,0,b)\)。记标准锥三角中的包含为 \(j:C_f\to\operatorname{Cone}(i_f)\)，最后投影为 \(q:\operatorname{Cone}(i_f)\to B[1]\)，则 \(rj=p_f\)、\(qs=-f[1]\) 严格成立，并且在 \(K(\mathcal A)\) 中
\[
 (-f[1])r=qsr=q,
\]
因为 \(sr\simeq1\)。因此 \((1_B,1_{C_f},r)\) 的三个方块都交换，它是标准锥三角到旋转三角的同构。

锥的平移同构
\[
 \operatorname{Cone}(f[r])\xrightarrow{\sim}\operatorname{Cone}(f)[r],
 \qquad (b,a)\longmapsto(b,(-1)^ra)
\]
与包含相容，把最后投影变为 \((-1)^rp_f[r]\)。所以好三角 \(A\xrightarrow{u}B\xrightarrow{v}C\xrightarrow{w}A[1]\) 平移到 \(r=-1\) 后，箭头为 \(u[-1],v[-1],-w[-1]\) 的三角仍为好三角。将它旋转两次，所得箭头为 \(-w[-1],-u,-v\)；在第三个对象上取负恒等同构，就得到逆向旋转
\[
 C[-1]\xrightarrow{-w[-1]}A\xrightarrow{u}B\xrightarrow{v}C.
\]
由此正向和逆向旋转都保持好三角，(TR2) 的两向均成立。

为核对 (TR3)，先取两个标准锥三角，设前两个分量为 \(a:A\to A'\)、\(b:B\to B'\)。同伦交换表示可以选择 \(H\) 使
\(bf-f'a=\mathrm{d}_{B'}H+H\mathrm{d}_A\)。要与包含、投影相容，第三个分量可取上三角形式
\[
 c=\begin{pmatrix}b&T\\0&a[1]\end{pmatrix}:C_f\longrightarrow C_{f'},
 \qquad T^n:A^{n+1}\longrightarrow B'^n.
\]
将它代入复形映射条件 \(\mathrm{d}_{C_{f'}}c=c\mathrm{d}_{C_f}\)，非对角块的等式为
\[
 \mathrm{d}_{B'}T+T\mathrm{d}_A=bf-f'a.
\]
因此取 \(T^n=H^{n+1}\) 就能用所给同伦修正原方块的误差，即定义
\[
 c:C_f\longrightarrow C_{f'},\qquad c(y,x)=(by+Hx,ax).
\]
第二分量的微分相容性由 \(a\) 为复形映射得到。它与包含和投影严格相容，所以补成三角态射。任意好三角先用同构换成标准锥即可。至此前三条公理成立，因而可用\cref{b4:lem:pretriangle-hom}：对固定首箭头，不同的好三角之间存在保持前两项的同构。

现在核对 (TR4)。设复形映射 \(f:A\to B\)、\(g:B\to C\)。从 \(f\) 到 \(gf\) 的前两项比较是 \((1_A,g)\)，从 \(gf\) 到 \(g\) 的比较是 \((f,1_C)\)；两个方块都严格交换，所以在刚才的公式中取 \(H=0\)，便得到前两张锥映射。第三张先由 \(p_g\) 投影到 \(B[1]\)，再由 \(i_f[1]\) 包含到 \(C_f[1]\)。于是三张映射为
\[
 \alpha(b,a)=(gb,a),\qquad
 \beta(c,a)=(c,fa),\qquad
 \gamma(c,b)=(b,0).
\]
最后一个目标为 \(C_f[1]\)。它们都是复形映射，且满足公理中四个相容等式和 \(\gamma=i_f[1]p_g\)。复合 \(\beta\alpha\) 的零伦为 \((b,a)\mapsto(0,b)\)；要将这三个箭头识别为好三角，还须把 \(\operatorname{Cone}(\alpha)\) 与 \(C_g\) 比较。

把 \(\operatorname{Cone}(\alpha)^n\) 按顺序写作
\(C^n\oplus A^{n+1}\oplus B^{n+1}\oplus A^{n+2}\)，则
\[
 \mathrm{d}(c,a,b,a')=(\mathrm{d}c+gf a+gb,\,-\mathrm{d}a+a',\,-\mathrm{d}b-fa',\,\mathrm{d}a').
\]
其中 \(a,a'\) 两个分量仍含有 \(A\) 的重复信息。令 \(t=b+fa\)，则 \(t\) 的微分分量为
\[
 -\mathrm{d}_Bb-fa'+f(-\mathrm{d}_Aa+a')
 =-\mathrm{d}_B(b+fa)=-\mathrm{d}_Bt.
\]
因此可逆坐标变换
\[
 (c,a,b,a')\longmapsto\bigl((c,b+fa),(a,a')\bigr)
\]
将微分变为两组互不混合的微分
\[
 (c,t)\longmapsto(\mathrm{d}_Cc+gt,-\mathrm{d}_Bt),\qquad
 (a,a')\longmapsto(-\mathrm{d}_Aa+a',\mathrm{d}_Aa').
\]
这正是严格的复形同构
\[
 \operatorname{Cone}(\alpha)\cong C_g\oplus\operatorname{Cone}(1_{A[1]}).
\]
第二个直和因子由\cref{b4:prop:cone-sign-compatibility}可缩，在同伦范畴中为零。取上述分解中第一因子的投影及包含，写回原坐标就是
\[
 r(c,a,b,a')=(c,b+fa),\qquad s(c,b)=(c,0,b,0).
\]
于是 \(rs=1\)、\(sr\simeq1\)，故 \(r,s\) 是同伦逆。若 \(j:C_{gf}\to\operatorname{Cone}(\alpha)\) 是包含，\(q:\operatorname{Cone}(\alpha)\to C_f[1]\) 是最后投影，则 \(rj=\beta\)、\(qs=\gamma\) 严格成立，并且在 \(K(\mathcal A)\) 中
\[
 \gamma r=qsr=q,
\]
因为 \(sr\simeq1\)。于是 \((1_{C_f},1_{C_{gf}},r)\) 的三个方块都交换，标准锥三角经此同构成为
\(C_f\xrightarrow{\alpha}C_{gf}\xrightarrow{\beta}C_g\xrightarrow{\gamma}C_f[1]\)，正是要求的好三角。给定的三个三角若非标准形式，利用前述保持首两项的同构运送所有箭头，四个相容等式随之保留。这完成八面体公理的验证。

所有构造只使用有限双积、平移和锥。它们保持两端有界、有上界或有下界，因此三个有界方向的满子范畴同样满足全部公理。
\end{proof}

由\cref{b4:thm:homotopy-triangulated}，\(K(\mathbf{Ab})\) 是三角范畴，却不是交换范畴。若它是交换范畴，把任意态射 \(f:X\to Y\) 分解为满态射 \(e:X\to I\) 和单态射 \(m:I\to Y\)，则由\cref{b4:prop:triangle-mono-splits}可取 \(s:I\to X\)、\(r:Y\to I\)，满足 \(es=1_I\)、\(rm=1_I\)。于是 \(t=sr:Y\to X\) 满足 \(ftf=f\)。对 \(\mathbb Z[0]\) 的乘二态射，这要求某个整数 \(t\) 满足 \(4t=2\)，矛盾：次数零复形的自同态同伦类就是整数，非零的链同伦无处取值。

\ProseParagraphBreak
尤其要区分“各次数均为单射”和“在同伦范畴中为单态射”。乘二 \(\mathbb Z[0]\to\mathbb Z[0]\) 逐次单，却不是同伦范畴中的单态射。令 \(C=\operatorname{Cone}(2)\)，则 \(C[-1]\to\mathbb Z[0]\) 的投影同伦类非零，而复合乘二后为零伦。这里 \(C[-1]\) 在次数零和一为 \(\mathbb Z\)，微分为 \(-2\)；投影为零伦会要求某个整数 \(t\) 满足 \(-2t=1\)，而投影复合乘二的零伦可取 \(h^1=-1\)。这是更换态射范畴后单态射性质的变化。

\begin{definition}\label{b4:def:triangle-functor}
设 \(\mathcal T,\mathcal U\) 为三角范畴。加性函子 \(F:\mathcal T\to\mathcal U\) 连同自然同构 \(\phi_X:F(X[1])\to F(X)[1]\)，若把每个好三角 \(X\xrightarrow{u}Y\xrightarrow{v}Z\xrightarrow{w}X[1]\) 送到好三角
\[
 F(X)\xrightarrow{F(u)}F(Y)\xrightarrow{F(v)}F(Z)
 \xrightarrow{\phi_XF(w)}F(X)[1],
\]
则称 \((F,\phi)\) 为\term[sanjiao-hanzi]{三角函子}[triangle functor]，也称三角范畴间的正合函子。具有三角准逆的此类函子称为三角等价。
\end{definition}

比较同构 \(\phi\) 要把 \(F(w)\) 的目标 \(F(X[1])\) 识别为 \(F(X)[1]\)，因此它也决定最后一条连接箭头及其符号。例如对平移函子 \(F=[1]\)，两者虽都等于 \(X[2]\)，好三角平移后的箭头却是 \(u[1],v[1],-w[1]\)。取 \(\phi_X=-1_{X[2]}\) 正好给出最后的 \(-w[1]\)；若取恒等比较，得到的将是 \(w[1]\)，这个三角未必仍是好三角。下章的局部化函子则采用与复形平移严格相容的比较。对于复形范畴上的逐项加性函子，有限双积及锥矩阵的保持直接给出三角函子。

% !TeX root = ../../Book4.tex
% 纲要标题：### 17.3 上同调函子
% 纲要标签：b4:sec:1603
% 纲要位置：计划纲要.md，第 3668 行。

\section{上同调函子}
\label{b4:sec:1603}

好三角经过上同调函子以后形成长正合列，Hom 本身也具有这种性质。这些正合列判断三角何时分裂，并说明它与交换范畴短正合列的关系。


% 纲要内容（保留纲要核验项目）：
%
% * 三角产生长正合列
% * Hom 函子的正合性
% * 三角与短正合列的对应范围
%

\subsection{三角与长正合列}
\label{b4:subsec:1603:01}

\begin{definition}\label{b4:def:cohomological-functor}
设 \(\mathcal T\) 为三角范畴，\(\mathcal B\) 为交换范畴。加性函子 \(H:\mathcal T\to\mathcal B\) 若将每个好三角 \(X\to Y\to Z\to X[1]\) 变为在中间正合的序列 \(H(X)\to H(Y)\to H(Z)\)，则称为\term[shangtongdiao-hanzi]{上同调函子}[cohomological functor]。记 \(H^n(X)=H(X[n])\)。
\end{definition}

由旋转和平移，定义中的三项正合自动延长成
\[
 \cdots\to H^n(X)\to H^n(Y)\to H^n(Z)
 \to H^{n+1}(X)\to H^{n+1}(Y)\to\cdots.
\]
旋转可能给某些箭头添负号，但不改变核和像；选定一致的符号后，这就是长正合列。三角态射诱导长正合列之间的交换图，所以正合列同时记录对象和态射的信息。

\begin{proposition}\label{b4:prop:cohomology-on-homotopy}
设 \(\mathcal A\) 为局部小交换范畴。通常的零次上同调 \(H^0:K(\mathcal A)\to\mathcal A\) 是上同调函子，其平移得到通常的各次上同调。
\end{proposition}
\begin{proof}
链同伦映射诱导相同上同调，故函子已在 \(K(\mathcal A)\) 上定义。对标准锥三角，正合列由逐次分裂短正合列 \(0\to B\to C_f\to A[1]\to0\) 得到，连接映射是 \(H^{n+1}(f)\)，见\cref{b4:prop:cone-long-exact}。因此其中的 \(H^n(A)\to H^n(B)\to H^n(C_f)\) 正合。好三角与标准锥三角同构，正合性随同构保留。
\end{proof}

\subsection{Hom 函子的正合性}
\label{b4:subsec:1603:02}

\begin{proposition}\label{b4:prop:triangle-hom-exact}
设 \(\mathcal T\) 为局部小三角范畴，\(W\in\mathcal T\)。函子 \(\operatorname{Hom}_{\mathcal T}(W,-)\) 是取值于交换群的上同调函子；\(\operatorname{Hom}_{\mathcal T}(-,W)\) 具有相应的反变正合性。特别地，好三角 \(X\to Y\to Z\to X[1]\) 给出
\[
 \cdots\to\operatorname{Hom}(W,X[n])\to\operatorname{Hom}(W,Y[n])
 \to\operatorname{Hom}(W,Z[n])\to\operatorname{Hom}(W,X[n+1])\to\cdots,
\]
以及箭头倒向的反变长正合列。
\end{proposition}
\begin{proof}
\cref{b4:lem:pretriangle-hom}已证明三项正合。对三角旋转并平移，可以把长列中的任意一项移到三项列的中间，故在每一项正合。反变情形应用该引理的第二列。
\end{proof}

\begin{proposition}\label{b4:prop:split-triangle-criterion}
设 \(\mathcal T\) 为局部小三角范畴，\(X\xrightarrow{u}Y\xrightarrow{v}Z\xrightarrow{w}X[1]\) 为好三角。则 \(w=0\) 当且仅当此三角同构于由包含、投影和零映射组成的分裂三角 \(X\to X\oplus Z\to Z\to X[1]\)。
\end{proposition}
\begin{proof}
若 \(w=0\)，协变 Hom 长正合列的片段
\[
 \operatorname{Hom}(Z,Y)\xrightarrow{v_*}\operatorname{Hom}(Z,Z)
 \xrightarrow{w_*}\operatorname{Hom}(Z,X[1])
\]
中最后一箭头为零，所以 \(1_Z\) 可提升为 \(s:Z\to Y\)，满足 \(vs=1_Z\)。另一方面，逆向旋转后取反变 Hom，得到
\[
 \operatorname{Hom}(Y,X)\xrightarrow{u^*}\operatorname{Hom}(X,X)
 \xrightarrow{(-w[-1])^*}\operatorname{Hom}(Z[-1],X).
\]
末箭头也为零，所以 \(1_X\) 可提升为 \(r:Y\to X\)，满足 \(ru=1_X\)。连接箭头消失因而同时提供了 \(u\) 的左逆和 \(v\) 的右逆。

要用它们构造 \(Y\cong X\oplus Z\)，还需要使两个分量相容，即 \(rs=0\)。以 \(s-urs\) 代替 \(s\)，由于 \(vu=0\) 且 \(ru=1_X\)，仍有 \(vs=1_Z\)，并且 \(rs\) 变为零。此时 \((r,v)(u,s)=1_{X\oplus Z}\)。为证明另一方向的复合也是恒等，令 \(e=1_Y-ur-sv\)，则 \(ve=0\)、\(re=0\)。对 \(\operatorname{Hom}(Y,-)\) 使用正合性，存在 \(t:Y\to X\) 使 \(e=ut\)；再由 \(t=rut=re=0\)，得 \(e=0\)。所以 \((u,s):X\oplus Z\to Y\) 与 \((r,v)\) 互逆，并保持三角的全部箭头。反向直接从分裂三角的最后箭头为零得到。
\end{proof}

这一判据说明，三角的连接箭头为零时，中间对象连同两端映射一起分裂；这与\cref{b4:thm:ext1-extensions}中短正合扩张的零类对应分裂扩张相呼应。

\ProseParagraphBreak
设 \(\mathcal A\) 为局部小加性范畴，\(C,D\) 为其中的复形。在 \(K(\mathcal A)\) 中还有一个计算形式：
\[
 \operatorname{Hom}_{K(\mathcal A)}(C,D[n])
 \cong H^n\operatorname{Hom}_{\mathcal A}^{\bullet}(C,D).
\]
右端按\cref{b4:def:hom-complex}的公式，把 Hom 换为 \(\mathcal A\) 的态射交换群。次数 \(n\) 的余圈条件为 \(\mathrm{d}_Df=(-1)^nf\mathrm{d}_C\)，恰是 \(C\to D[n]\) 的复形映射条件。还须比较两边所商掉的子群。若 \(g\) 是次数 \(n-1\) 的 Hom 元素，则
\[
 \mathrm{d}_{\operatorname{Hom}}g
 =\mathrm{d}_Dg-(-1)^{n-1}g\mathrm{d}_C
 =\mathrm{d}_Dg+(-1)^ng\mathrm{d}_C.
\]
把 \(h=(-1)^ng\) 看作从 \(C\) 到 \(D[n]\) 的降次态射族，得到
\[
\begin{aligned}
 \mathrm{d}_{D[n]}h+h\mathrm{d}_C
 &=(-1)^n\cdot\mathrm{d}_Dh+h\mathrm{d}_C\\
 &=\mathrm{d}_Dg+(-1)^ng\mathrm{d}_C
 =\mathrm{d}_{\operatorname{Hom}}g.
\end{aligned}
\]
反过来，每个这样的 \(h\) 都唯一写成 \((-1)^ng\)，故 Hom 复形的余边与到 \(D[n]\) 的零伦映射组成同一子群，取商便得到上述同构。此处必须取整个 Hom 复形的上同调，而不能只看某一个次数的普通 Hom 群。

\subsection{三角与短正合列的对应范围}
\label{b4:subsec:1603:03}

\begin{proposition}\label{b4:prop:degreewise-split-triangle}
设 \(\mathcal A\) 为局部小交换范畴，\(0\to A\xrightarrow{j}B\xrightarrow{q}C\to0\) 是逐次分裂的复形短正合列。则存在 \(\delta:C\to A[1]\)，使 \(A\xrightarrow{j}B\xrightarrow{q}C\xrightarrow{\delta}A[1]\) 是 \(K(\mathcal A)\) 中的好三角。
\end{proposition}
\begin{proof}
选逐次分裂，写 \(B^n=A^n\oplus C^n\)，微分为
\(\mathrm{d}_B=\left(\begin{smallmatrix}\mathrm{d}_A&t\\0&\mathrm{d}_C\end{smallmatrix}\right)\)。条件 \(\mathrm{d}_B^2=0\) 给出 \(\mathrm{d}_At+t\mathrm{d}_C=0\)，故 \(-t:C\to A[1]\) 为复形映射。自然映射
\[
 r:\operatorname{Cone}(j)\to C,\quad r((a,c),a')=c,
 \qquad s(c)=((0,c),-tc)
\]
满足 \(rs=1\)。取 \(h((a,c),a')=((0,0),a)\)，便有 \(1-sr=\mathrm{d}h+h\mathrm{d}\)。于是 \(r\) 为同伦等价，其与锥的包含复合是 \(q\)。记锥的最后投影为 \(p_j:\operatorname{Cone}(j)\to A[1]\)，则 \(p_js=-t\)。由于 \(sr=1\) 在同伦范畴中成立，还有
\[
 (-t)r=p_jsr=p_j\qquad\text{于}\;K(\mathcal A).
\]
这同时核对了三角同构的最后一个方块。因此把锥三角经 \(r\) 运送，得到 \(\delta=-t\)。
\end{proof}

为比较这里的 \(\delta=-t\) 与普通短正合列的连接同态，记 \(z_C^n:Z^n(C)\to C^n\) 为余圈嵌入。逐次截面将它送到 \(B^n\) 的第二分量；再作微分得到 \((t^nz_C^n,0)\)。由\cref{b4:prop:connecting-map-construction}，\(t^nz_C^n\) 在上同调上诱导通常的连接同态。在模中，这就是把余圈 \(c\) 提升为 \((0,c)\) 并计算
\[
 \mathrm{d}_B(0,c)=(t(c),0),\qquad
 \delta_{\mathrm{LES}}[c]=[t(c)],\qquad
 H^n(\delta)[c]=-[t(c)].
\]
这里使用 \(H^n(A[1])=H^{n+1}(A)\) 的通常识别。因此三角连接箭头在上同调上的作用为 \(-\delta_{\mathrm{LES}}\)；负号来自本书把锥三角最后箭头固定为正投影的约定。它不影响正合性，却必须在比较连接同态的具体公式时保留。

\ProseParagraphBreak
若短正合列不逐次分裂，则不能直接使用这个结论。例如次数零上的 \(0\to\mathbb Z\xrightarrow{2}\mathbb Z\to\mathbb Z/2\mathbb Z\to0\)，\(\operatorname{Cone}(2)\to\mathbb Z/2\mathbb Z\) 是拟同构，却不是同伦等价：其同伦逆会给出从有限群到自由交换群的非零映射。到导出范畴中，这个拟同构才成为同构，短正合列也才无须逐次分裂即可给出好三角，见\cref{b4:prop:derived-short-exact-triangle}。

% !TeX root = ../../Book4.tex
% 纲要标题：### 17.4 映射锥的函子性与选择
% 纲要标签：b4:sec:1604
% 纲要位置：计划纲要.md，第 2913 行。

\section{映射锥的函子性与选择}
\label{b4:sec:1604}

映射锥在复形及其严格交换方块组成的箭头范畴上构成函子。只保留同伦类及普通三角中的交换方块时，补全锥映射所用的同伦没有被指定，因而这些数据本身不决定典范的锥映射。


% 纲要内容（保留纲要核验项目）：
%
% * 三角范畴中锥的选择
% * 同构不等于典范选择
% * DG 增强的动机
%
% ---
%

\subsection{三角范畴中锥的选择}
\label{b4:subsec:1604:01}

在复形层面，选定 \(f\) 后，\(\operatorname{Cone}(f)=B\oplus A[1]\) 有明确公式。若方块 \(bf=f'a\) 严格交换，则 \((y,x)\mapsto(by,ax)\) 给出两个锥之间的映射，而且保持复合。因此严格交换方块构成的箭头范畴上确实有锥函子。

\ProseParagraphBreak
转到同伦范畴，方块只要求 \(bf-f'a=\mathrm{d}h+h\mathrm{d}\)。构造锥之间的映射便需要选 \(h\)，得到 \((y,x)\mapsto(by+hx,ax)\)。换一个同伦，结果未必同伦。这正是 (TR3) 只要求存在补全的原因：先忘掉同伦本身，再仅凭四条边，通常无法恢复原先的锥映射。

\subsection{同构与典范选择}
\label{b4:subsec:1604:02}

取域 \(k\)，令 \(X=k[-1]\)、\(Y=k[0]\)，考虑零映射 \(X\to Y\)。其锥是次数零上的 \(k\oplus k\)，三角为
\[
 X\xrightarrow{0}Y\longrightarrow Y\oplus X[1]\longrightarrow X[1].
\]
对每个 \(\lambda\in k\)，矩阵
\[
 c_\lambda=\begin{pmatrix}1&\lambda\\0&1\end{pmatrix}
\]
都与包含和投影相容，因此 \((1_X,1_Y,c_\lambda)\) 是三角自同构。不同 \(\lambda\) 在同伦范畴中仍不同，因为次数零复形之间没有能改变映射的非零链同伦。

\ProseParagraphBreak
这个例子中，固定前两个分量仍得到不同的第三分量，因而相容同构没有唯一性。若要让锥映射随方块函子性地变化，还须指定相容的选择，并检验这些选择保持恒等态射与复合。

\subsection{DG 增强的动机}
\label{b4:subsec:1604:03}

要保留上面的 \(h\)，可以不把 Hom 立即压缩成同伦类，而保存整个 Hom 复形。零次余圈是复形映射，次数负一的元素是候选同伦，等式 \(\mathrm{d}h=f-g\) 记录两映射为何同伦。进一步的负次数还记录同伦之间的相容关系。

\ProseParagraphBreak
这给出\term[DG-zengqiang]{DG 增强}[DG enhancement]的动机：一个普通范畴之外，还保留态射复形及与微分相容的复合。取这些复形的零次上同调后，才回到同伦范畴。\chapref{b4:ch:20}将给出正式定义和实例。


\begin{chaptersummary}
拟同构在同伦范畴中满足交换和消去条件，因此导出态射可写成屋顶，屋顶的相等由共同细化判断。这个计算既证明局部化的泛性质，也允许把三角公理从同伦范畴传下来。由此，一般复形短正合列都能给出导出范畴中的好三角。

\ProseParagraphBreak
截断负责把上同调有界的对象换成项有界的代表，投射和内射替换负责把态射计算化为同伦计算。它们解决的问题不同：有界对象可能仍需向一端无限延伸的特殊消解。Ext 因而成为到平移对象的态射群，而这些态射又保留了普通上同调对象之间映射看不到的扩张信息。
\end{chaptersummary}

\BookExercises

\begin{exercise}\label{b4:ex:17-cyclic-ext}
设 \(m,n\ge1\)。计算 \(\operatorname{Hom}_{D(\mathbf{Ab})}(\mathbb Z/m,\mathbb Z/n[r])\) 对所有整数 \(r\) 的值，并写出 \(r=0,1\) 的核与余核表达。
\end{exercise}
\begin{solution}
用 \(\mathbb Z\xrightarrow{m}\mathbb Z\) 替换第一变量，Hom 复形仅在零、一次非零，两项都是 \(\mathbb Z/n\)，微分为负乘 \(m\)。故零次为 \((\mathbb Z/n)[m]\)，一次为 \((\mathbb Z/n)/m(\mathbb Z/n)\)，二者均同构于 \(\mathbb Z/\gcd(m,n)\)，其余为零。零次与一次的具体子群、商群不同，不把这种同构类型误认成同一内部对象。
\end{solution}

\begin{exercise}\label{b4:ex:17-roof-localize-zero}
设 \(f:X\to Y\) 为复形映射。若它经某个零调复形分解，证明 \(Q(f)=0\)。反之，若 \(Q(f)=0\)，证明 \(f\) 在同伦范畴中经某个零调复形分解。
\end{exercise}
\begin{solution}
正向因为中间对象在 \(D\) 中为零。反向由屋顶判据取拟同构 \(s:U\to X\)，使 \(fs=0\)。补成三角 \(U\xrightarrow{s}X\to C_s\to U[1]\)。反变 Hom 正合性使 \(f\) 经 \(C_s\) 分解；拟同构的锥零调，故满足要求。
\end{solution}

\begin{exercise}\label{b4:ex:17-roof-two-denominators}
设两个屋顶 \((s,f)\)、\((s,g)\) 具有同一个左分母。证明它们在导出范畴中相等，当且仅当存在拟同构 \(u\) 使 \(fu=gu\)。
\end{exercise}
\begin{solution}
若有这样的 \(u\)，共同细化立即给出相等。反向，乘上可逆的 \(Q(s)\) 得 \(Q(f)=Q(g)\)，再用普通态射的屋顶相等判据。也可以直接从共同细化 \(sa=sb\)、\(fa=gb\) 出发，先用拟同构 \(s\) 的消去条件进一步使 \(a=b\)。
\end{solution}

\begin{exercise}\label{b4:ex:17-field-two-degrees}
设域 \(k\) 上的复形 \(X\) 仅有 \(H^{-1}(X)=k^2,H^0(X)=k\)，\(Y\) 仅有 \(H^{-1}(Y)=k,H^0(Y)=k^3\)。计算 \(\dim_k\operatorname{Hom}_D(X,Y[r])\)。所得维数是否依赖分裂的选择？
\end{exercise}
\begin{solution}
向量空间复形分解为上同调与可缩部分，故计算可在零微分复形上进行。\(r=0\) 时维数为 \(2\cdot1+1\cdot3=5\)；\(r=1\) 时从 \(X^{-1}=k^2\) 到 \(Y[1]^{-1}=k^3\)，维数为六；\(r=-1\) 时从 \(X^0=k\) 到 \(Y[-1]^0=k\)，维数为一；其余为零。分裂不典范，但 Hom 群本身及其维数由导出对象确定，不依赖选择。
\end{solution}

\begin{exercise}\label{b4:ex:17-injective-cyclic}
设整数 \(m\ge2\)。用内射消解
\[
 0\to\mathbb Z\to\mathbb Q\to\mathbb Q/\mathbb Z\to0
\]
计算 \(\operatorname{Hom}_D(\mathbb Z/m,\mathbb Z[1])\)，并描述一个生成元。
\end{exercise}
\begin{solution}
内射复形取 \(I^0=\mathbb Q,I^1=\mathbb Q/\mathbb Z\)。Hom 的零次项为零，一次项为 \((\mathbb Q/\mathbb Z)[m]\)。因此所求群为 \(\mathbb Z/m\)，生成元由 \(\bar1\mapsto1/m+\mathbb Z\) 给出。与投射屋顶的具体符号比较时须用正文的每度重标识，不能仅比较写下的两个数值。
\end{solution}

\begin{exercise}\label{b4:ex:17-projective-dimension}
设 \(R\) 为环，\(M\) 为左 \(R\)-模，\(d\ge0\)。证明 \(\operatorname{pd}_R M\le d\) 当且仅当对所有左模 \(N\) 都有 \(\operatorname{Hom}_D(M,N[d+1])=0\)。
\end{exercise}
\begin{solution}
由 Ext 识别，条件等价于所有 \(\operatorname{Ext}_R^{d+1}(M,N)=0\)。若有长度 \(d\) 的投射消解，消解计算直接给出消失。反向取任意投射消解，维数平移给出 \(\operatorname{Ext}_R^1(\Omega^dM,N)=0\) 对所有 \(N\) 成立，故 \(\Omega^dM\) 投射，将原消解在这里截短即可。\(d=0\) 时用一次 Ext 的投射判据。
\end{solution}

\begin{exercise}\label{b4:ex:17-cubic-truncated}
令 \(R=k[t]/(t^3)\)、\(k=R/(t)\)。构造 \(k\) 的周期投射消解，并计算它到各次移位 \(k[r]\) 的导出态射群。由此说明 \(k\) 不能由有界投射复形代表。
\end{exercise}
\begin{solution}
消解交替使用乘 \(t\) 与乘 \(t^2\)，靠近增广的箭头为乘 \(t\)。因为 \(\ker t=(t^2)\)、\(\ker t^2=(t)\)，它正合。施加 \(\operatorname{Hom}_R(-,k)\) 后微分均为零，故所求群在 \(r\ge0\) 为 \(k\)，在 \(r<0\) 为零。若存在有界投射代表，则到 \(k[r]\) 的同伦态射在充分大的 \(r\) 必为零，与计算矛盾。
\end{solution}

\begin{exercise}\label{b4:ex:17-cone-localized}
设 \(s:X\to Y\) 为拟同构，\(f:Y\to Z\) 为复形映射。证明 \(\operatorname{Cone}(fs)\) 与 \(\operatorname{Cone}(f)\) 在导出范畴中同构；这种结论在同伦范畴中是否总成立？
\end{exercise}
\begin{solution}
对严格方块 \((s,1_Z)\) 取锥映射。上同调长正合列与五引理使该锥映射拟同构，故在 \(D\) 中同构。若取零调但不可缩的 \(X\)，令 \(s:X\to0\)、\(f:0\to0\)，则两个锥分别为 \(X[1]\) 和零，在 \(K\) 中不同构。因此不能去掉局部化。
\end{solution}

\begin{exercise}\label{b4:ex:17-cohomology-detect}
设 \(R\) 为环，\(X\) 为任意左模复形。证明 \(\operatorname{Hom}_{D(R)}(R[-n],X)\cong H^n(X)\)。由此证明若所有这些群为零，则 \(X=0\)。
\end{exercise}
\begin{solution}
\(R[-n]\) 是只有一项的投射复形，可用同伦类计算导出态射。映射由 \(X^n\) 的一个余圈决定，两个映射同伦恰当其差为余边，故所求群为 \(H^n(X)\)。所有群为零表示 \(X\) 零调，等价于 \(X\) 在导出范畴中为零。
\end{solution}

\begin{exercise}\label{b4:ex:17-split-cohomology-not-split}
设整数 \(m\ge2\)，考虑由非分裂短正合列 \(0\to\mathbb Z\xrightarrow{m}\mathbb Z\to\mathbb Z/m\to0\) 给出的好三角。证明其最后箭头非零，尽管它诱导的每个上同调对象之间的映射都为零。
\end{exercise}
\begin{solution}
若最后箭头为零，由分裂三角判据，\(\mathbb Z\to\mathbb Z/m\) 在导出范畴中有截面。取 \(H^0\) 得交换群满射的截面，与 \(\mathbb Z\) 无有限阶元素矛盾。所以最后箭头非零。它从 \(\mathbb Z/m\) 指向 \(\mathbb Z[1]\)，两端非零上同调的次数不同，因此每个 \(H^n\) 上的映射均零。
\end{solution}

\begin{exercise}\label{b4:ex:17-quasi-between-models}
设 \(P,Q\) 为逐项投射且有上界的复形，\(s:P\to Q\) 为拟同构。证明它已经是同伦等价，并说明不能仅以“拟同构的锥正合”结束证明。
\end{exercise}
\begin{solution}
锥 \(C_s\) 也逐项投射且有上界。把零伦引理应用于其恒等映射 \(C_s\to C_s\)，因目标正合，得到锥可缩。三角 Hom 正合列说明 \(s\) 在 \(K\) 中同构，即同伦等价。正合本身不足以推出可缩，这里必须同时使用锥逐项投射及有上界。
\end{solution}

\begin{exercise}\label{b4:ex:17-truncation-object}
设复形 \(X\) 只有次数 \(n\) 的上同调可能非零。构造 \(D(\mathcal A)\) 中 \(X\cong H^n(X)[-n]\) 的同构，不选各项短正合列的分裂。
\end{exercise}
\begin{solution}
取 \(\tau_{\le n}X\to X\)，它是拟同构。再将 \(\tau_{\le n}X\) 的次数 \(n\) 映到 \(H^n(X)\)，其他次数映到零，得到复形映射 \(\tau_{\le n}X\to H^n(X)[-n]\)。它在唯一可能非零的上同调上为恒等，因而也是拟同构。这两个箭头给出所需同构；全部映射只用核和商，无须分裂。
\end{solution}

\begin{exercise}\label{b4:ex:17-localization-functor}
设 \(F:\mathcal A\to\mathcal B\) 为交换范畴间的正合加性函子。证明逐项 \(F\) 诱导三角函子 \(D(\mathcal A)\to D(\mathcal B)\)，并说明为何仅加性不足。
\end{exercise}
\begin{solution}
正合性使 \(H^n(FC)\cong F(H^nC)\)，故逐项 \(F\) 保持拟同构，由局部化泛性质下降。加性使它保持有限双积和平移的符号，也保持锥矩阵，所以保持好三角。仅加性时拟同构可能不再保持，例如把 \(\mathbb Z\xrightarrow{2}\mathbb Z\to\mathbb Z/2\) 的投射替换与 \(\mathbb Z/2\) 张量会新增负一次上同调，因而不能直接下降。
\end{solution}

\begin{exercise}\label{b4:ex:17-abelian-embedding}
设 \(\mathcal A\) 有足够内射对象。证明将对象放在次数零给出全忠实函子 \(\mathcal A\to D^b(\mathcal A)\)，但这不使 \(D^b(\mathcal A)\) 自动成为交换范畴。
\end{exercise}
\begin{solution}
次数零的 Ext 识别给出 \(\operatorname{Hom}_{D}(M,N)=\operatorname{Hom}_{\mathcal A}(M,N)\)，有界满嵌入使它也等于 \(D^b\) 中的 Hom，因此全忠实。可是三角范畴的所有范畴单态射必须分裂。例如 \(D^b(\mathbf{Ab})\) 中，乘二 \(\mathbb Z\to\mathbb Z\) 不分裂；若该范畴交换，其每个态射都有广义逆，便要求 \(4t=2\) 对整数 \(t\) 成立，矛盾。原交换范畴只是一个满子范畴，核余核不能无条件在较大范畴中照搬。
\end{solution}

% !TeX root = ../../Book4.tex
% 纲要标题：## 第18章　导出范畴的构造
% 纲要标签：b4:ch:17
% 纲要位置：计划纲要.md，第 3682 行。

\chapter{导出范畴的构造}
\label{b4:ch:17}
\BookChapterNumber{b4:ch:17}

\begin{chapterabstract}
本章把拟同构在同伦范畴中形式地变为可逆，建立屋顶演算与导出范畴的三角结构，再用有界投射、内射模型计算导出态射与 Ext。
\end{chapterabstract}

\begin{learninggoals}
\begin{itemize}
\item 用泛性质区别导出范畴与仅列出上同调对象的做法。
\item 通分、合成和比较屋顶，说明拟同构的消去条件。
\item 证明导出范畴的三角结构及短正合列所产生的好三角。
\item 区分项有界、上同调有界和特殊消解的有界方向。
\item 用投射或内射模型计算导出态射，并核对 Ext 的次数与符号。
\end{itemize}
\end{learninggoals}

自由模的两项复形 \(P=(\mathbb Z\xrightarrow{2}\mathbb Z)\) 向只在零度放置 \(\mathbb Z/2\mathbb Z\) 的复形有一个自然的链映射：末项取商，另一项送到零。它在同调上是同构，却没有逐项的逆映射。若计算结果本来就不应取决于选用 \(P\) 还是直接使用商模，便要允许这样的链映射具有逆，而不能要求逆仍由逐项映射给出。导出范畴把这类拟同构正式变为可逆；「屋顶」表示记录了先更换模型、再施行态射的过程。接下来的任务是证明这种新态射仍与锥和三角结构相容。

\ProseParagraphBreak
本章使用\chapref{b4:ch:16}的三角公理和 Hom 正合性。\secref{b4:sec:1704}的模型另外调用\chapref{b4:ch:07}的投射、内射消解，\chapref{b4:ch:08}的 Ext，\cref{b4:thm:ce-resolution}中有下界复形的 Cartan--Eilenberg 消解，以及\cref{b4:cor:double-complex-comparison}的双复形比较。前三节建立导出范畴及其三角结构；\secref{b4:sec:1704}进一步给出计算态射所需的有界模型。

\ProseParagraphBreak
局部化与屋顶可对照 Weibel\cite[Section 10.3; Proposition 10.4.1]{Weibel1994HomologicalAlgebra}；有界消解模型见该书\cite[Corollary 10.4.7; Theorem 10.4.8]{Weibel1994HomologicalAlgebra}。本章同时说明一般交换范畴的集合大小约定，并在模范畴中给出局部小性的具体证明。

% !TeX root = ../../Book4.tex
% 纲要标题：### 18.1 局部化
% 纲要标签：b4:sec:1701
% 纲要位置：计划纲要.md，第 3684 行。

\section{局部化}
\label{b4:sec:1701}

拟同构保存同调，却未必已有同伦逆。将全部拟同构正式变为可逆态射，能允许更换消解，同时仍保留态射的复合及其泛性质。


% 纲要内容（保留纲要核验项目）：
%
% * 把拟同构变成可逆态射
% * 局部化的泛性质
% * 集合大小与构造假设
%

\subsection{拟同构的可逆化}
\label{b4:subsec:1701:01}

链同伦等价必为拟同构，但反向不成立。消解 \(P\to M\) 的作用恰在于：虽然它可能不是同伦等价，上同调信息却与 \(M\) 完全相同。若希望既允许换消解又保持态射的复合，就应把所有拟同构都当作可逆态射。

\begin{definition}\label{b4:def:derived-category}
设 \(\mathcal A\) 为局部小交换范畴，\(S\) 为 \(K(\mathcal A)\) 中的拟同构，即在各次上同调上诱导同构的态射类。在适当的集合大小约定下，取局部化
\[
 D(\mathcal A)=K(\mathcal A)[S^{-1}].
\]
所得范畴称为 \(\mathcal A\) 的\term[daochu-fanchou]{导出范畴}[derived category]，其局部化函子记为 \(Q:K(\mathcal A)\to D(\mathcal A)\)。这里的逆是新范畴中的逆，不要求拟同构原本具有复形映射意义的逆。
\end{definition}
\symidx{\(D(\mathcal A)\)：复形的导出范畴}

一个零调复形在 \(D(\mathcal A)\) 中成为零，因为它到零复形的映射是拟同构。反过来，各次上同调函子能经 \(Q\) 分解，所以在 \(D(\mathcal A)\) 中为零的对象必零调。因此，复形 \(C\) 在 \(K(\mathcal A)\) 中为零对象当且仅当它可缩，在 \(D(\mathcal A)\) 中为零对象当且仅当它零调；零调但不可缩的复形也在局部化后与零复形同构。

\ProseParagraphBreak
例如，将正合列
\[
 0\longrightarrow\mathbb Z\xrightarrow{2}\mathbb Z
 \xrightarrow{q}\mathbb Z/2\mathbb Z\longrightarrow0
\]
的三个非零项依次放在次数零、一、二，得到零调复形 \(C\)。若它可缩，同伦等式在次数二要求 \(q h^2=1_{\mathbb Z/2\mathbb Z}\)，其中 \(h^2:\mathbb Z/2\mathbb Z\to\mathbb Z\)。但这样的同态只能为零，不可能给出 \(q\) 的右逆。因此 \(C\) 在同伦范畴中非零，在导出范畴中却为零。

\subsection{局部化的泛性质}
\label{b4:subsec:1701:02}

\begin{definition}\label{b4:def:category-localization}
设 \(\mathcal C\) 为范畴，\(S\) 为一类态射，给定函子 \(Q:\mathcal C\to\mathcal C[S^{-1}]\)。若 \(Q(s)\) 对每个 \(s\in S\) 都可逆，并且对每个范畴 \(\mathcal E\)，任何将 \(S\) 送为同构的函子 \(F:\mathcal C\to\mathcal E\) 都唯一写成 \(F=\overline FQ\)，自然变换也唯一经此分解，则称 \(Q\) 为 \(\mathcal C\) 在 \(S\) 处的\term[fanchou-jubuhua]{范畴局部化}[localization of a category]。采用对象不变的标准构造时，这里的唯一性是严格的；更换等价模型后则为唯一到自然同构。
\end{definition}
\symidx{\(\mathcal C[S^{-1}]\)：在指定态射类处的范畴局部化}

小范畴的一个直接构造，是在原有箭头之外为每个 \(s\in S\) 添一个反向符号 \(s^{-1}\)，取有限可复合箭头串，再按原范畴的复合关系以及 \(ss^{-1}=1,s^{-1}s=1\) 取商。复合由连接箭头串给出。函子 \(F\) 在新符号上只能取 \(F(s)^{-1}\)，因而唯一延拓；自然变换在新符号上的自然性由在 \(s\) 上的自然性乘逆得到。这也证明了自然变换的分解性质。

\ProseParagraphBreak
形式逆不必来自原范畴中的映射。取位于次数负一、零的 \(P=(\mathbb Z\xrightarrow{2}\mathbb Z)\)，以及拟同构 \(s:P\to(\mathbb Z/2\mathbb Z)[0]\)。任何反向复形映射 \(a:(\mathbb Z/2\mathbb Z)[0]\to P\) 在次数零都是 \(\mathbb Z/2\mathbb Z\to\mathbb Z\) 的零同态，故不能满足 \(sa=1\)，也不能成为 \(s\) 的同伦逆。局部化仍使 \(Q(s)^{-1}\) 存在；下一节的屋顶将通过把源对象换成 \(P\)，而非寻找这样的 \(a\)，来表示这个逆。

\ProseParagraphBreak
这种构造证明存在，却不便计算。下一节证明：同伦范畴中的拟同构允许把任意箭头串缩成一个分母在左的屋顶。这里不能先假定任意范畴都有这样的分式形式；它需要专门的交换与消去条件。

\ProseParagraphBreak
还可直接从复形范畴把拟同构倒过来。两种做法等价：拟同构包含链同伦等价，故\cref{b4:prop:homotopy-universal}迫使任何这样的函子先经 \(K(\mathcal A)\) 分解；之后再用局部化的泛性质即可。

\subsection{集合大小与构造假设}
\label{b4:subsec:1701:03}

若 \(\mathcal A\) 的对象不是一个集合，前述箭头串也可能形成真类，不能仅凭原范畴局部小就断言局部化仍局部小。本书对一般交换范畴使用固定的较大集合范围来作局部化；讨论 Hom 群时，要么假定局部化在所用范围内局部小，要么使用下面模范畴的大小控制，或\secref{b4:sec:1704}的有界消解模型。任意交换范畴的无界情形还需要额外的局部小性证明。

\begin{proposition}\label{b4:prop:module-derived-size}
设 \(R\) 为含幺环，\(X\) 为左 \(R\)-模复形。存在一集合的拟同构 \(X_i\to X\)，使每个拟同构 \(s:Y\to X\) 都能由某个 \(X_i\to X\) 经拟同构 \(X_i\to Y\) 细化。因此在允许屋顶演算后，\(D(R\text{-}\mathbf{Mod})\) 的每个态射类确实组成集合。
\end{proposition}
\begin{proof}
取无限基数 \(\kappa\) 至少为 \(|R|\)、\(\aleph_0\) 及所有 \(|X^n|\) 的上界。给定 \(s\)，为 \(X\) 每次上同调的每个类选一个 \(Y\) 中余圈，其像代表该类。所选元素生成一个各次大小至多 \(\kappa\) 的子复形 \(Y_0\subseteq Y\)。这一步保证 \(H(Y_0)\to H(X)\) 满，但缩小到子复形后，一个在 \(Y\) 中为余边的余圈未必仍在 \(Y_0\) 中为余边，所以还需消去这个映射的核。

归纳构造 \(Y_j\)：对 \(Y_j\) 中每个余圈 \(z\)，若 \(s(z)\) 在 \(X\) 中为余边，则 \(s\) 的上同调单射性保证 \(z=\mathrm{d}y\) 对某个 \(y\in Y\) 成立。把这些 \(y\) 加入并生成子复形，得 \(Y_{j+1}\)。新加入的元素还可能产生新的余圈，因此要重复这个过程。每一步至多加入 \(\kappa\) 个元素；取可数并 \(Y'=\bigcup_jY_j\)，其大小仍至多 \(\kappa\)。起初的选择保证 \(H(Y')\to H(X)\) 满。任一落入核的余圈 \(z\in Y'\) 已属于某个 \(Y_j\)，它在 \(Y_{j+1}\) 中成为余边，所以该映射也单。于是 \(Y'\to X\) 是拟同构；它是 \(Y'\to Y\to X\) 的复合，后二者中的 \(Y\to X\) 也为拟同构，故 \(Y'\to Y\) 同样为拟同构。

所有大小至多 \(\kappa\) 的模复形可以在固定集合上编码：每次的底集合取为一个固定的 \(\kappa\) 元集合的子集，模运算和微分都是这些集合上的函数，可能的选择组成集合。它们到 \(X\) 的复形映射也组成集合。取其代表即得所需的一集合分母。每个屋顶 \(X\leftarrow Y\to Z\) 都可换成以某个 \(X_i\) 为中间对象的屋顶，而 \(\bigcup_i\operatorname{Hom}_{K}(X_i,Z)\) 为集合，其商仍为集合。
\end{proof}

这里用 \(\le\kappa\) 控制大小，避免了“可数次合并以后仍严格小于某个基数”的不当推断。相关大小问题可对照 Weibel\cite[Set-Theoretic Considerations 10.3.6; Proposition 10.4.4]{Weibel1994HomologicalAlgebra}。

% !TeX root = ../../Book4.tex
% 纲要标题：### 18.2 屋顶与分式
% 纲要标签：b4:sec:1702
% 纲要位置：计划纲要.md，第 3690 行。

\section{屋顶与分式}
\label{b4:sec:1702}

导出态射可以先沿拟同构更换源对象，再作复形映射。屋顶的共同细化既给出相等判据，也使复合、加法与三角结构可以严格构造。


% 纲要内容（保留纲要核验项目）：
%
% * 屋顶图表示态射
% * 适当的分式演算
% * 同伦范畴到导出范畴
%

\subsection{态射的屋顶图表示}
\label{b4:subsec:1702:01}

\begin{definition}\label{b4:def:derived-roof}
设 \(\mathcal A\) 为局部小交换范畴，在固定的集合大小约定下记局部化函子为 \(Q:K(\mathcal A)\to D(\mathcal A)\)。\(K(\mathcal A)\) 中的图式
\[
 X\xleftarrow{s}U\xrightarrow{f}Y,\qquad s\;\text{为拟同构},
\]
称为一个\term[wuding]{屋顶}[roof]，表示 \(Q(f)Q(s)^{-1}:X\to Y\)。给定另一个屋顶 \(X\xleftarrow{s'}U'\xrightarrow{f'}Y\)，若存在对象 \(V\) 及拟同构 \(a:V\to U\)、\(b:V\to U'\)，使 \(sa=s'b\)、\(fa=f'b\)，则称两个屋顶有共同细化；所有等式都在同伦范畴中理解。
\end{definition}

共同细化把两个屋顶放在同一张图中：
\[
\begin{tikzcd}[column sep=3.2em,row sep=2.2em]
&U\arrow[dl,"s\;(\sim)"']\arrow[dr,"f"]&\\
X&V\arrow[u,"a\;(\sim)"]\arrow[d,"b\;(\sim)"']&Y\\
&U'\arrow[ul,"s'\;(\sim)"]\arrow[ur,"f'"']&
\end{tikzcd}
\]
左侧要求 \(sa=s'b\)，右侧要求 \(fa=f'b\)。这里 \(a,b\) 都是拟同构；也可以先要求共同左分母 \(sa=s'b\) 为拟同构，再由二取三性质推出这一点。符号 \(Q(s)^{-1}\) 是局部化以后得到的逆，图式本身只使用同伦范畴中实际存在的箭头，没有要求在那里选择 \(s\) 的逆映射。

\ProseParagraphBreak
在模消解的典型例子中，\(P\to M\) 是左分母，而 \(P\to N[n]\) 记录一个 Ext 余圈。屋顶因此不是额外的画图规则，而是“先换成可计算的模型，再作映射”的正式表达。接下来证明共同细化足以判别相等。

\subsection{分式演算}
\label{b4:subsec:1702:02}

屋顶的复合需要把指向同一对象的两张箭头通分；比较不同的通分选择，则需要消去两张箭头的差。下面第一项构造等式 \(ft=sg\)，使复合所需的方块交换；第二项从 \(sf=sg\) 出发，允许进一步前复合一个拟同构后得到 \(ft=gt\)，用来保证计算不依赖选择。第三、第四项分别是后复合方向的通分和消去，它们给出分母在右的分式演算。

\begin{lemma}\label{b4:lem:quasi-ore}
设 \(\mathcal A\) 为局部小交换范畴。\(K(\mathcal A)\) 中的拟同构包含恒等态射，对复合和平移封闭，并满足二取三性质。它们还满足以下四个条件。
\begin{enumerate}
\item 给定 \(f:X\to Y\) 和拟同构 \(s:Z\to Y\)，存在拟同构 \(t:W\to X\) 和 \(g:W\to Z\)，使 \(ft=sg\)。
\item 若拟同构 \(s:Y\to Y'\) 满足 \(sf=sg\)，其中 \(f,g:X\to Y\)，则存在拟同构 \(t:X'\to X\)，使 \(ft=gt\)。
\item 给定 \(f:X\to Y\) 和拟同构 \(s:X\to Z\)，存在拟同构 \(t:Y\to W\) 和 \(g:Z\to W\)，使 \(tf=gs\)。
\item 若拟同构 \(s:X'\to X\) 满足 \(fs=gs\)，其中 \(f,g:X\to Y\)，则存在拟同构 \(t:Y\to Y'\)，使 \(tf=tg\)。
\end{enumerate}
\end{lemma}
\begin{proof}
恒等态射、复合、平移及二取三性质的结论直接对各次上同调核对。对第一项，补全 \(Z\xrightarrow{s}Y\to C\to Z[1]\)。因 \(s\) 为拟同构，\(C\) 零调。再把复合 \(X\xrightarrow{f}Y\to C\) 补成旋转后的好三角 \(W\xrightarrow{t}X\to C\to W[1]\)。长正合列说明 \(t\) 是拟同构；对这两个三角使用 (TR3)，即可得到 \(g\) 及 \(ft=sg\)。

对第二项令 \(h=f-g\)。补全 \(s\) 的锥三角，记 \(C=\operatorname{Cone}(s)\)。由 \(sh=0\) 和\cref{b4:lem:pretriangle-hom}的协变 Hom 正合性，\(h\) 经 \(C[-1]\) 分解为 \(X\xrightarrow{a}C[-1]\to Y\)。这里 \(C[-1]\) 虽然零调，却未必是 \(K(\mathcal A)\) 中的零对象，不能仅由这一分解断言 \(h=0\)。把 \(a\) 补成好三角
\[
 X'\xrightarrow{t}X\xrightarrow{a}C[-1]\longrightarrow X'[1].
\]
由\cref{b4:prop:cohomology-on-homotopy}产生的上同调长正合列及 \(C[-1]\) 零调，\(t\) 是拟同构。相邻复合为零给出 \(at=0\)，所以 \(ht=0\)，即 \(ft=gt\)。零调锥在这里提供的是进一步拟同构后的消去，而非原同伦范畴中的直接消去。

对第三项，补全 \(X\xrightarrow{s}Z\to C\to X[1]\)，其中 \(C\) 零调。逆向旋转后记首箭头为 \(u:C[-1]\to X\)，再把 \(fu:C[-1]\to Y\) 补成好三角
\[
 C[-1]\xrightarrow{fu}Y\xrightarrow{t}W\to C.
\]
对逆向旋转的三角和这个三角应用 (TR3)，前两个分量取 \(1_{C[-1]},f\)，得到 \(g:Z\to W\)，并有 \(tf=gs\)。长正合列及 \(C\) 零调说明 \(t\) 是拟同构。

对第四项令 \(h=f-g\)，补全 \(X'\xrightarrow{s}X\xrightarrow{q}C\to X'[1]\)。由 \(hs=0\) 和反变 Hom 正合性，存在 \(b:C\to Y\) 使 \(h=bq\)。把 \(b\) 补成好三角 \(C\xrightarrow{b}Y\xrightarrow{t}Y'\to C[1]\)。\(C\) 零调，故 \(t\) 是拟同构，而 \(th=tbq=0\)。
\end{proof}

\begin{theorem}\label{b4:thm:roof-calculus}
设 \(\mathcal A\) 为局部小交换范畴，并采用\subsecref{b4:subsec:1701:03}的集合大小约定，记局部化函子为 \(Q:K(\mathcal A)\to D(\mathcal A)\)。\(D(\mathcal A)\) 中每个态射可由屋顶表示，两个屋顶表示同一态射当且仅当它们有共同细化。特别地，对 \(K(\mathcal A)\) 中的 \(f,g:X\to Y\)，\(Q(f)=Q(g)\) 当且仅当存在拟同构 \(t:X'\to X\) 使 \(ft=gt\)；也可用后复合的拟同构 \(Y\to Y'\) 检验。
\end{theorem}
\begin{proof}
先用共同细化定义屋顶的等价。自反、对称显然。对传递性，把两次细化到共同中间屋顶的分母用\cref{b4:lem:quasi-ore}第一项再作共同细化，二取三保证新添的箭头仍为拟同构；复合之后，左右两条边同时相等。这证明它确为等价关系。

复合 \((s,f):X\to Y\) 与 \((t,g):Y\to Z\) 时，对 \(f:U\to Y\) 和 \(t:V\to Y\) 取 \(a:W\to U\)、\(b:W\to V\)，其中 \(a\) 拟同构且 \(fa=tb\)。拼接图为
\[
\begin{tikzcd}[column sep=1.8em,row sep=2.1em]
&&W\arrow[dl,"a\;(\sim)"']\arrow[dr,"b"]&&\\
&U\arrow[dl,"s\;(\sim)"']\arrow[dr,"f"]&&V\arrow[dl,"t\;(\sim)"']\arrow[dr,"g"]&\\
X&&Y&&Z.
\end{tikzcd}
\]
由 \(fa=tb\)，在局部化中有
\[
 Q(t)^{-1}Q(f)=Q(b)Q(a)^{-1}.
\]
因此两个屋顶的复合为
\[
 Q(g)Q(t)^{-1}Q(f)Q(s)^{-1}
 =Q(gb)Q(sa)^{-1},
\]
正是沿图的外侧两条路径得到的屋顶 \((sa,gb)\)。图中标了 \(\sim\) 的 \(s,t,a\) 必须为拟同构，才能使用上述逆；辅助箭头 \(b\) 无须为拟同构。若另一选择为 \((a',b')\)，先共同细化拟同构 \(a,a'\)，得到 \(ac=a'c'\)。于是 \(tbc=tb'c'\)，用引理第二项再前复合一个拟同构，便使 \(bc=b'c'\) 也成立。因此两复合屋顶等价。更换原屋顶的代表时，把其共同细化与这一计算合并，同样得到等价的复合。

三重屋顶的两种括号分别使用有限个这样的交换方块。先将它们涉及的分母逐一共同细化，再用消去条件使所有指向同一对象的两条有限路径相等，就得到一个同时计算两种括号的屋顶；其左边是各分母的复合，右边是各分子的复合。因此结合律归结为 \(K(\mathcal A)\) 的结合律。这里每次只处理三重复合中的有限个方块，不涉及无限图式的同时提升。恒等屋顶为 \((1,1)\)，单位律由取恒等交换方块得到。

函子 \(Q\) 把 \(f\) 送到 \((1,f)\)。若 \(s\) 拟同构，\((s,1)\) 是 \(Q(s)\) 的逆。任意倒转拟同构的函子 \(F\) 必须把屋顶送到 \(F(f)F(s)^{-1}\)；共同细化保证此式与代表无关，复合公式保证它保持复合。这就给出唯一的分解函子 \(\overline F\)。

还需核对自然变换的分解。设 \(F,G:K(\mathcal A)\to\mathcal E\) 都把拟同构送为同构，\(\eta:F\Rightarrow G\) 为自然变换。屋顶范畴的对象不变，所以只能取 \(\overline\eta_X=\eta_X\)。对屋顶 \(X\xleftarrow{s}U\xrightarrow{f}Y\)，由 \(\eta\) 的自然性有
\[
 \eta_YF(f)=G(f)\eta_U,\qquad
 \eta_XF(s)=G(s)\eta_U.
\]
因 \(F(s),G(s)\) 可逆，得到
\[
 \eta_Y\overline F(s,f)
 =G(f)\eta_UF(s)^{-1}
 =G(f)G(s)^{-1}\eta_X
 =\overline G(s,f)\eta_X.
\]
因此这组分量确实给出唯一的自然变换 \(\overline\eta:\overline F\Rightarrow\overline G\)。屋顶范畴满足局部化的全部泛性质，正是 \(D(\mathcal A)\)。最后把两个普通态射看作分母恒等的屋顶，共同细化条件便是所写的相等判据；后复合的判据由引理第三、第四项得到。
\end{proof}

对两个屋顶 \(X\xleftarrow{s}U\xrightarrow{f}Y\)、\(X\xleftarrow{t}V\xrightarrow{g}Y\)，用通分条件取 \(a:W\to U\)、\(b:W\to V\)，使 \(sa=tb\) 为拟同构。二取三性质保证 \(a,b\) 也都是拟同构。两张右箭头此时有同一源对象，才可相加，所得屋顶为
\[
 X\xleftarrow{sa}W\xrightarrow{fa+gb}Y.
\]
同分母时，这就是 \((s,f)+(s,g)=(s,f+g)\)。不同的通分选择可进一步共同细化，并用消去条件使相应右箭头相等；预复合又保持加法，故结果与选择无关。复合公式及原来的双线性给出新复合的双线性。零对象和有限双积由 \(K(\mathcal A)\) 保留下来。因此 \(D(\mathcal A)\) 是加性范畴，\(Q\) 为加性函子。屋顶演算的一般形式可参阅 Weibel\cite[Theorem 10.3.7; Corollaries 10.3.8--10.3.11]{Weibel1994HomologicalAlgebra}。

\subsection{从同伦范畴到导出范畴}
\label{b4:subsec:1702:03}

\begin{theorem}\label{b4:thm:derived-triangulated}
设 \(\mathcal A\) 为局部小交换范畴，并采用\subsecref{b4:subsec:1701:03}的集合大小约定，记局部化函子为 \(Q:K(\mathcal A)\to D(\mathcal A)\)。在 \(D(\mathcal A)\) 中，以平移诱导的自等价为平移，以同构于 \(K(\mathcal A)\) 中好三角之像的三角为好三角，则 \(D(\mathcal A)\) 是三角范畴，\(Q\) 保持平移和好三角。
\end{theorem}
\begin{proof}
平移保持拟同构，所以经局部化诱导自等价。(TR1) 中的恒等三角和对同构封闭性已成立。任意屋顶 \(X\xleftarrow{s}U\xrightarrow{f}Y\) 的左箭头在 \(D\) 中可逆，因此 \(f\) 的锥三角运送后给出该屋顶的好三角。(TR2) 由 \(K\) 中的旋转直接得到。

核对 (TR3) 时，先把两个三角换成 \(K\) 中标准三角的像，首箭头记为 \(f:X\to Y\)、\(f':X'\to Y'\)。给定导出范畴中的交换方块 \((\alpha,\beta)\)，分别取屋顶
\(X\xleftarrow{s}S\xrightarrow{a}X'\)、\(Y\xleftarrow{t}T\xrightarrow{b}Y'\)。沿拟同构 \(t\) 通分 \(fs\)，得到拟同构 \(u:U\to S\) 和 \(k:U\to T\)，满足 \(t k=fsu\)。方块在 \(D\) 中交换，故 \(Q(f'au)=Q(bk)\)。由屋顶相等判据，再取拟同构 \(v:V\to U\)，使 \(f'auv=bkv\) 在 \(K\) 中成立。

上述构造给出 \(K\) 中的两个交换方块，它们共享上方箭头 \(kv:V\to T\)，其交换等式为
\[
 t(kv)=f(suv),\qquad b(kv)=f'(auv).
\]
把 \(kv\) 补成一条标准锥三角，再分别用 \(K\) 中的 (TR3) 映向 \(f\)、\(f'\) 的三角：前一对分量为 \((suv,t)\)，后一对为 \((auv,b)\)。这些方块的等式在 \(K\) 中成立，所选的复形映射代表可以只同伦交换；前章的锥映射公式正是用同伦修正它们。前一补全的第三个分量是拟同构，因为前两个分量是拟同构，对各次上同调长正合列使用\cref{b4:thm:five-lemma}即可。因此这一三角态射在 \(D\) 中成为同构；先取其逆，再复合后一三角态射，就得到原方块的补全。这证明 (TR3)，并使\cref{b4:lem:pretriangle-hom}现在也可在 \(D\) 中使用。

对 (TR4)，设 \(\alpha:X\to Y\)、\(\beta:Y\to Z\) 为两个导出态射。要对它们实际构造锥，需先用屋顶通分，将这一对导出态射换成 \(K\) 中可复合的复形映射。先取 \(\beta\) 的屋顶 \(Y\xleftarrow{t}T\to Z\)，再取 \(\alpha\) 的屋顶 \(X\xleftarrow{s}S\to Y\)，并沿 \(t\) 通分右边。得到 \(K\) 中一串 \(U\to T\to Z\)，以及拟同构 \(U\to X\)、\(T\to Y\)，其局部化像正好代表原来的可复合对。在 \(K\) 中对此串应用\cref{b4:thm:homotopy-triangulated}证明中的显式八面体，再施加 \(Q\) 并沿上述两个同构运送，便得到补全 \(\alpha,\beta,\beta\alpha\) 的三个好三角及它们的八面体。

这样得到的是由所选复形映射构造的三角，而 (TR4) 还要求八面体适配任意事先指定的三个好三角。对每个首箭头 \(f\in\{\alpha,\beta,\beta\alpha\}\)，在所得三角与指定三角的前两项取恒等态射，由已证的 (TR3) 补出第三分量 \(c_f:C_f\to C_f'\)。\cref{b4:lem:pretriangle-hom}只需前三条公理，故可在这里应用：前两个分量为同构，所以 \(c_f\) 也可逆。把所得八面体在第三对象之间的箭头记为 \(a:C_\alpha\to C_{\beta\alpha}\)、\(b:C_{\beta\alpha}\to C_\beta\)、\(\gamma:C_\beta\to C_\alpha[1]\)，定义运送后的箭头为
\[
 a'=c_{\beta\alpha}a c_\alpha^{-1},\qquad
 b'=c_\beta b c_{\beta\alpha}^{-1},\qquad
 \gamma'=c_\alpha[1]\gamma c_\beta^{-1}.
\]
各个 \(c_f\) 都与对应三角的后两条箭头相容，故八面体的四个相容等式保持成立；最后一个面经 \((c_\alpha,c_{\beta\alpha},c_\beta)\) 成为同构的好三角。这便满足针对任意指定三角的 (TR4)。
\end{proof}

\begin{proposition}\label{b4:prop:derived-short-exact-triangle}
设 \(\mathcal A\) 为局部小交换范畴，\(0\to A\xrightarrow{i}B\xrightarrow{q}C\to0\) 是 \(\mathcal A\) 中上链复形的短正合列，不要求逐次分裂。记 \(Q:K(\mathcal A)\to D(\mathcal A)\) 为局部化函子，\(p_i:\operatorname{Cone}(i)\to A[1]\) 为锥的投影，令 \(r:\operatorname{Cone}(i)\to C\) 为 \(r(b,a)=q(b)\)。则 \(r\) 是拟同构，且
\[
 A\xrightarrow{i}B\xrightarrow{q}C\xrightarrow{\delta}A[1],
 \qquad\delta=Q(p_i)Q(r)^{-1},
\]
是 \(D(\mathcal A)\) 中的好三角。
\end{proposition}
\begin{proof}
\(r\) 是逐次满的复形映射，因为 \(q\) 逐次满且 \(qi=0\)。利用 \(i^n\) 识别 \(\ker q^n\) 与 \(A^n\)，有
\[
 (\ker r)^n=\ker q^n\oplus A^{n+1}
 \cong A^n\oplus A^{n+1}.
\]
在此识别下，锥的微分为
\[
 \mathrm{d}_{\ker r}^n(a,a')
 =(\mathrm{d}_A^na+a',-\mathrm{d}_A^{n+1}a'),
\]
恰好是 \(\operatorname{Cone}(1_A)\) 的微分。该锥由\cref{b4:prop:cone-sign-compatibility}可缩，因而零调。对
\(0\to\ker r\to\operatorname{Cone}(i)\xrightarrow{r}C\to0\)
应用\cref{b4:thm:cohomology-long-exact}，得到每个 \(H^n(r)\) 为同构，故 \(r\) 拟同构。又因 \(r\) 与锥的包含复合等于 \(q\)，把 \(i\) 的标准锥三角经 \(Q(r)\) 运送，即得到所述好三角及最后箭头 \(Q(p_i)Q(r)^{-1}\)。这只使用 \(r\) 的拟同构性，无须原短正合列具有复形截面。
\end{proof}

若原短正合列逐次分裂，选 \(B^n=A^n\oplus C^n\)，并写
\[
 \mathrm{d}_B=\begin{pmatrix}\mathrm{d}_A&t\\0&\mathrm{d}_C\end{pmatrix}.
\]
由 \(\mathrm{d}_B^2=0\) 得 \(\mathrm{d}_At+t\mathrm{d}_C=0\)。前章\cref{b4:prop:degreewise-split-triangle}中的复形映射截面，在这里记为
\[
 \sigma^n:C^n\longrightarrow\operatorname{Cone}(i)^n,
 \qquad \sigma^n(c)=((0,c),-t^n(c)).
\]
它的微分满足
\[
 \mathrm{d}_{\operatorname{Cone}(i)}\sigma(c)
 =((0,\mathrm{d}_Cc),\mathrm{d}_At(c))
 =\sigma(\mathrm{d}_Cc),
\]
且 \(r\sigma=1_C\)。局部化后 \(Q(r)\) 可逆，因而 \(Q(\sigma)=Q(r)^{-1}\)，于是
\[
 \delta=Q(p_i)Q(r)^{-1}=Q(p_i\sigma)=Q(-t).
\]
负号来自 \(\sigma\) 的第二分量，与前章的正投影约定一致。在模的情形，普通上同调连接同态由余圈 \(c\) 的提升 \((0,c)\) 及其微分 \((t(c),0)\) 计算，给出 \([c]\mapsto[t(c)]\)；这里三角连接箭头在上同调上给出 \([c]\mapsto-[t(c)]\)。一般交换范畴中的比较已由\cref{b4:prop:degreewise-split-triangle}之后的余圈子对象计算给出，同样相差所选锥约定的负号。

\ProseParagraphBreak
此外，\(D\) 中态射为同构，当且仅当其锥为零：用三角的 Hom 正合列，锥为零使该态射对每个 Hom 群诱导同构，从而可逆；反向由三角上同调或同样的 Hom 列得到。对原来的复形映射，这恰好恢复拟同构判据。

% !TeX root = ../../Book4.tex
% 纲要标题：### 18.3 导出范畴中的有界性
% 纲要标签：b4:sec:1703
% 纲要位置：计划纲要.md，第 3696 行。

\section{导出范畴中的有界性}
\label{b4:sec:1703}

复形的项有界与其同调有界并非同一条件。典范截断允许删去正合尾部，却可能失去投射或内射性，因此不能把有界对象与有界特殊消解混同。


% 纲要内容（保留纲要核验项目）：
%
% * D^b、D^+ 与 D^-
% * 项有界和同调有界的区别
% * 模范畴中的基本实例
%

\subsection{\texorpdfstring{\(D^b\)、\(D^+\) 与 \(D^-\)}{Db、D+ 与 D−}}
\label{b4:subsec:1703:01}

\begin{definition}\label{b4:def:bounded-derived}
设 \(\mathcal A\) 为局部小交换范畴，并固定局部化的集合大小约定。分别将 \(K^b(\mathcal A),K^+(\mathcal A),K^-(\mathcal A)\) 中的拟同构倒转，得到\term[youjie-daochu-fanchou]{有界导出范畴}[bounded derived category] \(D^b(\mathcal A)\)、有下界导出范畴 \(D^+(\mathcal A)\) 和有上界导出范畴 \(D^-(\mathcal A)\)。有界在此先指复形的项。
\end{definition}
\symidx{\(D^b(\mathcal A),D^+(\mathcal A),D^-(\mathcal A)\)：有界、有下界与有上界导出范畴}

\begin{proposition}\label{b4:prop:bounded-derived-full}
设 \(\mathcal A\) 为局部小交换范畴，且相关局部化均采用同一集合大小约定。自然函子 \(D^b(\mathcal A),D^+(\mathcal A),D^-(\mathcal A)\to D(\mathcal A)\) 全忠实；它们的本质像分别由上同调两端有界、有下界、有上界的对象组成。三者各为三角子范畴。
\end{proposition}
\begin{proof}
截去尾部时必须保留截口上的余圈与余边条件。例如，\(\mathbb Z\xrightarrow{2}\mathbb Z\xrightarrow{q}\mathbb Z/2\mathbb Z\) 位于次数零、一、二，原复形零调；若直接删去次数大于一的项，余下的 \(\mathbb Z\xrightarrow{2}\mathbb Z\) 却有 \(H^1=\mathbb Z/2\mathbb Z\)。这是因为末项的全部元素都变成余圈。把次数一的项改为 \(\ker q=2\mathbb Z\)，所得 \(\mathbb Z\xrightarrow{2}2\mathbb Z\) 就仍零调。

一般地，复形 \(C\) 的 \(\tau_{\le n}C\) 在 \(n\) 次取 \(\ker \mathrm{d}_C^n\)，较低次数不变，较高次数取零；它自然映入 \(C\)，并保留次数不大于 \(n\) 的上同调。\(\tau_{\ge m}C\) 在 \(m\) 次取 \(\operatorname{coker}\mathrm{d}_C^{m-1}\)，较高次数不变，较低次数取零；\(C\) 自然映到它，并保留次数不小于 \(m\) 的上同调，其余次数的上同调为零。右端取核，是为了保留原来成为余圈的条件；左端取余核，是为了仍将来自被删项的余边模掉。这些断言直接从核、像的定义核对，因而两种截断都保持拟同构。链映射在截口分别限制到核、下降到余核，给出截断后的链映射，并保持恒等与复合。

为了在同伦范畴中截断屋顶及其共同细化，还必须保证同伦的映射经截断后仍同伦。截口取核或余核依赖整个复形，不能仅由逐项加性函子保持同伦的结论代替这一步。设 \(f,g:C\to E\) 及 \(h^i:C^i\to E^{i-1}\) 满足
\[
 f^i-g^i=\mathrm{d}_E^{i-1}h^i+h^{i+1}\mathrm{d}_C^i.
\]
对 \(\tau_{\le n}\)，在 \(i<n\) 保留 \(h^i\)，在 \(i=n\) 将 \(h^n\) 限制到 \(\ker\mathrm{d}_C^n\)，在更高次数取零。截口上 \(h^{n+1}\mathrm{d}_C^n\) 消失，剩下的 \(\mathrm{d}_E^{n-1}h^n\) 取值于 \(\ker\mathrm{d}_E^n\)，所以同伦等式仍成立；在 \(n-1\) 次，限制后的 \(h^n\) 与诱导微分的复合仍为 \(h^n\mathrm{d}_C^{n-1}\)。对 \(\tau_{\ge m}\)，记截口的商映射为
\[
 q_C:C^m\to\operatorname{coker}\mathrm{d}_C^{m-1},
 \qquad q_E:E^m\to\operatorname{coker}\mathrm{d}_E^{m-1}.
\]
在 \(i>m+1\) 保留 \(h^i\)，在 \(i=m+1\) 取 \(q_Eh^{m+1}\)，在 \(i\le m\) 取零。由于 \(q_E\mathrm{d}_E^{m-1}=0\)，截口上的等式成为
\[
 (\tau_{\ge m}f-\tau_{\ge m}g)^m q_C
 =q_Eh^{m+1}\mathrm{d}_C^m,
\]
正是余核上所需的同伦等式；在 \(m+1\) 次，余核诱导的微分与 \(q_Eh^{m+1}\) 复合仍为 \(\mathrm{d}_E^mh^{m+1}\)。因此两种截断都给出 \(K(\mathcal A)\) 上的函子。

若 \(X,Y\) 的项在次数 \(n\) 以上为零，则屋顶 \(X\leftarrow U\to Y\) 可前复合 \(\tau_{\le n}U\to U\)。这里使用左屋顶，是因为截断的自然映射指向 \(U\)。因 \(H^{>n}(U)=0\)，这是拟同构，所以得到有上界的屋顶。若两个有上界的屋顶在 \(D(\mathcal A)\) 中相等，可增大 \(n\)，使这两个屋顶的全部对象在 \(n\) 以上为零，再对它们的共同细化施加 \(\tau_{\le n}\)。截断保持拟同构及 \(K(\mathcal A)\) 中的等式，且不改变两个原屋顶，因此得到有上界的共同细化。这同时证明 \(D^-\) 的满性和忠实性。对 \(D^+\)，截断的自然映射方向为 \(U\to\tau_{\ge m}U\)，因此使用分母在右的屋顶 \(X\to U\leftarrow Y\)，并把两个箭头后复合到截断对象。比较两个屋顶时，取它们全部对象的公共下界 \(m\)，截断共同细化，得到相同的结论。

若 \(X,Y\) 都只在 \([m,n]\) 内有项，先把左屋顶截到 \(\tau_{\le n}U\)，再取它的 \(\tau_{\ge m}\)。指向 \(X,Y\) 的复形映射在次数 \(m\) 都消去来自 \(m-1\) 的像，因此均通过该余核下降，得到有界屋顶；其左边仍为拟同构。比较两个有界屋顶时，选取包含它们全部对象的非零项的公共区间 \([m,n]\)，令 \(T=\tau_{\ge m}\tau_{\le n}\)。若共同细化中的等式为 \(s_1a=s_2b\)、\(f_1a=f_2b\)，施加 \(T\) 后仍有
\[
 T(s_1)T(a)=T(s_2)T(b),
 \qquad T(f_1)T(a)=T(f_2)T(b).
\]
这里 \(T\) 不改变两个原屋顶，而 \(T(a),T(b)\) 仍为拟同构，故所得细化属于 \(K^b(\mathcal A)\)。这证明 \(D^b\) 也全忠实。同样的截断论证在单向有界的屋顶范畴内直接进行，还给出 \(D^b(\mathcal A)\to D^+(\mathcal A),D^-(\mathcal A)\) 的全忠实性。

若 \(H^i(X)=0\) 对所有 \(i\notin[m,n]\) 成立，则有拟同构之字形
\[
 X\xleftarrow{\simeq}\tau_{\le n}X
 \xrightarrow{\simeq}\tau_{\ge m}\tau_{\le n}X.
\]
最右的复形只在 \([m,n]\) 内有项，故 \(X\) 在导出范畴中与项有界的复形同构。只有上界或下界时，分别用 \(\tau_{\le n}X\to X\) 或 \(X\to\tau_{\ge m}X\) 即可。反向由上同调保持拟同构显然。平移和锥保持相应的上同调有界性，故这些满子范畴对好三角封闭。
\end{proof}

\subsection{项有界与同调有界}
\label{b4:subsec:1703:02}

“项有界”是某个具体代表的性质，“上同调有界”是导出范畴中对象的性质。上命题说明后者允许更换成项有界的代表，却没有说明该代表还能同时逐项投射或逐项内射。为获得特殊项，可能必须重新加上一段无限长但正合的尾部。

\ProseParagraphBreak
例如投射消解 \(\cdots\to P_2\to P_1\to P_0\to M\to0\)，删去增广项后，把 \(P_i\) 放在次数 \(-i\)，得到复形 \(P\)。它只有 \(H^0(P)=M\)，因此代表 \(D^b\) 的对象；但它的项可能向负次数无限延伸。把正合尾部简单截掉通常会新增上同调；正确截断需要把端点项换成核或余核，而后者未必仍投射。

\ProseParagraphBreak
这也解释了为什么两个有界条件在导出函子公式中各有用处：上同调有界用于识别对象属于哪一类导出范畴，项及其投射、内射性质则用于保证逐项计算有效。证明中不可把二者无说明地互换。

\subsection{模范畴中的基本实例}
\label{b4:subsec:1703:03}

对域 \(k\)，每个向量空间复形都能非典范地分解为零微分的上同调复形与可缩复形的直和。逐次选 \(B^n\subseteq Z^n\subseteq C^n\) 的补空间，则 \(C^n\cong B^n\oplus H^n\oplus B^{n+1}\)，微分只在最后一项到下一次数第一项之间为恒等。于是 \(K(k\text{-}\mathbf{Mod})\to D(k\text{-}\mathbf{Mod})\) 已是等价，每个拟同构都是同伦等价。这里依赖所有向量空间短正合列可分裂。

\ProseParagraphBreak
整数上情况不同。两项复形 \(P=(\mathbb Z\xrightarrow{m}\mathbb Z)\) 位于次数负一、零，其中 \(m\ge2\)，自然映到 \(\mathbb Z/m\mathbb Z\)，是拟同构而非同伦等价。它给出屋顶
\[
 \mathbb Z/m\mathbb Z\xleftarrow{\simeq}P\longrightarrow\mathbb Z[1],
\]
右边的复形映射 \(f\) 在次数负一取恒等、在次数零取零，满足链映射条件。若 \(f\) 同伦于零，同伦只能在次数零给出 \(h^0:\mathbb Z\to\mathbb Z\)，而次数负一的同伦等式要求 \(h^0\circ m=1_{\mathbb Z}\)，即某个整数乘 \(m\) 等于一，这是不可能的。\cref{b4:prop:resolution-computes-derived-hom}将证明，源为这种有上界投射复形时，局部化不再使新的同伦类变成零；因此该屋顶表示非零导出态射。\secref{b4:sec:1704}还将严格算出 \(\operatorname{Hom}_D(\mathbb Z/m\mathbb Z,\mathbb Z[1])\cong\mathbb Z/m\mathbb Z\)，说明其倍数按模 \(m\) 区分。它在所有上同调对象之间诱导零映射，因为两端非零上同调位于不同次数。这说明“知道全部上同调对象和其上诱导的映射”仍不足以恢复导出范畴的态射。

% !TeX root = ../../Book4.tex
% 纲要标题：### 18.4 消解模型
% 纲要标签：b4:sec:1704
% 纲要位置：计划纲要.md，第 2941 行。

\section{消解模型}
\label{b4:sec:1704}

有上界投射复形到零调复形的映射，以及零调复形到有下界内射复形的映射，都在同伦范畴中为零。这个消失性质使导出态射可以回到同伦类计算，并把 Ext 识别为移位 Hom。为建立单向有界模型，下面使用\cref{b4:thm:ce-resolution}的消解存在性与比较，以及\cref{b4:cor:double-complex-comparison}。


% 纲要内容（保留纲要核验项目）：
%
% * 在足够多投射对象的情形，用有界上方的投射复形处理 \(D^-\)；在足够多内射对象的情形，用有界下方的内射复形处理 \(D^+\)
% * 导出范畴中计算态射
% * Ext 与导出 Hom 的一致性
% * \(D^b\) 中的对象一般仍可能需要向一端无限延伸的投射或内射消解；不能无条件把 \(D^b\) 等同于有界投射复形的同伦范畴
%
% ---
%

\subsection{有界方向下的投射与内射复形}
\label{b4:subsec:1704:01}

\begin{lemma}\label{b4:lem:bounded-resolution-nullhomotopy}
设 \(\mathcal A\) 为局部小交换范畴，\(A\) 为任意零调复形。若 \(P\) 是逐项投射的有上界复形，则每个复形映射 \(P\to A\) 零伦；若 \(I\) 是逐项内射的有下界复形，则每个复形映射 \(A\to I\) 零伦。这里对 \(A\) 没有有界性要求。
\end{lemma}
\begin{proof}
设 \(P^n=0\) 对 \(n>N\) 成立，\(f:P\to A\)。从 \(h^{N+1}=0\) 向下构造 \(h^n:P^n\to A^{n-1}\)。已构造高次数部分时，令 \(r^n=f^n-h^{n+1}\mathrm{d}_P^n\)。由高一度的同伦等式及 \(f\) 与微分相容，得 \(\mathrm{d}_A^nr^n=0\)。因为 \(A\) 正合，\(A^{n-1}\to Z^n(A)\) 满；\(P^n\) 投射使 \(r^n\) 提升为 \(h^n\)。于是 \(f^n=\mathrm{d}_A^{n-1}h^n+h^{n+1}\mathrm{d}_P^n\)。对每个固定次数只需有限次从 \(N\) 下降，所以构造定义了全部同伦。

内射情形把范畴改为 \(\mathcal A^{\mathrm{op}}\)，把次数反向：有下界内射复形变成有上界投射复形，正合性保留。前一构造便逐度向上延拓同伦，回到 \(\mathcal A\) 后得到第二个结论。
\end{proof}

\begin{theorem}\label{b4:thm:bounded-projective-model}
设 \(\mathcal A\) 为有足够投射对象的局部小交换范畴，\(\mathcal P\) 为其投射对象满子范畴。每个有上界复形 \(X\) 都有拟同构 \(P\to X\)，其中 \(P\in K^-(\mathcal P)\)，且自然函子
\(K^-(\mathcal P)\to D^-(\mathcal A)\) 是三角等价。
\end{theorem}
\begin{theorem}\label{b4:thm:bounded-injective-model}
设 \(\mathcal A\) 为有足够内射对象的局部小交换范畴，\(\mathcal I\) 为其内射对象满子范畴。每个有下界复形 \(X\) 都有拟同构 \(X\to I\)，其中 \(I\in K^+(\mathcal I)\)，且自然函子
\(K^+(\mathcal I)\to D^+(\mathcal A)\) 是三角等价。
\end{theorem}
\begin{proof}
先证内射模型的存在。对有下界 \(X\)，应用\cref{b4:thm:ce-resolution}已经构造并比较的 Cartan--Eilenberg 内射消解，取总复形 \(I\)。若 \(X^p=0\) 对 \(p<m\) 成立，双复形可取在 \(p\ge m,q\ge0\)，所以每个总次数只含有限多个项。有限直和内射，故 \(I\) 逐项内射且有下界。增广 \(X\to I\) 是拟同构，由逐列消解的正合性与\cref{b4:cor:double-complex-comparison}得到；有限对角线使这里不涉及无限总化。投射模型则把同一构造在反范畴中对次数反向应用：所得双复形位于 \(p\le N,q\le0\)，每条对角线仍有限，给出逐项投射的 \(P\to X\)。

直接在有界方向的屋顶范畴中证明全忠实性。记 \(K^-=K^-(\mathcal A)\)，以 \(Q^-:K^-\to D^-(\mathcal A)\) 表示局部化函子，令 \(P\in K^-(\mathcal P)\)。对 \(K^-\) 中的拟同构 \(s:U\to V\)，其锥及各平移零调，故由\cref{b4:lem:bounded-resolution-nullhomotopy}及三角的 Hom 正合列，
\[
 s_*:\operatorname{Hom}_{K^-}(P,U)
 \xrightarrow{\sim}\operatorname{Hom}_{K^-}(P,V).
\]
给定 \(D^-(\mathcal A)\) 中的屋顶 \(P\xleftarrow{s}U\xrightarrow{f}Y\)，取唯一的同伦类 \(a:P\to U\)，使 \(sa=1_P\)。局部化后 \(Q^-(a)=Q^-(s)^{-1}\)，所以屋顶等于 \(Q^-(fa)\)。若 \(Q^-(g)=0\)，屋顶相等判据给出拟同构 \(t:V\to P\)，使 \(gt=0\)；取 \(b:P\to V\) 满足 \(tb=1_P\)，便得 \(g=gtb=0\)。因此对每个有上界 \(Y\)，
\[
 \operatorname{Hom}_{K^-}(P,Y)
 \xrightarrow{\sim}\operatorname{Hom}_{D^-(\mathcal A)}(P,Y).
\]

内射情形记 \(K^+=K^+(\mathcal A)\)，以 \(Q^+:K^+\to D^+(\mathcal A)\) 表示局部化函子，令 \(I\in K^+(\mathcal I)\)。同一引理的内射部分与反变 Hom 正合列说明，对每个 \(K^+\) 中的拟同构 \(s:U\to V\)，
\[
 s^*:\operatorname{Hom}_{K^+}(V,I)
 \xrightarrow{\sim}\operatorname{Hom}_{K^+}(U,I).
\]
对分母在右的屋顶 \(X\xrightarrow{f}U\xleftarrow{s}I\)，取 \(a:U\to I\) 满足 \(as=1_I\)，屋顶便由 \(af\) 表示。若 \(Q^+(g)=0\)，右屋顶的相等判据给出拟同构 \(t:I\to V\)，使 \(tg=0\)；取 \(b:V\to I\) 满足 \(bt=1_I\)，得 \(g=btg=0\)。因此对每个有下界 \(X\)，
\[
 \operatorname{Hom}_{K^+}(X,I)
 \xrightarrow{\sim}\operatorname{Hom}_{D^+(\mathcal A)}(X,I).
\]
把 \(Y\) 也取为投射模型、把 \(X\) 也取为内射模型，即得两个模型函子的全忠实性；已证的拟同构提供本质满性。由于 \(\mathcal A\) 局部小，复形映射组成集合，取同伦类后仍为集合，故上述模型也说明 \(D^-(\mathcal A)\)、\(D^+(\mathcal A)\) 在原来的集合范围内局部小。

有限直和保持投射性和内射性，所以模型范畴均对平移和映射锥封闭。上述自然函子正是局部化函子 \(Q^-\)、\(Q^+\) 在模型满子范畴上的限制，与平移相容，并把每个标准锥三角送到相应有界导出范畴的好三角。因此所得等价是三角等价。
\end{proof}

\subsection{导出范畴中态射的计算}
\label{b4:subsec:1704:02}

\begin{proposition}\label{b4:prop:resolution-computes-derived-hom}
设 \(\mathcal A\) 为局部小交换范畴，并对局部化采用\subsecref{b4:subsec:1701:03}的大小约定，简记 \(K=K(\mathcal A)\)、\(D=D(\mathcal A)\)。若 \(\mathcal A\) 中的上链复形 \(P\) 逐项投射且有上界，\(I\) 逐项内射且有下界，则对任意 \(\mathcal A\) 中的上链复形 \(X,Y\)，自然映射
\[
 \operatorname{Hom}_{K}(P,Y)\xrightarrow{\sim}\operatorname{Hom}_{D}(P,Y),
 \qquad
 \operatorname{Hom}_{K}(X,I)\xrightarrow{\sim}\operatorname{Hom}_{D}(X,I)
\]
为双射。这里不要求 \(Y\) 有上界或 \(X\) 有下界。
\end{proposition}
\begin{proof}
若 \(s:U\to V\) 是拟同构，其锥及各平移零调。前一引理给出
\[
 \operatorname{Hom}_K(P,\operatorname{Cone}(s)[n])=0\qquad(n\in\mathbb Z).
\]
三角的 Hom 正合列因而说明 \(s_*:\operatorname{Hom}_K(P,U)\to\operatorname{Hom}_K(P,V)\) 为同构。

给定屋顶 \(P\xleftarrow{s}U\xrightarrow{f}Y\)，取唯一的同伦类 \(a:P\to U\)，使 \(sa=1_P\)。局部化后 \(Q(a)=Q(s)^{-1}\)，所以该屋顶等于 \(Q(fa)\)，证明满射。若 \(Q(g)=0\)，则由屋顶判据存在拟同构 \(t:V\to P\)，使 \(gt=0\)。同样取 \(b:P\to V\) 满足 \(tb=1\)，得 \(g=gtb=0\)，证明单射。

对内射版本，使用反变 Hom 正合列、引理第二项及分母在右的屋顶。于是 \(s^*:\operatorname{Hom}_K(V,I)\to\operatorname{Hom}_K(U,I)\) 为同构；完全相同的代表与消去论证给出第二个双射。
\end{proof}

若同时有 \(P\to X\) 和 \(Y\to I\)，便可任选
\[
 \operatorname{Hom}_D(X,Y)
 \cong\operatorname{Hom}_K(P,Y)
 \cong\operatorname{Hom}_K(P,I)
 \cong\operatorname{Hom}_K(X,I)
\]
来计算。中间比较来自自然的拟同构及上命题，故不同消解不会产生互不相关的答案。

\subsection{Ext 与导出 Hom 的一致性}
\label{b4:subsec:1704:03}

\begin{theorem}\label{b4:thm:derived-ext-identification}
设 \(R\) 为含幺环，\(M,N\) 为左 \(R\)-模，视为次数零复形。则对 \(n\ge0\) 有自然同构
\[
 \operatorname{Ext}_R^n(M,N)\cong
 \operatorname{Hom}_{D(R\text{-}\mathbf{Mod})}(M,N[n]).
\]
对 \(n<0\)，右端为零；\(n=0\) 时该同构是通常的 Hom 识别。若 \(\mathcal A\) 是有足够内射对象的局部小交换范畴，\(M,N\in\mathcal A\)，则在固定的集合大小约定下也有 \(\operatorname{Ext}_{\mathcal A}^n(M,N)\cong\operatorname{Hom}_{D(\mathcal A)}(M,N[n])\)，其中 Ext 用右导出函子定义，负次数的右端同样为零。
\end{theorem}
\begin{proof}
取 \(N\to I\) 的非负次数内射消解。\(I[n]\) 仍逐项内射且有下界，所以前一命题及 Hom 复形公式给出
\[
 \operatorname{Hom}_D(M,N[n])\cong\operatorname{Hom}_K(M,I[n])
 \cong H^n\operatorname{Hom}^{\bullet}(M,I)
 =\operatorname{Ext}^n_R(M,N).
\]
因为 \(M\) 仅在零次非零，此处的 Hom 微分就是后复合 \(\mathrm{d}_I\)，与\chapref{b4:ch:08}内射定义一致。负次数的 Hom 复形项为零，所以所得群为零。比较映射唯一到同伦，故同构对两个变量自然。

也可以取 \(P\to M\) 的非正次数投射消解，得到 \(H^n\operatorname{Hom}^{\bullet}(P,N)\)。须注意\chapref{b4:ch:08}投射计算采用 \(\delta f=f\mathrm{d}_P\)，而内部 Hom 在次数 \(n\) 的微分为 \((-1)^{n+1}f\mathrm{d}_P\)。每个次数乘 \((-1)^{n(n+1)/2}\) 给出前者到后者的复形同构。连同\chapref{b4:ch:08}的平衡比较，即与上面的内射识别相容；不能只换微分而保留所有余圈代表的符号。
\end{proof}

因此 \(\operatorname{Hom}_D(\mathbb Z/m,\mathbb Z[1])\cong\mathbb Z/m\)，而 \(\operatorname{Hom}_D(\mathbb Z/m,\mathbb Z)=0\)。前一群记录扩张，后一群只记录普通同态。按照\chapref{b4:ch:09}用第一变量的内射连接同态识别扩张类的约定，一次 Yoneda 扩张对应其好三角的最后箭头 \(\delta\)。

\ProseParagraphBreak
可以直接核对这一符号。给定 \(0\to N\xrightarrow{i}E\xrightarrow{q}M\to0\)，取内射消解 \(j:N\to I\)，并延拓为 \(a:E\to I^0\)。由于 \(\mathrm{d}_Ia\) 消去 \(N\)，它唯一下降为余圈 \(f:M\to I^1\)，满足 \(fq=\mathrm{d}_Ia\)；这就是\chapref{b4:ch:09}的内射连接类。在 \(\operatorname{Cone}(i)\to I[1]\) 上，\(j[1]p_i-fr\) 的零伦取 \(h^0=a\)：次数负一的等式是 \(ai=j\)，次数零的等式是 \(-\mathrm{d}_Ia=-fq\)。所以 \(Q(j[1])\delta=Q(f)\)。若改用无符号投射余圈，必须连同前述每度重标识和平衡比较一起换算，不能据次数一的重标识孤立地改变这个扩张对应。

\subsection{有界导出对象与单向无界消解}
\label{b4:subsec:1704:04}

有界导出对象未必有有界投射代表。取域 \(k\) 及交换环 \(R=k[\varepsilon]/(\varepsilon^2)\)，令 \(k=R/(\varepsilon)\)。复形
\[
 \cdots\xrightarrow{\varepsilon}R\xrightarrow{\varepsilon}R
 \xrightarrow{\varepsilon}R\longrightarrow k\longrightarrow0
\]
是自由消解，因为 \(\ker(\varepsilon:R\to R)=\varepsilon R\)。施加 \(\operatorname{Hom}_R(-,k)\) 后所有微分为零，故
\[
 \operatorname{Ext}_R^n(k,k)\cong k\qquad(n\ge0).
\]
若 \(k\) 在导出范畴中同构于一个有界投射复形 \(Q\)，则前面的态射计算给出 \(\operatorname{Ext}_R^n(k,k)\cong\operatorname{Hom}_K(Q,k[n])\)。当 \(n\) 足够大时，\(k[n]\) 的唯一非零项已移到 \(Q\) 的范围以外，后者为零，矛盾。

\ProseParagraphBreak
所以 \(k\in D^b(R\text{-}\mathbf{Mod})\)，但不能用 \(K^b(R\text{-}\mathbf{Proj})\) 中的对象代表。与此相对，\(\mathbb Z/m\) 有长度一的自由消解，可以由有界投射复形代表。是否允许这样的有限模型，与环和模的同调维数有关；\chapref{b4:ch:20}将把有限生成投射项构成的有限模型称为完美复形。

\ProseParagraphBreak
本章的结论不包括“逐项投射的任意无界复形都能计算导出态射”。向下构造零伦时必须有起始次数；没有这个起点，逐次投射性本身不能启动归纳。无界情形所需的更强条件在\secref{b4:sec:1804}介绍。


\begin{chaptersummary}
拟同构在同伦范畴中满足交换和消去条件，因此导出态射可写成屋顶，屋顶的相等由共同细化判断。这个计算既证明局部化的泛性质，也允许把三角公理从同伦范畴传下来。由此，一般复形短正合列都能给出导出范畴中的好三角。

\ProseParagraphBreak
截断负责把上同调有界的对象换成项有界的代表，投射和内射替换负责把态射计算化为同伦计算。它们解决的问题不同：有界对象可能仍需向一端无限延伸的特殊消解。Ext 因而成为到平移对象的态射群，而这些态射又保留了普通上同调对象之间映射看不到的扩张信息。
\end{chaptersummary}

\BookExercises

\begin{exercise}\label{b4:ex:17-cyclic-ext}
设 \(m,n\ge1\)。计算 \(\operatorname{Hom}_{D(\mathbf{Ab})}(\mathbb Z/m\mathbb Z,\mathbb Z/n\mathbb Z[r])\) 对所有整数 \(r\) 的值，并写出 \(r=0,1\) 的核与余核表达。
\end{exercise}
\begin{solution}
用 \(\mathbb Z\xrightarrow{m}\mathbb Z\) 替换第一变量，Hom 复形仅在零、一次非零，两项都是 \(\mathbb Z/n\mathbb Z\)，微分为负乘 \(m\)。故零次为 \((\mathbb Z/n\mathbb Z)[m]\)，一次为 \((\mathbb Z/n\mathbb Z)/m(\mathbb Z/n\mathbb Z)\)，二者均同构于 \(\mathbb Z/\gcd(m,n)\mathbb Z\)，其余为零。零次与一次的具体子群、商群不同，不把这种同构类型误认成同一内部对象。
\end{solution}

\begin{exercise}\label{b4:ex:17-roof-localize-zero}
设 \(f:X\to Y\) 为复形映射。若它经某个零调复形分解，证明 \(Q(f)=0\)。反之，若 \(Q(f)=0\)，证明 \(f\) 在同伦范畴中经某个零调复形分解。
\end{exercise}
\begin{solution}
正向因为中间对象在 \(D\) 中为零。反向由屋顶判据取拟同构 \(s:U\to X\)，使 \(fs=0\)。补成三角 \(U\xrightarrow{s}X\to C_s\to U[1]\)。反变 Hom 正合性使 \(f\) 经 \(C_s\) 分解；拟同构的锥零调，故满足要求。
\end{solution}

\begin{exercise}\label{b4:ex:17-roof-two-denominators}
设两个屋顶 \((s,f)\)、\((s,g)\) 具有同一个左分母。证明它们在导出范畴中相等，当且仅当存在拟同构 \(u\) 使 \(fu=gu\)。
\end{exercise}
\begin{solution}
若有这样的 \(u\)，共同细化立即给出相等。反向，乘上可逆的 \(Q(s)\) 得 \(Q(f)=Q(g)\)，再用普通态射的屋顶相等判据。也可以直接从共同细化 \(sa=sb\)、\(fa=gb\) 出发，先用拟同构 \(s\) 的消去条件进一步使 \(a=b\)。
\end{solution}

\begin{exercise}\label{b4:ex:17-field-two-degrees}
设域 \(k\) 上的复形 \(X\) 仅有 \(H^{-1}(X)=k^2,H^0(X)=k\)，\(Y\) 仅有 \(H^{-1}(Y)=k,H^0(Y)=k^3\)。计算 \(\dim_k\operatorname{Hom}_D(X,Y[r])\)。所得维数是否依赖分裂的选择？
\end{exercise}
\begin{solution}
向量空间复形分解为上同调与可缩部分，故计算可在零微分复形上进行。\(r=0\) 时维数为 \(2\cdot1+1\cdot3=5\)；\(r=1\) 时从 \(X^{-1}=k^2\) 到 \(Y[1]^{-1}=k^3\)，维数为六；\(r=-1\) 时从 \(X^0=k\) 到 \(Y[-1]^0=k\)，维数为一；其余为零。分裂不典范，但 Hom 群本身及其维数由导出对象确定，不依赖选择。
\end{solution}

\begin{exercise}\label{b4:ex:17-injective-cyclic}
设整数 \(m\ge2\)。用内射消解
\[
 0\to\mathbb Z\to\mathbb Q\to\mathbb Q/\mathbb Z\to0
\]
计算 \(\operatorname{Hom}_D(\mathbb Z/m\mathbb Z,\mathbb Z[1])\)，并描述一个生成元。
\end{exercise}
\begin{solution}
内射复形取 \(I^0=\mathbb Q,I^1=\mathbb Q/\mathbb Z\)。Hom 的零次项为零，一次项为 \((\mathbb Q/\mathbb Z)[m]\)。因此所求群为 \(\mathbb Z/m\mathbb Z\)，生成元由 \(\bar1\mapsto1/m+\mathbb Z\) 给出。与投射屋顶的具体符号比较时须用正文的每度重标识，不能仅比较写下的两个数值。
\end{solution}

\begin{exercise}\label{b4:ex:17-projective-dimension}
设 \(R\) 为环，\(M\) 为左 \(R\)-模，\(d\ge0\)。证明 \(\operatorname{pd}_R M\le d\) 当且仅当对所有左模 \(N\) 都有 \(\operatorname{Hom}_D(M,N[d+1])=0\)。
\end{exercise}
\begin{solution}
由 Ext 识别，条件等价于所有 \(\operatorname{Ext}_R^{d+1}(M,N)=0\)。若有长度 \(d\) 的投射消解，消解计算直接给出消失。反向取任意投射消解，维数平移给出 \(\operatorname{Ext}_R^1(\Omega^dM,N)=0\) 对所有 \(N\) 成立，故 \(\Omega^dM\) 投射，将原消解在这里截短即可。\(d=0\) 时用一次 Ext 的投射判据。
\end{solution}

\begin{exercise}\label{b4:ex:17-cubic-truncated}
令 \(R=k[t]/(t^3)\)、\(k=R/(t)\)。构造 \(k\) 的周期投射消解，并计算它到各次移位 \(k[r]\) 的导出态射群。由此说明 \(k\) 不能由有界投射复形代表。
\end{exercise}
\begin{solution}
消解交替使用乘 \(t\) 与乘 \(t^2\)，靠近增广的箭头为乘 \(t\)。因为 \(\ker t=(t^2)\)、\(\ker t^2=(t)\)，它正合。施加 \(\operatorname{Hom}_R(-,k)\) 后微分均为零，故所求群在 \(r\ge0\) 为 \(k\)，在 \(r<0\) 为零。若存在有界投射代表，则到 \(k[r]\) 的同伦态射在充分大的 \(r\) 必为零，与计算矛盾。
\end{solution}

\begin{exercise}\label{b4:ex:17-cone-localized}
设 \(s:X\to Y\) 为拟同构，\(f:Y\to Z\) 为复形映射。证明 \(\operatorname{Cone}(fs)\) 与 \(\operatorname{Cone}(f)\) 在导出范畴中同构；这种结论在同伦范畴中是否总成立？
\end{exercise}
\begin{solution}
对严格方块 \((s,1_Z)\) 取锥映射。上同调长正合列与五引理使该锥映射拟同构，故在 \(D\) 中同构。若取零调但不可缩的 \(X\)，令 \(s:X\to0\)、\(f:0\to0\)，则两个锥分别为 \(X[1]\) 和零，在 \(K\) 中不同构。因此不能去掉局部化。
\end{solution}

\begin{exercise}\label{b4:ex:17-cohomology-detect}
设 \(R\) 为环，\(X\) 为任意左模复形。证明 \(\operatorname{Hom}_{D(R)}(R[-n],X)\cong H^n(X)\)。由此证明若所有这些群为零，则 \(X=0\)。
\end{exercise}
\begin{solution}
\(R[-n]\) 是只有一项的投射复形，可用同伦类计算导出态射。映射由 \(X^n\) 的一个余圈决定，两个映射同伦恰当其差为余边，故所求群为 \(H^n(X)\)。所有群为零表示 \(X\) 零调，等价于 \(X\) 在导出范畴中为零。
\end{solution}

\begin{exercise}\label{b4:ex:17-split-cohomology-not-split}
设整数 \(m\ge2\)，考虑由非分裂短正合列 \(0\to\mathbb Z\xrightarrow{m}\mathbb Z\to\mathbb Z/m\mathbb Z\to0\) 给出的好三角。证明其最后箭头非零，尽管它诱导的每个上同调对象之间的映射都为零。
\end{exercise}
\begin{solution}
若最后箭头为零，由分裂三角判据，\(\mathbb Z\to\mathbb Z/m\mathbb Z\) 在导出范畴中有截面。取 \(H^0\) 得交换群满射的截面，与 \(\mathbb Z\) 无有限阶元素矛盾。所以最后箭头非零。它从 \(\mathbb Z/m\mathbb Z\) 指向 \(\mathbb Z[1]\)，两端非零上同调的次数不同，因此每个 \(H^n\) 上的映射均零。
\end{solution}

\begin{exercise}\label{b4:ex:17-quasi-between-models}
设 \(P,Q\) 为逐项投射且有上界的复形，\(s:P\to Q\) 为拟同构。证明它已经是同伦等价，并说明不能仅以“拟同构的锥正合”结束证明。
\end{exercise}
\begin{solution}
锥 \(C_s\) 也逐项投射且有上界。把零伦引理应用于其恒等映射 \(C_s\to C_s\)，因目标正合，得到锥可缩。三角 Hom 正合列说明 \(s\) 在 \(K\) 中同构，即同伦等价。正合本身不足以推出可缩，这里必须同时使用锥逐项投射及有上界。
\end{solution}

\begin{exercise}\label{b4:ex:17-truncation-object}
设复形 \(X\) 只有次数 \(n\) 的上同调可能非零。构造 \(D(\mathcal A)\) 中 \(X\cong H^n(X)[-n]\) 的同构，不选各项短正合列的分裂。
\end{exercise}
\begin{solution}
取 \(\tau_{\le n}X\to X\)，它是拟同构。再将 \(\tau_{\le n}X\) 的次数 \(n\) 映到 \(H^n(X)\)，其他次数映到零，得到复形映射 \(\tau_{\le n}X\to H^n(X)[-n]\)。它在唯一可能非零的上同调上为恒等，因而也是拟同构。这两个箭头给出所需同构；全部映射只用核和商，无须分裂。
\end{solution}

\begin{exercise}\label{b4:ex:17-localization-functor}
设 \(F:\mathcal A\to\mathcal B\) 为交换范畴间的正合加性函子。证明逐项 \(F\) 诱导三角函子 \(D(\mathcal A)\to D(\mathcal B)\)，并说明为何仅加性不足。
\end{exercise}
\begin{solution}
正合性使 \(H^n(FC)\cong F(H^nC)\)，故逐项 \(F\) 保持拟同构，由局部化泛性质下降。加性使它保持有限双积和平移的符号，也保持锥矩阵，所以保持好三角。仅加性时拟同构可能不再保持，例如把 \(\mathbb Z\xrightarrow{2}\mathbb Z\to\mathbb Z/2\mathbb Z\) 的投射替换与 \(\mathbb Z/2\mathbb Z\) 张量会新增负一次上同调，因而不能直接下降。
\end{solution}

\begin{exercise}\label{b4:ex:17-abelian-embedding}
设 \(\mathcal A\) 为有足够内射对象的局部小交换范畴。证明将对象放在次数零给出全忠实函子 \(\mathcal A\to D^b(\mathcal A)\)，但这不使 \(D^b(\mathcal A)\) 自动成为交换范畴。
\end{exercise}
\begin{solution}
用有下界内射消解在 \(D^+(\mathcal A)\) 中计算，次数零给出 \(\operatorname{Hom}_{D^+}(M,N)=\operatorname{Hom}_{\mathcal A}(M,N)\)。有界嵌入的全忠实性使它也等于 \(D^b\) 中的 Hom，因此对象放在次数零的函子全忠实。

若 \(D^b(\mathbf{Ab})\) 同时为交换范畴，则任意态射 \(f:X\to Y\) 可分解为满态射 \(e:X\to\operatorname{im}f\) 和单态射 \(m:\operatorname{im}f\to Y\)。由\cref{b4:prop:triangle-mono-splits}，可取 \(a:\operatorname{im}f\to X\)、\(b:Y\to\operatorname{im}f\)，满足 \(ea=1\)、\(bm=1\)。于是 \(t=ab:Y\to X\) 满足
\[
 ftf=(me)(ab)(me)=m(ea)(bm)e=me=f.
\]
对次数零的乘二态射 \(\mathbb Z\to\mathbb Z\)，全忠实性使 \(t\) 仍由一个整数表示，因而要求 \(4t=2\)，矛盾。可见原交换范畴作为满子范畴嵌入，并不使核余核的性质在较大范畴中自动成立。
\end{solution}

% !TeX root = ../../Book4.tex
% 纲要标题：## 第19章　全导出函子、导出张量积与导出 Hom
% 纲要标签：b4:ch:18
% 纲要位置：计划纲要.md，第 3711 行。

\chapter{全导出函子、导出张量积与导出 Hom}
\label{b4:ch:18}
\BookChapterNumber{b4:ch:18}

\begin{chapterabstract}
本章以合适的复形替换构造全导出函子、导出张量积和导出 Hom，在明确有界条件与模的侧别后证明结合律、伴随和基变换，并讨论无界替换。
\end{chapterabstract}

\begin{learninggoals}
\begin{itemize}
\item 区分全导出函子、各次导出函子与超上同调，写出替换方向及泛性质中的比较态射。
\item 按右模与左模的侧别计算导出张量，并核对总次数及移位符号。
\item 用投射或内射替换计算导出 Hom，证明其与导出张量的伴随。
\item 分别检验基变换中的平坦性和有限生成性，识别普通基变换的失效。
\item 区分逐项投射、内射、平坦与相应的 K-性质。
\end{itemize}
\end{learninggoals}

若直接计算 \(\mathbb Z/2\mathbb Z\otimes_{\mathbb Z}\mathbb Z/2\mathbb Z\)，只能得到一个 \(\mathbb Z/2\mathbb Z\)。但将第一个模换成自由消解 \(0\rightarrow\mathbb Z\xrightarrow{2}\mathbb Z\rightarrow0\)，再与第二个模张量，乘 \(2\) 的微分变为零，所得复形在两个相邻次数各有一个 \(\mathbb Z/2\mathbb Z\) 的同调。直接张量丢掉的信息因此不是凭空添上的；它原本就在消解的另一度中。要让更换消解成为合法计算，还须说明何种复形可以替换、替换的选择是否改变结果，以及 Hom 的导出版本怎样与张量相配。

\ProseParagraphBreak
本章使用\chapref{b4:ch:17}中导出范畴的局部化和投射、内射模型，也使用\cref{b4:def:hom-complex}、\cref{b4:def:tensor-complex}的微分约定。与复合导出相关的谱序列已在\cref{b4:thm:grothendieck-spectral-sequence}建立。前三节的主要计算分别在有上界或有下界的范围进行；\subsecref{b4:subsec:2002:03}的无界余积与紧性讨论会调用\secref{b4:sec:1804}中的半自由替换。

\ProseParagraphBreak
对应阅读是 Weibel 的《An Introduction to Homological Algebra》\cite[Sections 10.5--10.8]{Weibel1994HomologicalAlgebra}：全导出构造见 Existence Theorem 10.5.6，张量与 Tor 见 Theorem 10.6.3 和 Corollary 10.6.4，导出 Hom 见 Theorem 10.7.4，复合条件见 Composition Theorem 10.8.2。本书在模范畴中展开基变换算例，并自行给出\secref{b4:sec:1804}所需的半自由替换构造。

% !TeX root = ../../Book4.tex
% 纲要标题：### 19.1 全导出函子
% 纲要标签：b4:sec:1801
% 纲要位置：计划纲要.md，第 3713 行。

\section{全导出函子}
\label{b4:sec:1801}

一个加性函子不一定保持拟同构，因而不能直接作用于导出态射。先选择能被它正确处理的复形模型，才能把导出群、自然变换与复合一起提升到导出范畴。

% 纲要内容（仅作写作计划，尚非教材正文）：
%
% * RF 和 LF 的构造及假设
% * 对象级导出函子与复形级导出的关系
% * 复合与自然变换
%

\subsection{\texorpdfstring{\(RF\) 与 \(LF\) 的构造条件}{RF 与 LF 的构造条件}}
\label{b4:subsec:1801:01}

对一个左正合函子逐次取导出，得到一列函子；但连接同态、自然变换及它们的复合仍需另外追踪。导出范畴允许保留产生这些群的整个复形。困难是原来的加性函子不一定保持拟同构，因而不能直接作用于分式表示。

\begin{definition}\label{b4:def:total-derived-functor}
设 \(\mathcal A,\mathcal B\) 是局部小交换范畴，\(F:\mathcal A\to\mathcal B\) 是加性函子，并固定导出范畴的集合大小约定。记 \(q_{\mathcal A}:K^+(\mathcal A)\to D^+(\mathcal A)\)、\(q_{\mathcal B}:K^+(\mathcal B)\to D^+(\mathcal B)\) 为局部化函子。三角函子
\(RF:D^+(\mathcal A)\to D^+(\mathcal B)\) 连同自然变换
\[
\eta:q_{\mathcal B}F\longrightarrow RFq_{\mathcal A}
\]
若满足下列泛性质，就称为 \(F\) 的\term[quan-you-daochu-hanzi]{全右导出函子}[total right derived functor]：对每个三角函子 \(T:D^+(\mathcal A)\to D^+(\mathcal B)\) 及自然变换 \(\xi:q_{\mathcal B}F\to Tq_{\mathcal A}\)，唯一存在自然变换 \(\theta:RF\to T\)，使 \(\xi=(\theta q_{\mathcal A})\eta\)。这里 \(F\) 在复形和同伦类上逐项作用。

对三角函子 \(LF:D^-(\mathcal A)\to D^-(\mathcal B)\)，将 \(q_{\mathcal A},q_{\mathcal B}\) 改为相应有上界范畴的局部化函子，并把自然变换方向改为
\(\epsilon:LFq_{\mathcal A}\to q_{\mathcal B}F\)。若对每个三角函子 \(T:D^-(\mathcal A)\to D^-(\mathcal B)\)，每个自然变换 \(Tq_{\mathcal A}\to q_{\mathcal B}F\) 都唯一经 \(T\to LF\) 分解，则称 \(LF\) 连同 \(\epsilon\) 为\term[quan-zuo-daochu-hanzi]{全左导出函子}[total left derived functor]。
\end{definition}
\symidx{\(RF,LF\)：全右导出函子与全左导出函子}

定义中的两个自然变换方向不同。内射替换给出 \(X\to I_X\)，因而右导出的比较来自 \(F(X)\to F(I_X)\)；投射替换给出 \(P_X\to X\)，因而左导出的比较来自 \(F(P_X)\to F(X)\)。加性函子未必保持这些拟同构，所以比较态射未必是同构。全导出函子修正的正是逐项 \(F\) 不能在导出范畴中直接使用的问题，不能要求修正前后的值总是相同。\cref{b4:prop:additive-right-derived-zero}将给出一个非零加性函子，其全右导出却为零。

\begin{theorem}\label{b4:thm:total-derived-functors}
设 \(\mathcal A,\mathcal B\) 是局部小交换范畴，\(F:\mathcal A\to\mathcal B\) 是加性函子，并在固定的集合大小约定下取导出范畴。
若 \(\mathcal A\) 有足够内射对象，则全右导出函子 \(RF:D^+(\mathcal A)\to D^+(\mathcal B)\) 存在，并对每个有下界上链复形 \(X\) 由
\[
RF(q_{\mathcal A}X)\cong q_{\mathcal B}F(I_X),\qquad X\xrightarrow{\simeq}I_X,
\]
计算，其中 \(I_X\) 是逐项内射的有下界复形。若 \(\mathcal A\) 有足够投射对象，则全左导出函子 \(LF:D^-(\mathcal A)\to D^-(\mathcal B)\) 存在，并对每个有上界上链复形 \(X\) 由
\[
LF(q_{\mathcal A}X)\cong q_{\mathcal B}F(P_X),\qquad P_X\xrightarrow{\simeq}X,
\]
计算，其中 \(P_X\) 是逐项投射的有上界复形。结果及其比较映射在导出范畴中不依赖替换的选择。
\end{theorem}
\begin{proof}
\cref{b4:thm:bounded-injective-model}给出的三角等价
\[
D^+(\mathcal A)\simeq K^+(\operatorname{Inj}\mathcal A)
\]
说明每个导出对象都有有下界的内射模型，而内射模型之间的导出态射可以按链同伦类计算。因此可先把导出对象换成内射模型，在模型上逐项作用 \(F\)，再进入 \(D^+(\mathcal B)\)。
记 \(\mathcal I=K^+(\operatorname{Inj}\mathcal A)\)，以
\(j:\mathcal I\to K^+(\mathcal A)\) 表示包含函子，并记此三角等价为
\(E=q_{\mathcal A}j:\mathcal I\to D^+(\mathcal A)\)。选择三角准逆
\(U:D^+(\mathcal A)\to\mathcal I\)，并将等价的自然同构取为伴随等价 \(U\dashv E\) 的单位与余单位：
\[
\rho:1_{D^+(\mathcal A)}\xrightarrow{\sim}EU,
\qquad
\kappa:UE\xrightarrow{\sim}1_{\mathcal I}.
\]
它们满足 \(E(\kappa_I)\rho_{E(I)}=1_{E(I)}\) 及
\(\kappa_{U(D)}U(\rho_D)=1_{U(D)}\)。定义
\(RF=q_{\mathcal B}FjU\)。加性函子保持有限双积、移位及锥的微分矩阵，所以逐项 \(F\) 是三角函子，所得复合也保持好三角。

对 \(X\in K^+(\mathcal A)\)，记 \(I_X=jU(q_{\mathcal A}X)\)。上一模型定理证明中的双射
\[
\operatorname{Hom}_{K^+(\mathcal A)}(X,I_X)
\xrightarrow{\sim}
\operatorname{Hom}_{D^+(\mathcal A)}(q_{\mathcal A}X,E U(q_{\mathcal A}X))
\]
将 \(\rho_{q_{\mathcal A}X}\) 唯一提升为同伦类 \(i_X:X\to I_X\)。因为其局部化是同构，\(i_X\) 的复形映射代表是拟同构。\(\rho\) 的自然性及上述双射的单射性，使 \(i\) 成为同伦范畴上的自然变换。这并不要求预先自然地选定严格链映射代表：更换代表只会改变一个链同伦，而加性 \(F\) 保持链同伦。因此
\(\eta_X=q_{\mathcal B}F(i_X)\) 不依赖代表，给出所需自然比较。

给定 \(T,\xi\)，先看泛分解等式怎样决定所需变换。第一条三角恒等式及提升唯一性给出 \(i_{jI}=j(\kappa_I^{-1})\)，第二条三角恒等式因而说明
\[
RF(\rho_D)=q_{\mathcal B}Fj(U(\rho_D))
=q_{\mathcal B}Fj(\kappa_{U(D)}^{-1})
=\eta_{jU(D)}.
\]
若 \(\theta:RF\to T\) 满足泛分解等式，沿 \(\rho_D:D\xrightarrow{\sim}EU(D)\) 的自然性便要求
\[
T(\rho_D)\theta_D
=\theta_{EU(D)}RF(\rho_D)
=\theta_{EU(D)}\eta_{jU(D)}
=\xi_{jU(D)}.
\]
由于 \(T(\rho_D)\) 可逆，\(\theta_D\) 只能取为
\[
\theta_D=(T(\rho_D))^{-1}\xi_{jU(D)}:
RF(D)\longrightarrow T(D).
\]
这里 \(\xi_{jU(D)}\) 的目标是 \(T(EU(D))\)；它与 \(T(\rho_D)\) 的逆可以复合。\(\xi\) 与 \(\rho\) 的自然性说明 \(\theta\) 自然，而
\[
\theta_{q_{\mathcal A}X}\eta_X
=(T(\rho_{q_{\mathcal A}X}))^{-1}
T(q_{\mathcal A}i_X)\xi_X=\xi_X.
\]
这证明存在性；前面的强制公式同时证明唯一性。

左导出采用 \(P_X\to X\) 的投射替换，所以比较方向是 \(LF\to F\)。给定的变换是 \(\xi:Tq_{\mathcal A}\to q_{\mathcal B}F\)，要找的是 \(\theta:T\to LF\)，使 \(\epsilon_X\theta_{q_{\mathcal A}X}=\xi_X\)；它的泛分解方向与右导出相反。记
\(\mathcal P=K^-(\operatorname{Proj}\mathcal A)\)，令
\(j^-:\mathcal P\to K^-(\mathcal A)\) 为包含，
\(E^-=q_{\mathcal A}j^-\)。由\cref{b4:thm:bounded-projective-model}选择三角准逆
\(V:D^-(\mathcal A)\to\mathcal P\)，以及伴随等价 \(E^-\dashv V\) 的单位、余单位
\[
\nu:1_{\mathcal P}\xrightarrow{\sim}VE^-,
\qquad \lambda:E^-V\xrightarrow{\sim}1_{D^-(\mathcal A)}.
\]
投射模型的 Hom 双射将 \(\lambda_{q_{\mathcal A}X}\) 唯一提升为自然的同伦类
\(p_X:j^-V(q_{\mathcal A}X)\to X\)。定义
\(LF=q_{\mathcal B}Fj^-V\)、\(\epsilon_X=q_{\mathcal B}F(p_X)\)。三角恒等式给出
\(LF(\lambda_D)=\epsilon_{j^-V(D)}\)。沿 \(\lambda_D:E^-V(D)\xrightarrow{\sim}D\) 的自然性和泛分解等式因而强制
\[
\theta_D T(\lambda_D)
=LF(\lambda_D)\theta_{E^-V(D)}
=\epsilon_{j^-V(D)}\theta_{E^-V(D)}
=\xi_{j^-V(D)}.
\]
所以所需变换在导出对象 \(D\) 上为
\[
\theta_D=\xi_{j^-V(D)}(T(\lambda_D))^{-1}:
T(D)\longrightarrow LF(D).
\]
自然性及 \(\epsilon_X\theta_{q_{\mathcal A}X}=\xi_X\) 由 \(\xi\) 与 \(\lambda\) 的自然性方块得到；前面的强制公式证明唯一性。

任意另一内射或投射替换都经模型的 Hom 双射得到与给定替换相容的唯一同伦比较；它在局部化后为同构，因模型等价全忠实而是同伦等价。逐项 \(F\) 保持同伦等价，故这些替换计算的对象典范同构于上述函子的值。泛性质也使不同模型选择所得函子连同比较变换唯一自然同构。
\end{proof}

这里构造存在只需加性，不要求左正合或右正合。正合性将用于把零次同调重新识别为 \(F\) 本身。以上泛性质及构造可对照 Weibel\cite[Definition 10.5.1, Existence Theorem 10.5.6]{Weibel1994HomologicalAlgebra}。

\subsection{对象级导出与复形级导出}
\label{b4:subsec:1801:02}

把对象 \(A\in\mathcal A\) 看作零次复形，替换 \(A\to I^\bullet\) 就是通常的内射消解。因此，当 \(F\) 左正合且 \(\mathcal A\) 有足够内射对象时，
\[
H^n(RF(A))\simeq R^nF(A)\quad(n\ge0),\qquad
H^n(RF(A))=0\quad(n<0).
\]
尤其 \(H^0(RF(A))=F(A)\)。右正合函子在有足够投射对象时给出
\(H^{-n}(LF(A))=L_nF(A)\)。负号来自本书始终使用升次微分，而通常的投射消解用非负链次数编号。

\ProseParagraphBreak
若 \(F\) 左正合、\(\mathcal A\) 有足够内射对象，而 \(X\) 为有下界复形，则 \(H^n(RF(X))\) 就是\cref{b4:def:hypercohomology}中的 \(\mathbb H^n(F,X)\)：两种计算选择同一内射复形后相同。\cref{b4:thm:hypercohomology-spectral-sequences}的第二套谱序列于是写成
\[
E_2^{p,q}=R^pF(H^q(X))
\Longrightarrow H^{p+q}(RF(X)).
\]
这里 \(p\ge0\)，\(q\) 有下界，每个总次数的滤过有限。即使知道所有第二页数据，仍须经过后续微分，并由终页数据得到目标的逐层商，不能一般地把它们直接相加。\(R^nF(A)\) 是单个对象的第 \(n\) 次导出，\(RF(X)\) 是整个复形的导出对象，\(H^n(RF(X))\) 才是该对象的第 \(n\) 次上同调。特别地，\(RF(X)\) 不能只由 \(R^nF(H^0X)\) 代替；即使 \(F\) 是恒等函子，\(RF(X)=X\) 仍保留每个 \(H^nX\)。

下文沿用\cref{b4:def:functor-acyclic}的“\(F\)-零调对象”：当 \(F\) 左正合时，它只要求 \(R^jF(Q)=0\) 对 \(j>0\) 成立；“无上同调对象”是同一性质的别称。它与整个复形的零调性是两个不同条件。

\begin{proposition}\label{b4:prop:total-derived-acyclic-complex}
设 \(F:\mathcal A\to\mathcal B\) 是局部小交换范畴间的加性左正合函子，\(\mathcal A\) 有足够内射对象。若 \(X\) 是 \(\mathcal A\) 中的有下界上链复形，且每项均满足
\(R^jF(X^n)=0\) 对所有 \(j>0\) 成立，则自然映射 \(F(X)\to RF(X)\) 是拟同构。若 \(F\) 为加性右正合函子，\(\mathcal A\) 有足够投射对象，\(X\) 为 \(\mathcal A\) 中有上界的上链复形，并且 \(L_jF(X^n)=0\) 对所有 \(n\in\mathbb Z\)、\(j>0\) 成立，则对偶的自然映射 \(LF(X)\to F(X)\) 是拟同构。
\end{proposition}
\begin{proof}
先设 \(E\) 是有下界的正合复形，每项为 \(F\)-零调对象。记 \(Z^n=\ker \mathrm{d}_E^n\)。从低于所有非零项的次数开始，\(Z^n=0\)。短正合列
\[
0\longrightarrow Z^n\longrightarrow E^n\longrightarrow Z^{n+1}\longrightarrow0
\]
给出导出长正合列。若 \(Z^n\) 已 \(F\)-零调，其中的片段
\[
F(E^n)\longrightarrow F(Z^{n+1})\longrightarrow R^1F(Z^n)=0
\]
使第一个映射满；而 \(E^n\) 的零调性给出
\[
R^iF(Z^{n+1})\cong R^{i+1}F(Z^n)=0\qquad(i\ge1),
\]
使下一层 \(Z^{n+1}\) 也 \(F\)-零调。左正合性再给出
\(0\to FZ^n\to FE^n\to FZ^{n+1}\to0\) 正合。由起始的 \(Z^n=0\) 逐次向上归纳，\(F(E)\) 便正合。

取 \(X\to I\) 的有下界内射替换，其锥正合、有下界，且每项是内射对象与某个 \(X^n\) 的有限直和，故逐项 \(F\)-零调。上段说明 \(F(\operatorname{Cone}(X\to I))\) 正合。加性保证此复形是 \(\operatorname{Cone}(FX\to FI)\)，锥判据得到拟同构。对偶情形从最高次数向下归纳。
\end{proof}

上一个证明说明了有下界假设的实际用途：它给余圈对象的归纳提供起点。逐项对 \(F\) 零调是对象性质，整个复形经 \(F\) 后仍正合是另一个性质；无界时不能仅凭前者推出后者。

\subsection{复合与自然变换}
\label{b4:subsec:1801:03}

若 \(\alpha:F\to G\) 是两个加性函子间的自然变换，在同一个内射替换上逐项取 \(\alpha_{I_X}\)，便得到
\(R\alpha:RF\to RG\)。比较态射使它自然，且
\(R(\beta\alpha)=(R\beta)(R\alpha)\)、\(R(1_F)=1_{RF}\)。左导出使用同一个投射替换，也有相同结论。因而自然变换不是在每个导出群上重新猜测，而是先在复形层面构造。

\begin{theorem}\label{b4:thm:total-derived-composition}
设 \(\mathcal A,\mathcal B,\mathcal C\) 是局部小交换范畴，\(\mathcal A,\mathcal B\) 有足够内射对象，\(F:\mathcal A\to\mathcal B\)、\(G:\mathcal B\to\mathcal C\) 为加性左正合函子。总有自然比较
\[
R(GF)\longrightarrow RG\circ RF
\quad\text{在}\;D^+(\mathcal A)\;\text{上}.
\]
若 \(F\) 把内射对象送到 \(G\)-零调对象，则这个比较是自然同构。
\end{theorem}
\begin{proof}
取 \(X\to I\)，再取 \(F(I)\to J\) 的有下界内射替换。\(R(GF)(X)\) 由 \(GF(I)\) 表示，\(RG(RF(X))\) 由 \(G(J)\) 表示，比较就是
\(G(F(I))\to G(J)\)：第一次导出使用 \(I\)，第二次导出还要把 \(F(I)\) 换成内射模型 \(J\)。同伦比较保证选择独立与自然性；等价地，它由全右导出的泛性质作用于两次比较变换的复合产生。虽然 \(F(I)\to J\) 为拟同构，加性左正合的 \(G\) 未必保持它，所以这个比较并不自动可逆。

若每个 \(F(I^n)\) 都 \(G\)-零调，则\cref{b4:prop:total-derived-acyclic-complex}应用于有下界复形 \(F(I)\)，说明上述比较是拟同构。这也给出两种复合方式之间的自然同构。
\end{proof}

若 \(F\) 把内射对象送到 \(G\)-零调对象，则对一般
\(X\in D^+(\mathcal A)\)，将\cref{b4:thm:hypercohomology-spectral-sequences}的第二套谱序列应用于 \(RF(X)\) 的有下界复形模型，得到自然的强收敛谱序列
\[
E_2^{p,q}=R^pG\bigl(H^q(RF(X))\bigr)
\Longrightarrow H^{p+q}\bigl(R(GF)(X)\bigr).
\]
这里 \(p\ge0\)，\(q\) 有下界，所以每个总次数的滤过仍有限。若特别取
\(X=A[0]\)，其中 \(A\in\mathcal A\)，则
\(H^q(RF(A[0]))=R^qF(A)\) 对 \(q\ge0\) 成立，负次数为零。于是恢复\cref{b4:thm:grothendieck-spectral-sequence}的第一象限谱序列
\[
E_2^{p,q}=R^pG(R^qF(A))
\Longrightarrow R^{p+q}(GF)(A).
\]
它提供目标的滤过及其逐层商，而不自动给出直和分解
\[
R^n(GF)(A)\cong\bigoplus_{p+q=n}R^pG(R^qF(A)).
\]
全导出复合的比较与零调性条件见 Weibel\cite[Composition Theorem 10.8.2, Corollary 10.8.3]{Weibel1994HomologicalAlgebra}。

\ProseParagraphBreak
对右正合函子及投射替换有对偶版本：若 \(F\) 把投射对象送到 \(G\)-零调对象，则 \(L(GF)\simeq LG\circ LF\)；一般比较的方向为 \(LG\circ LF\to L(GF)\)。如果不检查零调性，就没有理由让比较成为同构。

\begin{proposition}\label{b4:prop:additive-right-derived-zero}
在交换群范畴 \(\mathbf{Ab}\) 上取加性函子
\[
F(M)=M\otimes_{\mathbb Z}\mathbb Z/2\mathbb Z.
\]
其全右导出函子 \(RF:D^+(\mathbf{Ab})\to D^+(\mathbf{Ab})\) 为零函子，但 \(F(\mathbb Z)=\mathbb Z/2\mathbb Z\ne0\)。因此在零次复形 \(\mathbb Z[0]\) 上，比较态射 \(F(\mathbb Z)[0]\to RF(\mathbb Z[0])\) 不是同构。仅有加性不能保证对每个交换群 \(M\) 都有 \(H^0(RF(M[0]))\cong F(M)\)。
\end{proposition}
\begin{proof}
设 \(I\) 是内射交换群。对任意 \(x\in I\)，将同态 \(2\mathbb Z\to I\)、\(2\mapsto x\) 沿包含 \(2\mathbb Z\hookrightarrow\mathbb Z\) 延拓为 \(f:\mathbb Z\to I\)。令 \(y=f(1)\)，便有 \(2y=f(2)=x\)。故乘 \(2\) 在 \(I\) 上满射，\(I/2I=0\)，从而
\[
F(I)\cong I/2I=0.
\]
对任意有下界复形 \(X\)，取有下界的逐项内射替换 \(X\to I_X\)。\(F(I_X)\) 逐项为零，所以\cref{b4:thm:total-derived-functors}给出的 \(RF\) 是零函子。另一方面，\(F(\mathbb Z)=\mathbb Z/2\mathbb Z\) 非零，其零次复形的 \(H^0\) 非零，不能同构于零导出对象 \(RF(\mathbb Z[0])\)。该函子为右正合函子，但它把单射 \(\mathbb Z\xrightarrow{2}\mathbb Z\) 送为 \(\mathbb Z/2\mathbb Z\) 上的零映射，故不是左正合函子。因此它的全右导出不受前面左正合情形的零次识别约束。
\end{proof}

% !TeX root = ../../Book4.tex
% 纲要标题：### 19.2 导出张量积
% 纲要标签：b4:sec:1802
% 纲要位置：计划纲要.md，第 3719 行。

\section{导出张量积}
\label{b4:sec:1802}

普通张量积可能在更换拟同构模型以后改变同调。平坦或投射替换修正这一问题，并使张量积的结合律、单位同构及 Tor 计算相互一致。有上界平坦复形的论证使用\cref{b4:prop:filtered-colimits-cohomology}中滤过余极限与同调交换的性质。

% 纲要内容（仅作写作计划，尚非教材正文）：
%
% * 平坦或投射替换
% * 结合律与适当的单位同构
% * Tor 作为导出张量的同调
%

\subsection{平坦替换与投射替换}
\label{b4:subsec:1802:01}

普通张量积不能保持所有拟同构。例如整数模复形 \(\mathbb Z\xrightarrow{2}\mathbb Z\)（次数 \(-1,0\)）拟同构于 \(\mathbb Z/2\mathbb Z\)，但与 \(\mathbb Z/2\mathbb Z\) 张量后前者多出一次负次数同调。为了让张量积在导出范畴中有意义，需要能保存拟同构的替换。

\begin{lemma}\label{b4:lem:bounded-flat-tensor-preserves}
设 \(R\) 为环，\(P\) 为逐项平坦且有上界的右 \(R\)-模复形，\(A\) 为任意正合的左 \(R\)-模复形。采用直和总化时，\(P\otimes_RA\) 正合。因此 \(P\otimes_R-\) 保持任意复形之间的拟同构。交换左右侧别后也成立。
\end{lemma}
\begin{proof}
若 \(P\) 只有一项，平坦性说明张量保留 \(A\) 的各个核像短列，所以结果正合。若 \(P\) 有有限个非零项，按最高非零项逐次剥去，所得复形短列在分次模意义下分裂；张量后仍是短正合列，用同调长正合列归纳，得到正合性。

一般情形取粗截断 \(P_{\ge m}\)，即低于 \(m\) 的项换成零。它是 \(P\) 的子复形，且因有上界而只有有限个非零项。截口可能改变同调，故这些粗截断一般不拟同构于 \(P\)；这里使用的是随着 \(m\) 向负方向移动，\(P\) 为它们的逐项滤过余极限。张量与直和总化都与此余极限交换，而模范畴的滤过余极限正合并与同调交换，见\cref{b4:prop:filtered-colimits-cohomology}。因此
\[
 H^n(P\otimes_RA)
 \cong\varinjlim_{m\to-\infty}H^n(P_{\ge m}\otimes_RA)=0.
\]

若 \(u:A\to B\) 是拟同构，其锥正合。张量锥与张量映射的锥由下面的分量识别自然同构：
\[
 \begin{gathered}
 \theta:P\otimes_R\operatorname{Cone}(u)
 \longrightarrow\operatorname{Cone}(1_P\otimes u),\\
 x\otimes(b,a)\longmapsto
 \bigl(x\otimes b,(-1)^p x\otimes a\bigr),
 \end{gathered}
\]
其中 \(x\in P^p\)、\(b\in B^q\)、\(a\in A^{q+1}=A[1]^q\)。第二分量的符号区别了先移位右因子与先张量再移位。具体地，在 \(x\otimes(0,a)\) 上，两边的微分计算均给出
\[
 \begin{aligned}
 \theta\mathrm{d}\bigl(x\otimes(0,a)\bigr)
 &=\mathrm{d}\theta\bigl(x\otimes(0,a)\bigr)\\
 &=\Bigl((-1)^p x\otimes u(a),\;
 (-1)^{p+1}\cdot\mathrm{d}_Px\otimes a-x\otimes\mathrm{d}_Aa\Bigr).
 \end{aligned}
\]
在 \(x\otimes(b,0)\) 上，\(\theta\) 不改变分量，微分相容性直接由张量微分得到。\(\theta\) 的逆使用同一符号。于是张量后的锥正合，得到后一结论。
\end{proof}

\begin{theorem}\label{b4:thm:derived-tensor}
设 \(R\) 是环。对有上界的右 \(R\)-模复形 \(X\) 和有上界的左 \(R\)-模复形 \(Y\)，选择逐项投射的有上界替换 \(p:P\to X\)、\(q:Q\to Y\)。公式
\[
X\otimes_R^{\mathbf L}Y\coloneqq P\otimes_RQ
\]
定义双函子
\[
-\otimes_R^{\mathbf L}-:
D^-(\operatorname{Mod}\text{-}R)\times D^-(R\text{-}\operatorname{Mod})
\longrightarrow D^-(\mathbf{Ab}).
\]
它称为\term[daochu-zhangliangji]{导出张量积}[derived tensor product]。也可以只替换一个变量，使用 \(P\otimes_RY\) 或 \(X\otimes_RQ\)。若 \(P'\to X\) 或 \(Q'\to Y\) 是有上界的逐项平坦替换，则 \(P'\otimes_RY\) 或 \(X\otimes_RQ'\) 分别典范同构于这个双函子的值；同时使用两侧的平坦替换也得到同一导出对象。
\end{theorem}
\symidx{\(\operatorname{Mod}\text{-}R\)：幺性右模范畴}
\begin{proof}
上一引理分别应用于 \(P\) 和 \(Q\)，说明自然映射
\[
P\otimes_RY\xleftarrow{\,1_P\otimes q\,}P\otimes_RQ
\xrightarrow{\,p\otimes1_Q\,}X\otimes_RQ
\]
都是拟同构。左箭头使用 \(P\) 有上界且逐项投射，从而与 \(P\) 张量保持 \(q\) 的拟同构；右箭头使用 \(Q\) 的对应性质，从而与 \(Q\) 张量保持 \(p\) 的拟同构。因此固定一侧的适当模型，另一侧就可以保留原复形。投射模型中的态射由同伦类表示，由\cref{b4:prop:tensor-hom-preserve-homotopy}，张量积在两个变量都保持同伦，所以先在两个投射同伦范畴上得到双函子，再由两侧的投射模型等价下降到所写导出范畴。

若 \(p':P'\to X\) 是有上界逐项平坦替换，由\cref{b4:thm:bounded-projective-model}证明中的比较结论，存在唯一的同伦类 \(t:P\to P'\)，使 \(p't=p\) 于同伦范畴中。由于 \(p,p'\) 都是拟同构，\(t\) 也是拟同构。于是
\[
P\otimes_RQ
\xrightarrow{\,t\otimes1_Q\,}P'\otimes_RQ
\xrightarrow{\,1_{P'}\otimes q\,}P'\otimes_RY
\]
的两条箭头都是拟同构：第一条使用 \(Q\) 的平坦性及引理的左右对偶，第二条使用 \(P'\) 的平坦性。张量保持同伦，所以这个合成不依赖 \(t\) 的代表元。若更换 \(P\) 或 \(Q\)，投射模型的比较同伦类分别满足与 \(p\) 或 \(q\) 的相容等式；\(t\) 的唯一性使上述合成与模型间的比较交换。因此它给出 \(P'\otimes_RY\) 与已定义双函子值的典范同构。交换左右侧别得到第二变量的结论，再依次替换两侧即可同时使用 \(P',Q'\)。双函子在态射上的作用始终由投射模型定义。两变量有上界使每个总次数只涉及有限条对角项，且输出仍有上界。
\end{proof}
\symidx{\(X\otimes_R^{\mathbf L}Y\)：导出张量积}

对于非交换环，上式一般只给交换群复形；不能在结果上无说明添一个左 \(R\)-模结构。若 \(X\) 还带与右 \(R\)-作用交换的左 \(S\)-作用，则张量复形上的作用为 \(s\cdot(x\otimes y)=(sx)\otimes y\)。可以固定 \(X\)，仅用左 \(R\)-模投射复形 \(Q\) 替换 \(Y\)，使 \(X\otimes_RQ\) 按这个公式保留左 \(S\)-作用，且微分仍为 \(S\)-线性的。若仅将 \(X\) 看作右 \(R\)-模而任取替换 \(P\to X\)，\(P\) 未必带有与微分相容的左 \(S\)-作用，此时 \(P\otimes_RY\) 也不能按上述公式成为左 \(S\)-模复形。若要替换 \(X\)，就应说明替换如何保留所需双模结构。模的有上界导出张量积可对照 Weibel\cite[Definition 10.6.1, Lemma 10.6.2, Theorem 10.6.3]{Weibel1994HomologicalAlgebra}。

\subsection{结合律与单位同构}
\label{b4:subsec:1802:02}

\begin{proposition}\label{b4:prop:derived-tensor-associativity}
设 \(R\) 为交换环，\(X,Y,Z\in D^-(R\text{-}\operatorname{Mod})\)。导出张量积有自然同构
\[
(X\otimes_R^{\mathbf L}Y)\otimes_R^{\mathbf L}Z
\simeq X\otimes_R^{\mathbf L}(Y\otimes_R^{\mathbf L}Z),
\qquad
R\otimes_R^{\mathbf L}X\simeq X\simeq X\otimes_R^{\mathbf L}R.
\]
还有对称同构 \(X\otimes_R^{\mathbf L}Y\simeq Y\otimes_R^{\mathbf L}X\)。若 \(P,Q\) 分别是 \(X,Y\) 的有上界逐项平坦替换，它在齐次元上由 \(x\otimes y\mapsto(-1)^{pq}y\otimes x\) 表示，其中 \(x\in P^p\)、\(y\in Q^q\)。
\end{proposition}
\begin{proof}
取有上界逐项投射替换 \(P,Q,U\)。交换环上两个投射模的张量积仍投射：把二者写成自由模的直和项，其张量就是自由模张量的直和项。每条总次数对角线有限，故 \(P\otimes_RQ\) 及 \(Q\otimes_RU\) 均逐项投射、有上界，可以直接用来计算下一次导出张量。

两边都由三重张量总复形表示，齐次纯张量的微分为
\[
\mathrm{d}(p\otimes q\otimes u)=\mathrm{d}p\otimes q\otimes u
+(-1)^{|p|}p\otimes \mathrm{d}q\otimes u
+(-1)^{|p|+|q|}p\otimes q\otimes \mathrm{d}u.
\]
普通张量的结合映射保持此公式，给出复形同构；对四个因子，所有结合映射都只改变括号，所以它们的不同复合相同。\(R[0]\) 本身投射，通常的单位映射逐项给出单位同构。交换两个因子的公式中，微分越过次数 \(p\) 或 \(q\) 时产生的符号正好补偿 \((-1)^{pq}\) 的改变，直接代入便得到复形映射；两次交换的符号是 \((-1)^{2pq}=1\)，故为同构。这些复形映射与投射模型中的比较态射相容，故在导出范畴中仍自然。
\end{proof}

单位结论不需要交换性：对任意环 \(R\)，右模复形 \(X\) 满足 \(X\otimes_R^{\mathbf L}R\simeq X\)，左模复形 \(Y\) 满足 \(R\otimes_R^{\mathbf L}Y\simeq Y\)，此处 \(R\) 带通常的 \((R,R)\)-双模结构。非交换环上的多重张量也有结合映射，但每两个相邻因子必须具有匹配的左右作用，并须核对保留这些作用的替换；上面的交换环命题没有免除这些要求。

\subsection{Tor 与导出张量的同调}
\label{b4:subsec:1802:03}

\begin{corollary}\label{b4:cor:tor-derived-tensor}
设 \(R\) 为环，\(M\) 为右 \(R\)-模，\(N\) 为左 \(R\)-模，二者视作零次复形。则对每个 \(n\ge0\)，有自然同构
\[
H^{-n}(M\otimes_R^{\mathbf L}N)\simeq\operatorname{Tor}_n^R(M,N),
\]
而正次数同调为零。特别地，零次同调为 \(M\otimes_RN\)。
\end{corollary}
\begin{proof}
把 \(M\) 的投射链消解 \(P_\bullet\to M\) 重编号为 \(P^{-n}=P_n\)，并只替换第一变量。张量复形在 \(-n\) 次的同调正是\cref{b4:thm:tor-balanced}所计算的 \(\operatorname{Tor}_n\)。其复形没有正次数项，零次同调由张量的右正合性识别。
\end{proof}

例如，设 \(R\) 为交换环，\(a\in R\) 不是零因子。\(R/(a)\) 由 \(R\xrightarrow{a}R\) 的两项自由复形替换，因此对任意 \(R\)-模 \(N\)，
\[
R/(a)\otimes_R^{\mathbf L}N\simeq
\bigl(N\xrightarrow{a}N\bigr),\qquad
H^{-1}=\{\,n\in N\mid an=0\,\},\quad H^0=N/aN.
\]
若 \(a\) 同时零化 \(N\)，微分为零，故在导出范畴中这是 \(N\oplus N[1]\)。若 \(a\) 在 \(N\) 上仍为非零因子，则负次同调消失，此时才可用普通张量积代替。\(a\) 在 \(R\) 上不是零因子的条件保证所用两项自由复形确为消解。

\ProseParagraphBreak
把这个公式用于两条直线，便能看出负次同调记录了什么。设 \(k\) 为域，\(R=k[x,y]\)，并取 \(M=R/(x)\)。若 \(N=R/(y)=k[x]\)，则乘 \(x\) 单射，因而
\[
 H^{-1}(M\otimes_R^{\mathbf L}N)=0,\qquad
 H^0(M\otimes_R^{\mathbf L}N)=k.
\]
两条坐标轴的方程彼此独立，普通张量积已经给出它们的交点。若改取 \(N=M=k[y]\)，则乘 \(x\) 为零，故
\[
 M\otimes_R^{\mathbf L}M\simeq M\oplus M[1],\qquad
 H^{-1}=M,
\]
而普通张量积只有 \(H^0=M\)。这一次两个模使用同一条方程，新增的负次同调保留了方程重复所产生的关系。导出张量积既记录共同的商，也记录形成这个商时关系之间的依赖。

% !TeX root = ../../Book4.tex
% 纲要标题：### 19.3 导出 Hom
% 纲要标签：b4:sec:1803
% 纲要位置：计划纲要.md，第 2964 行。

\section{导出 Hom}
\label{b4:sec:1803}

导出 Hom 保留全部 Ext 次数，导出 Hom--张量伴随则联系两个不同的计算模型。基变换需要额外假设，双数代数的周期消解给出非平坦情形的明确反例。

% 纲要内容（仅作写作计划，尚非教材正文）：
%
% * RHom 与 Ext
% * 导出 Hom–张量伴随
% * 基变换的充分条件及失败例子
% * 以 \(R=k[\varepsilon]/(\varepsilon^2)\) 和 \(k=R/(\varepsilon)\) 对照普通张量与导出张量：周期自由消解给出非零的高阶 Tor，说明非平坦基变换丢失的信息
%

\subsection{RHom 与 Ext}
\label{b4:subsec:1803:01}

\cref{b4:def:hom-complex}已经定义
\[
\operatorname{Hom}_R^n(X,Y)=
\prod_p\operatorname{Hom}_R(X^p,Y^{p+n}),\qquad
\delta f=\mathrm{d}_Yf-(-1)^nf\mathrm{d}_X.
\]
零次余圈是复形映射，零次余边是零同伦映射。更一般地，次数 \(n\) 的余圈满足 \(\mathrm{d}_Yf=(-1)^nf\mathrm{d}_X\)，正好给出 \(X\to Y[n]\) 的复形映射，因此
\(H^n\operatorname{Hom}_R^\bullet(X,Y)=\operatorname{Hom}_{K(R)}(X,Y[n])\)。把同伦范畴换成导出范畴，也需要修正 Hom 复形。

\begin{theorem}\label{b4:thm:derived-hom}
设 \(R\) 为环，记 \(D(R)=D(R\text{-}\mathbf{Mod})\)。设 \(X\) 为有上界的左 \(R\)-模上链复形，\(Y\) 为有下界的左 \(R\)-模上链复形。取有上界投射替换 \(P\to X\) 与有下界内射替换 \(Y\to I\)。三个交换群复形
\[
\operatorname{Hom}_R^\bullet(P,Y),\qquad
\operatorname{Hom}_R^\bullet(P,I),\qquad
\operatorname{Hom}_R^\bullet(X,I)
\]
自然拟同构，定义
\[
\mathbf R\operatorname{Hom}_R:
D^-(R\text{-}\operatorname{Mod})^{\mathrm{op}}
\times D^+(R\text{-}\operatorname{Mod})
\longrightarrow D^+(\mathbf{Ab}).
\]
这称为\term[daochu-Hom]{导出 Hom}[derived Hom]。其同调满足
\[
H^n\mathbf R\operatorname{Hom}_R(X,Y)
\simeq\operatorname{Hom}_{D(R)}(X,Y[n])
\quad(n\in\mathbb Z).
\]
若 \(R\) 交换，所构造的 Hom 复形带自然的 \(R\)-模结构，输出可取在 \(D^+(R\text{-}\operatorname{Mod})\) 中。
\end{theorem}
\begin{proof}
\cref{b4:thm:bounded-projective-model}的比较结论说明，对于任意正合复形 \(A\)，所有 \(P\to A[n]\) 都零同伦，所以 \(\operatorname{Hom}_R^\bullet(P,A)\) 正合。\cref{b4:thm:bounded-injective-model}对偶地说明 \(\operatorname{Hom}_R^\bullet(A,I)\) 正合。因此对 \(Y\to I\) 和 \(P\to X\) 的正合锥施加 Hom，便得到所写的两条拟同构。Hom 与锥相容时使用其定义中的符号，或在每次同调上直接用上述同伦类解释，都会给出同一结论。

投射及内射比较态射在同伦意义下唯一，因此此构造在两个变量上自然，并把拟同构送到同构，从而定义所写双函子。若 \(P^p=0\) 对 \(p>a\)，\(I^j=0\) 对 \(j<b\)，则 \(\operatorname{Hom}^n(P,I)=0\) 对 \(n<b-a\)，所以输出有下界。最后
\[
H^n\operatorname{Hom}_R^\bullet(P,I)
=\operatorname{Hom}_{K(R)}(P,I[n])
\simeq\operatorname{Hom}_{D(R)}(X,Y[n]),
\]
最后一步仍是投射或内射比较结论。交换环时，标量与所有映射及比较相容，故得到 \(R\)-模版本。
\end{proof}
\symidx{\(\mathbf R\operatorname{Hom}_R(X,Y)\)：导出 Hom}

对模 \(M,N\) 置于零次，以上公式化为
\[
H^n\mathbf R\operatorname{Hom}_R(M,N)=\operatorname{Ext}_R^n(M,N)
\quad(n\ge0),
\]
负次为零，与\cref{b4:thm:derived-ext-identification}一致。把所有 Ext 群单独列出，通常不能恢复导出 Hom 复形；后者还保留它在三角和复合中的相容关系。只替换第二变量时可以允许第一变量无界，只替换第一变量时可以允许第二变量无界；本节的共同有界范围使两种计算可以直接比较。参阅 Weibel\cite[Definition 10.7.2, Theorem 10.7.4, Corollary 10.7.5]{Weibel1994HomologicalAlgebra}。

\subsection{导出 Hom–张量伴随}
\label{b4:subsec:1803:02}

普通的 Hom--张量伴随把双线性映射写成一族线性映射。复形版本仍用同一个代入公式，不过必须核对各个次数和微分。

\begin{theorem}\label{b4:thm:derived-hom-tensor-adjunction}
设 \(R\) 为交换环，记 \(D(R)=D(R\text{-}\mathbf{Mod})\)。设 \(X,Y\in D^-(R\text{-}\operatorname{Mod})\)，\(Z\in D^+(R\text{-}\operatorname{Mod})\)。存在自然同构
\[
\mathbf R\operatorname{Hom}_R(X\otimes_R^{\mathbf L}Y,Z)
\simeq
\mathbf R\operatorname{Hom}_R
\bigl(X,\mathbf R\operatorname{Hom}_R(Y,Z)\bigr).
\]
特别地，在零次同调上得到
\[
\operatorname{Hom}_{D(R)}(X\otimes_R^{\mathbf L}Y,Z)
\simeq
\operatorname{Hom}_{D(R)}
\bigl(X,\mathbf R\operatorname{Hom}_R(Y,Z)\bigr).
\]
\end{theorem}
\begin{proof}
取有上界投射替换 \(P\to X\)、\(Q\to Y\)，以及有下界内射替换 \(Z\to I\)。\(P\otimes_RQ\) 是有上界逐项投射复形：交换环上两个投射模的张量是投射模，因为它是两个自由模张量的直和项。故左侧由 \(\operatorname{Hom}^\bullet_R(P\otimes_RQ,I)\) 表示。

内层 \(\mathbf R\operatorname{Hom}_R(Y,Z)\) 由 \(\operatorname{Hom}^\bullet_R(Q,I)\) 表示，而且有下界。外层可只替换第一变量，故右侧由
\(\operatorname{Hom}^\bullet_R(P,\operatorname{Hom}^\bullet_R(Q,I))\) 表示。逐次定义
\[
\Phi(f)(p)(q)=f(p\otimes q).
\]
若 \(f\) 的次数为 \(n\)，\(p,q\) 的次数为 \(a,b\)，则值在 \(I^{a+b+n}\)，与两边的次数规定一致。展开右侧的微分得到
\[
\mathrm{d}_I f(p\otimes q)
-(-1)^{a+n}f(p\otimes \mathrm{d}_Qq)
-(-1)^n f(\mathrm{d}_Pp\otimes q),
\]
它正是左侧微分后再施加 \(\Phi\) 的结果。普通张量泛性质逐项给出逆映射；Hom 中的积允许对所有次数同时指定映射，所以这个对应没有有限生成假设。它是自然的复形同构。取零次同调，并用\cref{b4:thm:derived-hom}，得到第二个公式。
\end{proof}

\begin{proposition}\label{b4:prop:derived-extension-restriction}
设 \(f:R\to S\) 为含幺环同态，不要求交换。简记左 \(R\)-模、左 \(S\)-模的导出范畴为 \(D(R),D(S)\)。对有上界的左 \(R\)-模复形 \(X\) 与有下界的左 \(S\)-模复形 \(Y\)，有自然同构
\[
\operatorname{Hom}_{D(S)}(S\otimes_R^{\mathbf L}X,Y)
\simeq\operatorname{Hom}_{D(R)}(X,\operatorname{Res}_R^S Y).
\]
其中导出标量扩张由 \(S\otimes_RP\) 计算，\(P\to X\) 为有上界投射替换，且 \(S\) 带左 \(S\)、右 \(R\) 的通常双模结构。
\end{proposition}
\begin{proof}
\(S\otimes_RP\) 逐项是投射左 \(S\)-模，因为标量扩张把自由模送到自由模并保持直和项。普通伴随逐项给出复形同构
\[
\operatorname{Hom}_S^\bullet(S\otimes_RP,Y)
\simeq\operatorname{Hom}_R^\bullet(P,\operatorname{Res}_R^S Y),
\]
映射为 \(g\mapsto(p\mapsto g(1\otimes p))\)，逆为 \(h\mapsto(s\otimes p\mapsto s\cdot h(p))\)。二者保持 Hom 微分。投射比较把两边的零次同调分别识别为所需导出态射群。
\end{proof}

这一伴随不要求 \(S\) 在 \(R\) 上平坦，因为已经替换 \(X\)。忘却函子本来正合，不另取导出；它与导出标量扩张的关系可对照 Weibel\cite[Theorem 10.7.6]{Weibel1994HomologicalAlgebra}，该书在此使用交换环版本。

\subsection{基变换的充分条件与反例}
\label{b4:subsec:1803:03}

基变换有两种不同的问题。一个是把 \(R\) 换成 \(S\) 后，普通张量是否已能计算导出张量；另一个是张量能否穿过 Hom。前者由平坦性控制，后者还涉及有限生成条件，不能只核对其中一项。

\begin{proposition}\label{b4:prop:flat-base-change-hom}
设 \(R\to S\) 是交换环同态，\(S\) 作为 \(R\)-模平坦。设 \(X\in D^-(R\text{-}\mathbf{Mod})\) 在该导出范畴中同构于有限生成投射 \(R\)-模的有上界复形 \(P\)，\(Y\) 为有下界的 \(R\)-模复形。则自然映射
\[
S\otimes_R\mathbf R\operatorname{Hom}_R(X,Y)
\longrightarrow
\mathbf R\operatorname{Hom}_S(S\otimes_RX,S\otimes_RY)
\]
是同构。左侧及右侧变量中的普通张量均已代表导出标量扩张。
\end{proposition}
\begin{proof}
平坦性使 \(S\otimes_R-\) 保持拟同构，故可用 \(P\) 替换 \(X\)。在这些模型上定义复形映射
\[
\alpha:S\otimes_R\operatorname{Hom}_R^\bullet(P,Y)
\longrightarrow
\operatorname{Hom}_S^\bullet(S\otimes_RP,S\otimes_RY),
\]
其中对次数 \(n\) 的齐次映射 \(f\)、\(x\in P^i\) 及 \(s,s'\in S\)，规定
\[
\alpha(s\otimes f)(s'\otimes x)=ss'\otimes f(x).
\]
\(R,S\) 的交换性保证该公式对张量关系良定义，且所得映射为 \(S\)-线性；将 Hom 微分代入也得到 \(\alpha\delta=\delta\alpha\)。这个公式与两个变量的复形映射相容，因而传到导出范畴后就是命题中的自然映射。

对有限生成投射模 \(P^i\)，它在每个位置给出的映射
\[
S\otimes_R\operatorname{Hom}_R(P^i,Y^j)
\longrightarrow
\operatorname{Hom}_S(S\otimes_RP^i,S\otimes_RY^j)
\]
为同构：有限自由模时是有限个副本的恒等识别，直和项给出投射情形。固定总次数 \(n=j-i\)，有上界的 \(P\) 与有下界的 \(Y\) 只留下有限多个非零位置，所以张量可以与该次数中的 Hom 积交换。以上逐项同构与微分相容，给出复形同构。\(S\otimes_RP\) 逐项投射，故其 Hom 计算右侧。
\end{proof}

\begin{proposition}\label{b4:prop:derived-base-change-perfect}
设 \(R\to S\) 是任意交换环同态，\(X\in D(R\text{-}\mathbf{Mod})\) 在该导出范畴中同构于有限生成投射 \(R\)-模的有界复形，\(Y\in D^-(R\text{-}\mathbf{Mod})\)。则存在自然同构
\[
S\otimes_R^{\mathbf L}\mathbf R\operatorname{Hom}_R(X,Y)
\simeq
\mathbf R\operatorname{Hom}_S
\bigl(S\otimes_R^{\mathbf L}X,S\otimes_R^{\mathbf L}Y\bigr).
\]
\end{proposition}
\begin{proof}
用给定的有界有限生成投射复形 \(P\) 表示 \(X\)，取有上界投射替换 \(Q\to Y\)。由于 \(P\) 有界，\(\operatorname{Hom}_R^\bullet(P,Q)\) 有上界，每项是有限个投射模的直和：\(\operatorname{Hom}_R(P^i,Q^j)\) 是有限个 \(Q^j\) 的直和项。因此它既表示 \(\mathbf R\operatorname{Hom}_R(X,Y)\)，又可直接计算左侧的导出张量。

\(S\otimes_RP\) 仍为有界投射复形，故右侧可由
\(\operatorname{Hom}_S^\bullet(S\otimes_RP,S\otimes_RQ)\) 计算。上一命题中定义的比较映射对有限生成投射模逐项给出 Hom--基变换同构，这一步不需要平坦性；此处每次只有有限个项，把它们合起来便得到自然的复形同构。第二变量无需有下界：\(P\) 与 \(S\otimes_RP\) 均为有界投射复形，\cref{b4:prop:resolution-computes-derived-hom}对任意目标复形的比较结论保证上述 Hom 模型有效，而 \(Q\) 有上界保证左侧的导出张量仍在有上界范围内计算。
\end{proof}

平坦性不能在第一个命题中直接删去。取 \(R=\mathbb Z\)、\(S=\mathbb Z/2\)、\(M=\mathbb Z/2\)、\(N=\mathbb Z\)。则
\[
S\otimes_{\mathbb Z}\operatorname{Hom}_{\mathbb Z}(M,N)=0,
\qquad
\operatorname{Hom}_S(S\otimes M,S\otimes N)=S.
\]
在一次 Ext 上，左侧 \(S\otimes\operatorname{Ext}_{\mathbb Z}^1(M,N)=S\)，右侧 \(\operatorname{Ext}_S^1(S,S)=0\)。这些计算已显示普通基变换不能直接移入 Hom 或 Ext。对前一命题中的全复形比较，用自由消解构造的 Hom 模型，\(\mathbf R\operatorname{Hom}_{\mathbb Z}(S,\mathbb Z)\) 由零、一次的 \(\mathbb Z\xrightarrow{-2}\mathbb Z\) 表示；对这个模型作普通张量，得到两次均为 \(S\)、微分为零的复形，而右侧 \(\mathbf R\operatorname{Hom}_S(S,S)=S[0]\)。失配出现在一次上同调。由于普通张量不保持拟同构，离开平坦假设后，左侧若不指定模型，甚至不是导出范畴上的良定义运算；改成导出张量才能消除这一问题。

\ProseParagraphBreak
有限生成条件也有实质作用。即使 \(\mathbb Q\) 在 \(\mathbb Z\) 上平坦，对 \(P=\bigoplus_{j\ge1}\mathbb Z\)、\(N=\mathbb Z\)，比较映射为
\[
\mathbb Q\otimes_{\mathbb Z}\prod_{j\ge1}\mathbb Z
\longrightarrow\prod_{j\ge1}\mathbb Q.
\]
左侧像中的序列有一个共同分母，右侧的 \((1/j!)_{j\ge1}\) 没有共同分母，所以不满。无限自由虽然投射，仍不能把张量从无限 Hom 积外直接搬入。

\subsection{双数代数上的周期消解与非平坦基变换}
\label{b4:subsec:1803:04}

\begin{example}\label{b4:exm:dual-numbers-derived-base-change}
设 \(k\) 是域，\(R=k[\varepsilon]/(\varepsilon^2)\)，\(k=R/(\varepsilon)\)。取每项为 \(R\) 的有上界自由复形
\[
P=\bigl(\cdots\xrightarrow{\varepsilon}R
\xrightarrow{\varepsilon}R\xrightarrow{\varepsilon}R\bigr),
\qquad P^i=R\ (i\le0),
\]
并用零次的商映射增广到 \(k\)。在 \(R\) 中，乘 \(\varepsilon\) 的核与像都是 \((\varepsilon)\)，故增广正合，\(P\to k\) 是自由替换。与 \(k\) 张量后微分全为零，因此
\[
H^{-n}(k\otimes_R^{\mathbf L}k)=k\quad(n\ge0),\qquad
k\otimes_Rk=k.
\]
普通张量只保留零次项，所有负次同调均被遗漏。这同时说明 \(k\) 在 \(R\) 上不平坦，且投射维数无限；这里是在导出对象上重新观察\cref{b4:exm:infinite-dimension-dual-numbers}的同一个周期消解。

\ProseParagraphBreak
令 \(f:R\to k\)。限制标量本来正合，而再作导出标量扩张会得到 \(k\otimes_R^{\mathbf L}k\)，不是恒等对象 \(k\)。因此“先忘却，再扩张就回到原对象”即使对商环上的自由秩一模也不成立。伴随的余单位确实给出
\(k\otimes_R^{\mathbf L}k\to k\)，在上述代表中它保留零次项、把负次数送到零；它的存在不表示它是同构。

\ProseParagraphBreak
同一个 \(P\) 计算导出 Hom 时，\(\operatorname{Hom}_R(P,k)\) 的微分也全为零，于是
\[
\operatorname{Ext}_R^n(k,k)=k\quad(n\ge0).
\]
张量的额外同调向负次数延伸，Hom 的额外同调向正次数延伸，二者都在记录未能用有限长自由消解代替 \(k\) 的事实。把 \(R\) 直接换成 \(k\) 后，\(k\) 成为自由模，高次 \(\operatorname{Tor}^k\) 和 \(\operatorname{Ext}_k\) 全部消失；这正是普通非平坦基变换丢失的信息。
\end{example}

% !TeX root = ../../Book4.tex
% 纲要标题：### 19.4* 无界情形
% 纲要标签：b4:sec:1804
% 纲要位置：计划纲要.md，第 2971 行。

\OptionalSection{无界情形}{b4:sec:1804}

无界复形没有可供逐度归纳的起点，逐项投射或内射便不再自动足够。K-投射、K-内射与 K-平坦直接要求复形正确处理零调复形，并支持无界替换。\subsecref{b4:subsec:2002:03}中任意余积的构造和紧对象的反向刻画则以这里的半自由替换及其 K-投射性质为前提。

% 纲要内容（仅作写作计划，尚非教材正文）：
%
% * K-内射、K-投射与 K-平坦
% * 无界消解的必要修正
% * 不把有界证明直接无条件推广
%
% ---
%

\subsection{K-内射、K-投射与 K-平坦}
\label{b4:subsec:1804:01}

有上界投射复形和有下界内射复形之所以有效，关键不是每一项的名称，而是整个复形对正合复形的作用。无界时应直接把这种作用列为条件。

\begin{definition}\label{b4:def:k-resolutions}
设 \(R\) 为环，以下均为上链复形。
\begin{enumerate}
\item 左 \(R\)-模复形 \(P\) 若对每个正合左 \(R\)-模复形 \(A\)，都有 \(\operatorname{Hom}_R^\bullet(P,A)\) 正合，则称 \(P\) 为\term[K-toushe]{K-投射复形}[K-projective complex]。
\item 左 \(R\)-模复形 \(I\) 若对每个正合左 \(R\)-模复形 \(A\)，都有 \(\operatorname{Hom}_R^\bullet(A,I)\) 正合，则称 \(I\) 为\term[K-neishe]{K-内射复形}[K-injective complex]。
\item 左 \(R\)-模复形 \(F\) 若对每个正合右 \(R\)-模复形 \(A\)，直和总复形 \(A\otimes_RF\) 都正合，则称 \(F\) 为\term[K-pingtan]{K-平坦复形}[K-flat complex]。右模版本交换两个侧别。
\end{enumerate}
\end{definition}

前两项分别等价于所有 \(P\to A[n]\)、\(A\to I[n]\) 都零同伦，因为 Hom 复形的同调正是这些同伦类。由映射锥判据，这也等价于相应 Hom 函子把拟同构送到拟同构；K-平坦性则等价于张量函子具有这个性质。

\ProseParagraphBreak
\cref{b4:thm:bounded-projective-model}的提升论证说明，有上界的逐项投射复形是 K-投射的；其对偶说明有下界的逐项内射复形是 K-内射的。\cref{b4:lem:bounded-flat-tensor-preserves}说明有上界的逐项平坦复形是 K-平坦的。最后一个结论允许另一个变量无界，原因在于直和总化与滤过余极限的正合性，而不是对角线自动有限。

\begin{example}\label{b4:exm:unbounded-free-not-k-projective}
仍取 \(R=k[\varepsilon]/(\varepsilon^2)\)，令 \(E^n=R\) 对每个整数 \(n\)，所有微分均为乘 \(\varepsilon\)。\(E\) 正合且每项自由，但不是 K-投射：若 \(1_E\) 零同伦，则每度存在乘某元素 \(a_n\) 的同伦，使
\(1=\varepsilon a_n+a_{n+1}\varepsilon\)，右侧属于 \((\varepsilon)\)，矛盾。因此 \(\operatorname{Hom}_R(E,E)\) 的零次同调非零。

\ProseParagraphBreak
它也不是 K-平坦。取\cref{b4:exm:dual-numbers-derived-base-change}中的自由替换 \(P\to k\)。若 \(E\) 是 K-平坦，则
\(E\otimes_RP\to E\otimes_Rk\) 为拟同构。\(P\) 本身有上界且平坦，所以 \(P\otimes_RE\) 正合；交换因子的分次符号给出 \(E\otimes_RP\) 正合。但 \(E\otimes_Rk\) 每度为 \(k\)、微分为零，不正合，矛盾。

\ProseParagraphBreak
这个例子的每项甚至内射。取线性泛函 \(\lambda:R\to k\)，\(\lambda(a+b\varepsilon)=b\)。映射
\(R\to\operatorname{Hom}_k(R,k)\)、\(r\mapsto(s\mapsto\lambda(sr))\)
是 \(R\)-模同构：双线性型在基 \(1,\varepsilon\) 下的矩阵为 \(\left(\begin{smallmatrix}0&1\\1&0\end{smallmatrix}\right)\)。又有
\(\operatorname{Hom}_R(M,\operatorname{Hom}_k(R,k))\simeq\operatorname{Hom}_k(M,k)\)，右侧在向量空间短正合列上正合，故 \(R\) 内射。尽管如此，\(E\) 不是 K-内射，因为 \(E\) 正合而 \(\operatorname{Hom}_R(E,E)\) 不正合。
\end{example}

\subsection{无界消解的修正}
\label{b4:subsec:1804:02}

在模范畴中，任意无界复形仍可得到足够好的替换，但须按复形的全部循环关系构造。下面在模范畴中构造这样的替换。

\begin{definition}\label{b4:def:semifree-complex}
设 \(R\) 为环，\(P\) 是左 \(R\)-模复形。若存在子复形滤过
\[
0=P_{-1}\subseteq P_0\subseteq P_1\subseteq\cdots,
\qquad P=\bigcup_{s\ge0}P_s,
\]
使每个包含在分次模意义下分裂，每个商 \(P_s/P_{s-1}\) 是若干移位 \(R[n]\) 的直和且微分为零，则称此 \(P\) 连同滤过为\term[banziyou-fuxing]{半自由复形}[semi-free complex]。等价地，每一层可添一批齐次自由生成元，新生成元的微分只涉及更早层的生成元。
\end{definition}

\begin{theorem}\label{b4:thm:semifree-resolution}
设 \(R\) 为环，\(X\) 为任意左 \(R\)-模上链复形。存在半自由复形 \(P\) 及拟同构 \(P\to X\)。每个半自由复形都是 K-投射且 K-平坦的；这里不要求 \(X\) 或 \(P\) 有任何次数界。
\end{theorem}
\begin{proof}
先对 \(X\) 每个次数的每个余圈 \(z\in Z^nX\)，添一个次数 \(n\) 的自由生成元 \(e_z\)，令其微分为零，映到 \(z\)。这些余圈由集合 \(\coprod_{n\in\mathbb Z}Z^nX\) 索引，所得 \(P_0\to X\) 在同调上满射。

假定已构造 \(f_s:P_s\to X\)。对每对数据
\[
z\in Z^nP_s,\qquad x\in X^{n-1},\qquad f_s(z)=\mathrm{d}_Xx,
\]
添一个次数 \(n-1\) 的自由生成元 \(e_{z,x}\)；本层全部数据是集合 \(\coprod_{n\in\mathbb Z}(Z^nP_s\times X^{n-1})\) 的子集。规定
\(\mathrm{d} e_{z,x}=z\)、\(f_{s+1}(e_{z,x})=x\)。因为 \(\mathrm{d}z=0\) 且 \(f_s(z)=\mathrm{d}x\)，复形条件与映射条件同时成立。这定义 \(P_{s+1}\)，其分次模是 \(P_s\) 与所添自由模的直和，商微分为零。取并得到半自由复形 \(P\) 及映射 \(f:P\to X\)。

每个 \(X\) 的同调类在 \(P_0\) 已有余圈代表，因而 \(H(f)\) 满。若 \(z\in Z^nP\) 映成 \(X\) 中的边界，则 \(z\) 只含有限个生成元，属于某个 \(P_s\)；选定 \(x\) 使 \(f(z)=\mathrm{d}x\)，下一层就添入 \(e_{z,x}\) 使 \(z\) 成为边界。因此 \(H(f)\) 单，\(f\) 是拟同构。

为证 K-投射，设 \(A\) 正合，\(g:P\to A\) 是复形映射。逐层构造次数 \(-1\) 的映射 \(h\)，使 \(g=\mathrm{d}h+h\mathrm{d}\)。若 \(h\) 已在旧层定义，对次数 \(n\) 的新生成元 \(e\)，\(\mathrm{d}e\) 属于旧层，而且
\[
\mathrm{d}_A\bigl(g(e)-h(\mathrm{d}e)\bigr)=g(\mathrm{d}e)-\mathrm{d}_Ah(\mathrm{d}e)=0,
\]
后一等式用旧层上的同伦公式及 \(\mathrm{d}^2e=0\)。由 \(A\) 正合，每个 \(g(e)-h(\mathrm{d}e)\) 都有原像。按\subsecref{b4:subsec:0101:03}的集合指标选择约定，可为本层全部新生成元同时选取 \(h(e)\in A^{n-1}\)，使其微分为 \(g(e)-h(\mathrm{d}e)\)。自由性允许线性延拓到本层。逐层取并即得全局同伦。把 \(A\) 换成任意移位，得到所有 Hom 同调消失。

为证 K-平坦，设 \(A\) 是正合右模复形。每个商
\(A\otimes_R(P_s/P_{s-1})\) 是 \(A\) 的移位直和，所以正合。滤过包含逐次分裂，张量后保持短正合列，归纳得 \(A\otimes_RP_s\) 正合。最后由滤过余极限的正合性，\(A\otimes_RP\) 正合。
\end{proof}

这里的归纳不是沿复形次数开始，而是沿添加生成元的层次开始。即使次数向两边无限延伸，每个新生成元的微分只依赖已经处理的部分，所以构造同伦仍有起点。

\subsection{有界论证向无界情形推广的条件}
\label{b4:subsec:1804:03}

\begin{proposition}\label{b4:prop:k-resolution-comparison}
设 \(R\) 为环，简记左 \(R\)-模的同伦范畴和导出范畴为 \(K(R)=K(R\text{-}\mathbf{Mod})\)、\(D(R)=D(R\text{-}\mathbf{Mod})\)。若 \(P\) 为 K-投射左 \(R\)-模上链复形，则对任意左 \(R\)-模上链复形 \(Y\)，有
\[
\operatorname{Hom}_{K(R)}(P,Y)\simeq\operatorname{Hom}_{D(R)}(P,Y).
\]
若 \(I\) 为 K-内射左 \(R\)-模上链复形，则对任意左 \(R\)-模上链复形 \(X\)，有对偶同构
\(\operatorname{Hom}_{K(R)}(X,I)\simeq\operatorname{Hom}_{D(R)}(X,I)\)。
\end{proposition}
\begin{proof}
设 \(s:Y'\to Y\) 为拟同构。\(\operatorname{Cone}(s)\) 正合，K-投射性及 Hom 的锥长正合列说明
\(\operatorname{Hom}_{K(R)}(P,Y')\to\operatorname{Hom}_{K(R)}(P,Y)\)
是同构。因此从 \(P\) 出发的同伦态射函子已经把所有分母变成同构。

更具体地，导出态射的屋顶 \(P\xleftarrow{s}U\xrightarrow{a}Y\) 中，\(s\) 对 \(\operatorname{Hom}_{K(R)}(P,-)\) 给出同构，故唯一有 \(t:P\to U\) 使 \(st=1_P\) 在同伦范畴中成立。由拟同构的二取三性质，\(t\) 也是拟同构。因此原屋顶与 \(P\xleftarrow{1_P}P\xrightarrow{at}Y\) 以 \(P\) 为共同细化对象，细化映射分别为 \(t:P\to U\) 与 \(1_P:P\to P\)：左侧由 \(st=1_P\) 交换，右侧由 \(a\circ t=at\) 交换。故原屋顶由普通态射 \(at\) 表示。若两个这样的态射在导出范畴中相等，分式等价关系给出一个共同分母；对该分母同样施加上述 Hom 同构，便得它们已在同伦范畴中相等。对偶证明使用反向屋顶及 \(\operatorname{Hom}_{K(R)}(-,I)\)。
\end{proof}

右 \(R\)-模等同于左 \(R^{\mathrm{op}}\)-模，所以\cref{b4:def:semifree-complex,b4:thm:semifree-resolution}应用于反环 \(R^{\mathrm{op}}\) 时，给出右半自由复形及其替换；滤过商是移位右正则模 \(R_R[n]\) 的直和。给定任意无界右 \(R\)-模复形 \(X\) 与左 \(R\)-模复形 \(Y\)，分别取右半自由替换 \(P\to X\) 与左半自由替换 \(Q\to Y\)。直和总复形 \(P\otimes_RQ\) 表示 \(X\otimes_R^{\mathbf L}Y\)。K-平坦性使 \(P\otimes_RY\) 或 \(X\otimes_RQ\) 也能计算同一导出对象，K-投射比较保证态射和不同替换之间的相容性。

\ProseParagraphBreak
对任意无界左 \(R\)-模复形 \(X,Y\)，导出 Hom 可在第一变量用左半自由替换 \(P\to X\)，定义 \(\mathbf R\operatorname{Hom}_R(X,Y)=\operatorname{Hom}_R^\bullet(P,Y)\)，其中 Hom 仍取各次数的直积。无界替换的作用来自 K-投射、K-平坦性质，直和总化的对角线上可以有无限多个非零项。

\ProseParagraphBreak
对一般交换范畴，未必有自由生成元，也未必有所需的无限直和或精确余极限。若要用 K-内射复形构造无界 \(RF\)，必须另证每个对象有相应替换，并检查 \(F\) 在这些替换间保持拟同构；仅有“足够内射对象”并不自动提供这里使用的整个复形性质。可比较 Weibel\cite[Generalized Existence Theorem 10.5.9]{Weibel1994HomologicalAlgebra}所列出的替换与零调性条件；有界性有限制时，前面各节已经给出了可直接使用的构造。


\begin{chaptersummary}
全导出运算先选择能控制拟同构的复形，再进行原有加性运算。内射模型给出全右导出，投射模型给出全左导出；各项对相应函子零调的单侧有界复形也能计算相应结果。复合导出之间总有比较，只有满足相应零调性条件时才得到同构。

\ProseParagraphBreak
导出张量的负次同调是 Tor，导出 Hom 的正次同调在模上是 Ext。它们之间的伴随来自带分次符号的复形同构。平坦基变换和有限投射复形分别提供两类可验证的基变换条件；双数代数说明非平坦的普通张量可能同时丢掉无限多项。无界情形的修正，则是用 K-性质控制整个复形，而不只检查各项。
\end{chaptersummary}

\BookExercises

\begin{exercise}\label{b4:ex:18-additive-right-derived}
取加性函子 \(F:\mathbf{Ab}\to\mathbf{Ab}\)，\(F(M)=\operatorname{Hom}_{\mathbb Z}(\mathbb Z/2\mathbb Z,M)\)。证明 \(LF:D^-(\mathbf{Ab})\to D^-(\mathbf{Ab})\) 是零函子，并计算 \(RF(\mathbb Z)\)。分别判定比较态射 \(F(\mathbb Z)[0]\to RF(\mathbb Z)\) 与 \(LF(\mathbb Z/2\mathbb Z)\to F(\mathbb Z/2\mathbb Z)[0]\) 是否为同构，说明两种导出构造与 \(F\) 的正合方向有何关系。
\end{exercise}
\begin{solution}
投射交换群是自由交换群的直和因子，因而无挠；从 \(\mathbb Z/2\mathbb Z\) 到它的同态只能为零。所以任意有上界投射替换逐项施加 \(F\) 后都为零，\(LF=0\)。另一方面，用\cref{b4:exm:integer-injective-resolution}的内射消解
\[
0\longrightarrow\mathbb Z\longrightarrow\mathbb Q\longrightarrow\mathbb Q/\mathbb Z\longrightarrow0
\]
计算全右导出。\(F(\mathbb Q)=0\)，而 \(F(\mathbb Q/\mathbb Z)\) 是由 \(\dfrac{1}{2}+\mathbb Z\) 生成的二阶子群。因此 \(RF(\mathbb Z)\simeq(\mathbb Z/2\mathbb Z)[-1]\)，其唯一非零同调在一次。

\(F(\mathbb Z)=0\)，故第一个比较是零对象到非零对象的态射，不是同构。\(F(\mathbb Z/2\mathbb Z)\simeq\mathbb Z/2\mathbb Z\)，而 \(LF(\mathbb Z/2\mathbb Z)=0\)，第二个比较也不是同构。\(F\) 左正合，所以模上的 \(H^0RF\) 恢复 \(F\)，但更高次仍可能非零。它不右正合：满射 \(\mathbb Z\to\mathbb Z/2\mathbb Z\) 被送到 \(0\to\mathbb Z/2\mathbb Z\)，不再满射，因而不能期待 \(H^0LF\) 恢复 \(F\)。这与\cref{b4:prop:additive-right-derived-zero}中右正合函子的现象相对照。
\end{solution}

\begin{exercise}\label{b4:ex:18-derived-annihilation}
固定正整数 \(m\)，对任意交换群 \(N\)，用两项复形表示 \(\mathbf R\operatorname{Hom}_{\mathbb Z}(\mathbb Z/m\mathbb Z,N)\)。证明在这个导出对象上乘 \(m\) 的态射为零，并给出同伦。
\end{exercise}
\begin{solution}
自由替换 \(\mathbb Z\xrightarrow{m}\mathbb Z\) 位于 \(-1,0\) 次。其 Hom 复形在 \(0,1\) 次为 \(N\xrightarrow{-m}N\)；把一次项乘 \(-1\)，可写成 \(C=(N\xrightarrow{m}N)\)。取唯一可能的同伦分量 \(h^1:C^1\to C^0\) 为恒等，其余为零，则 \(\mathrm{d}h+h\mathrm{d}=m1_C\)。所以乘 \(m\) 在同伦范畴中已为零，而不只是对每个 Ext 群诱导零映射。
\end{solution}

\begin{exercise}\label{b4:ex:18-integer-derived-tensor}
设 \(m,n>0\)，\(d=\gcd(m,n)\)。求 \(\mathbb Z/m\mathbb Z\otimes_{\mathbb Z}^{\mathbf L}\mathbb Z/n\mathbb Z\) 的一个复形代表及全部同调，并判定它何时为零对象。
\end{exercise}
\begin{solution}
代表为 \(\mathbb Z/n\mathbb Z\xrightarrow{m}\mathbb Z/n\mathbb Z\)，次数 \(-1,0\)。乘 \(m\) 的核由 \(n/d\) 的类生成，阶为 \(d\)，余核也为 \(\mathbb Z/d\mathbb Z\)。因此 \(H^{-1}=H^0=\mathbb Z/d\mathbb Z\)，其余为零。导出对象为零当且仅当 \(d=1\)。当 \(d>1\) 时，仅列出普通张量 \(\mathbb Z/d\mathbb Z\) 不足以表示该导出对象。
\end{solution}

\begin{exercise}\label{b4:ex:18-tensor-shift-sign}
设 \(R\) 为环，\(C\) 为右 \(R\)-模复形，\(D\) 为左 \(R\)-模复形，\(a,b\in\mathbb Z\)。构造复形同构
\(C[a]\otimes_RD[b]\simeq(C\otimes_RD)[a+b]\)，并核对符号。
\end{exercise}
\begin{solution}
若 \(x\in C[a]^p=C^{p+a}\)、\(y\in D[b]^q=D^{q+b}\)，定义 \(x\otimes y\mapsto(-1)^{bp}x\otimes y\)。两边总次数都为 \(p+q\)。先在左边微分再映射，两个系数分别为 \((-1)^{a+b(p+1)}\) 与 \((-1)^{p+b+bp}\)。先映射再在右边微分，系数分别为 \((-1)^{bp+a+b}\) 与 \((-1)^{bp+a+b+p+a}\)，逐项相同。再乘相同的符号即为逆。采用投射或平坦替换后得到相同的导出移位同构。
\end{solution}

\begin{exercise}\label{b4:ex:18-extension-tower}
设 \(R\to S\to T\) 为含幺环同态，\(X\) 为有上界的左 \(R\)-模复形。证明
\[
T\otimes_S^{\mathbf L}(S\otimes_R^{\mathbf L}X)
\simeq T\otimes_R^{\mathbf L}X,
\]
其中每次均按标量扩张保留左侧环的作用，不假定平坦性。
\end{exercise}
\begin{solution}
取 \(P\to X\) 的有上界投射替换。\(S\otimes_RP\) 逐项投射于 \(S\)，所以外层导出张量无需再次替换。左侧由 \(T\otimes_S(S\otimes_RP)\) 表示，普通结合映射
\(t\otimes(s\otimes p)\mapsto ts\otimes p\)
给出到 \(T\otimes_RP\) 的复形同构，逆为 \(t\otimes p\mapsto t\otimes(1\otimes p)\)。两边已是各自导出计算，故无须要求 \(S\) 或 \(T\) 平坦。
\end{solution}

\begin{exercise}\label{b4:ex:18-integer-derived-hom}
设 \(m,n>0\)。计算 \(\mathbf R\operatorname{Hom}_{\mathbb Z}(\mathbb Z/m\mathbb Z,\mathbb Z/n\mathbb Z)\)，并与上一道整数导出张量题比较非零次数。
\end{exercise}
\begin{solution}
自由替换给出 \(\mathbb Z/n\mathbb Z\xrightarrow{-m}\mathbb Z/n\mathbb Z\)，次数 \(0,1\)。故零次为乘 \(m\) 的核，一次为余核，均同构于 \(\mathbb Z/\gcd(m,n)\mathbb Z\)，其他次数为零。张量的修正项位于 \(-1\) 次，Hom 的修正项位于 \(1\) 次；这来自 Hom 的第一变量反变，而不是两者采用不同的微分升降约定。
\end{solution}

\begin{exercise}\label{b4:ex:18-finite-projective-bidual}
设 \(R\) 为交换环，\(P\) 为有限生成投射 \(R\)-模的有界复形，\(P^\vee=\operatorname{Hom}_R^\bullet(P,R)\)。证明带符号的求值给出 \(P\simeq(P^\vee)^\vee\)，并解释相应的导出双对偶同构。
\end{exercise}
\begin{solution}
对 \(x\in P^p\)、\(f\in(P^\vee)^{-p}\)，定义 \(\eta(x)(f)=(-1)^pf(x)\)。因为底复形 \(R\) 只有零次，\(\mathrm{d}_{P^\vee}f=-(-1)^{|f|}f\mathrm{d}_P\)。对次数 \(-p-1\) 的 \(g\)，代入两次对偶微分得到
\((\mathrm{d}\eta(x))(g)=(-1)^{p+1}g(\mathrm{d}_Px)=\eta(\mathrm{d}_Px)(g)\)。所以 \(\eta\) 是复形映射。各项有限生成投射，通常的求值到双对偶为同构，乘符号不改变可逆性。\(P\) 及 \(P^\vee\) 都有界投射，所以可直接计算两次导出 Hom，得到 \(P\simeq\mathbf R\operatorname{Hom}_R(\mathbf R\operatorname{Hom}_R(P,R),R)\)。
\end{solution}

\begin{exercise}\label{b4:ex:18-tensor-cone}
设 \(R\) 为交换环，\(f:P\to Q\) 是有上界投射复形间的映射，\(U\) 为有上界投射复形。证明 \(\operatorname{Cone}(f)\otimes_RU\) 自然同构于 \(\operatorname{Cone}(f\otimes1_U)\)，并推出固定一个变量的导出张量保持好三角。
\end{exercise}
\begin{solution}
按分量识别两边为 \((Q\otimes U)\oplus(P\otimes U)[1]\)。若 \(x\in P^{p+1}\)、\(u\in U^q\)，源中把 \(x\) 看作 \(P[1]^p\)，其内部微分为 \(-\mathrm{d}_Px\otimes u+(-1)^px\otimes \mathrm{d}_Uu\)，恰为目标移位总微分。跨分量映射均为 \(f(x)\otimes u\)，故该直接识别就是复形同构。导出范畴每个好三角同构于锥三角；先在投射模型中表示，再作上述张量，所得仍为锥三角。移位的相容性由上一移位题的公式核对。
\end{solution}

\begin{exercise}\label{b4:ex:18-restriction-not-full}
设 \(S=\mathbb Z/2\mathbb Z\)。对商环映射 \(\mathbb Z\twoheadrightarrow S\)，通常模范畴中的限制函子 \(S\text{-}\operatorname{Mod}\to\mathbb Z\text{-}\operatorname{Mod}\) 全忠实。证明导出限制函子在有界导出范畴中未必全忠实。
\end{exercise}
\begin{solution}
商环映射满射，所以两个 \(S\)-模之间的每个 \(\mathbb Z\)-线性映射自动保持 \(S\)-作用，通常的限制函子因而全忠实。再取两个导出对象 \(S\) 和 \(S[1]\)。在 \(D(S)\) 中，\(\operatorname{Hom}(S,S[1])=\operatorname{Ext}_S^1(S,S)=0\)，因为 \(S\) 是自由 \(S\)-模。在 \(D(\mathbb Z)\) 中，同一态射群为 \(\operatorname{Ext}_{\mathbb Z}^1(S,S)=S\ne0\)。其中的非零类由短正合列
\[
0\longrightarrow\mathbb Z/2\mathbb Z
\xrightarrow{\ \bar1\mapsto\bar2\ }\mathbb Z/4\mathbb Z
\longrightarrow\mathbb Z/2\mathbb Z\longrightarrow0
\]
表示，最后一个映射是模 \(2\)。它不分裂，因为 \(\bar1\) 的每个提升都是四阶元，不可能是二阶元的像；中间项也不是 \(S\)-模，因为乘 \(2\) 在其中不为零。有界嵌入全忠实，故这两个计算也分别是 \(D^b(S)\)、\(D^b(\mathbb Z)\) 中的态射群。因此导出限制函子在这对对象上不满。零次模同态没有改变，但允许扩张以后，底环提供了新的态射。
\end{solution}

\begin{exercise}\label{b4:ex:18-bounded-not-tensor-closed}
设 \(k\) 为域，\(R=k[\varepsilon]/(\varepsilon^2)\)。证明 \(D^b(R\text{-}\operatorname{Mod})\) 对导出张量未必封闭。再给出保证 \(X\otimes_R^{\mathbf L}Y\) 有界的一项充分条件，其中 \(X,Y\) 为 \(R\)-模上链复形，\(Y\) 有界。
\end{exercise}
\begin{solution}
取 \(X=Y=k\)，二者只有零次同调，故属于有界导出范畴。但\cref{b4:exm:dual-numbers-derived-base-change}给出每个 \(n\ge0\) 的 \(H^{-n}(k\otimes_R^{\mathbf L}k)=k\)，结果不有界。一个充分条件是 \(X\) 同构于有界平坦复形 \(F\)：此时 \(F\otimes_RY\) 计算导出张量，两个有界次数区间之和仍是有限区间，所以结果有界。
\end{solution}

\begin{exercise}\label{b4:ex:18-cohomology-does-not-determine}
设 \(k\) 为域，\(R=k[\varepsilon]/(\varepsilon^2)\)，\(C=(R\xrightarrow{\varepsilon}R)\) 位于 \(-1,0\) 次，\(T=k[1]\oplus k\)，其中 \(k=R/(\varepsilon)\)。对每个整数 \(n\)，计算 \(\operatorname{Hom}_{D(R)}(C,k[n])\) 与 \(\operatorname{Hom}_{D(R)}(T,k[n])\) 这两个 \(k\)-向量空间，并用结果区分 \(C\) 与 \(T\)。
\end{exercise}
\begin{solution}
\(C\) 有界自由，可以直接用 \(\operatorname{Hom}_R^\bullet(C,k)\) 计算。这是位于 \(0,1\) 次的两项 \(k\)，微分为零，因为 \(\varepsilon\) 在 \(k\) 上作用为零。所以
\[
\operatorname{Hom}_{D(R)}(C,k[n])\simeq
\begin{cases}k,&n=0,1,\\0,&\text{其他 }n.\end{cases}
\]
由\cref{b4:exm:dual-numbers-derived-base-change}中的周期自由消解，\(\operatorname{Ext}_R^j(k,k)=k\) 对每个 \(j\ge0\) 成立；负次导出 Hom 为零。移位第一变量给出
\[
\operatorname{Hom}_{D(R)}(T,k[n])
\simeq\operatorname{Ext}_R^{n-1}(k,k)\oplus\operatorname{Ext}_R^n(k,k)
\simeq\begin{cases}0,&n<0,\\k,&n=0,\\k\oplus k,&n\ge1.\end{cases}
\]
在 \(n=1\) 时两者的维数不同，故 \(C\not\simeq T\)。这用移位态射群重新检验了\cref{b4:prop:derived-cohomology-not-split}：这里同调仅在两度非零，问题是这些同调之间的扩张信息；前一题则是两个有界对象的导出张量产生了无限多度同调。
\end{solution}

\begin{exercise}\label{b4:ex:18-rational-vanishing}
设 \(M\) 为交换群。证明 \(\mathbb Q\otimes_{\mathbb Z}^{\mathbf L}M\) 为零对象，当且仅当 \(M\) 的每个元素都有有限阶。
\end{exercise}
\begin{solution}
\(\mathbb Q\) 是 \(\mathbb Z\) 关于非零整数的局部化，所以平坦，导出张量就是零次的 \(\mathbb Q\otimes M\)。局部化的零判据说明 \(1\otimes m=0\) 当且仅当存在非零整数 \(s\) 使 \(sm=0\)。这些元素生成张量积，故它为零当且仅当所有 \(m\) 都为挠元。此处不要求存在同时零化整个群的同一个整数。
\end{solution}

\begin{optionalexercise}\label{b4:ex:18-acyclic-k-projective}
证明正合的 K-投射复形可缩，正合的 K-内射复形也可缩。反过来，证明每个可缩模复形都同时具有 K-投射、K-内射与适当侧别的 K-平坦性。
\end{optionalexercise}
\begin{solution}
若 \(P\) 正合且 K-投射，在定义中取 \(A=P\)，则 \(\operatorname{Hom}_R^\bullet(P,P)\) 正合。其零次余圈 \(1_P\) 因而是余边，即存在次数 \(-1\) 的 \(h\) 使
\[
1_P=\mathrm{d}_{\operatorname{Hom}}h=\mathrm{d}_Ph+h\mathrm{d}_P.
\]
这就是收缩同伦。K-内射时在定义中取 \(A=I\)，同理得到 \(1_I=\mathrm{d}_Ih+h\mathrm{d}_I\)。反向若 \(1_C=\mathrm{d}h+h\mathrm{d}\)，与任意映射复合即可使从 \(C\) 出发或到 \(C\) 的复形映射零同伦，对各次移位同样成立，所以相应 Hom 复形正合。对右模复形 \(A\)，\(A\otimes C\) 的收缩为 \(a\otimes c\mapsto(-1)^{|a|}a\otimes h(c)\)；代入总微分，两个混合项相消，剩下恒等，故 K-平坦。
\end{solution}

\begin{optionalexercise}\label{b4:ex:18-acyclic-k-flat}
设 \(R\) 为交换环，\(F\) 是正合且 K-平坦的复形。证明对任意模复形 \(Y\)，\(F\otimes_RY\) 都正合。为什么仅假定 \(F\) 逐项平坦不够？
\end{optionalexercise}
\begin{solution}
取半自由替换 \(Q\to Y\)。\(Q\) 是 K-平坦，而 \(F\) 正合，故 \(F\otimes_RQ\) 正合。另一方面，\(F\) 的 K-平坦性使 \(F\otimes_RQ\to F\otimes_RY\) 为拟同构，故目标正合。对\cref{b4:exm:unbounded-free-not-k-projective}的双向周期复形 \(E\)，各项自由且 \(E\) 正合，但 \(E\otimes_Rk\) 零微分、每度非零，因此逐项平坦并不足够。
\end{solution}

\begin{optionalexercise}\label{b4:ex:18-field-unbounded}
设 \(k\) 是域。证明任意无界 \(k\)-向量空间复形都同时是 K-投射、K-内射和 K-平坦的，并说明这一结论为什么不能推广到任意环。
\end{optionalexercise}
\begin{solution}
对任意正合向量空间复形 \(A\)，逐次选补空间 \(A^n=Z^nA\oplus L^n\)。微分把 \(L^n\) 同构到 \(Z^{n+1}A\)，其逆在 \(Z^{n+1}A\) 上定义收缩，补空间上取零，所以 \(A\) 可缩。任意复形 \(X\) 与这个收缩复合，表明 \(\operatorname{Hom}(X,A)\)、\(\operatorname{Hom}(A,X)\) 可缩；张量后也由带分次符号的同伦可缩。因而三个 K-条件都成立。一般环上的正合列未必分裂，正合复形未必可缩；双数代数的双向周期复形已经给出失败例子。
\end{solution}

% !TeX root = ../../Book4.tex
% 纲要标题：## 第20章*　DG 范畴与导出等价
% 纲要位置：计划纲要.md，第 3024 行。

\OptionalChapter{DG 范畴与导出等价}{b4:ch:20}
\BookChapterNumber{b4:ch:20}

\begin{chapterabstract}
导出范畴把拟同构变成同构，也把许多具体同伦隐藏在态射的表示之中。本章保留整个 Hom 复形，建立 DG 代数、DG 范畴及增强的基本语言，再讨论有限投射数据如何在导出范畴中表现为完美复形。倾斜对象把自身态射的消失与厚生成结合起来，由此引出导出 Morita 理论。三个顶点的路径代数给出一个可以逐项计算的例子：模范畴并不等价，导出范畴却等价。最后从锥映射的同伦选择和一个导出张量计算说明进一步学习所需的问题与语言。
\end{chapterabstract}

\begin{learninggoals}
\begin{itemize}
\item 正确展开 DG 乘法、Hom 微分与复合中的符号，区分 DG 微分与普通指定导子。
\item 说明 \(H^0\) 忘掉了哪些资料，并在有界投射复形中识别一个具体增强。
\item 区分完美性、同调有界、每项有限生成和紧性，核对厚生成的有限构造。
\item 在三个顶点的表示上计算倾斜条件、自同态代数以及缺失投射模的恢复。
\item 说明紧对象刻画中怎样提取有限模型，以及导出 Morita 等价中双模替换的作用。
\end{itemize}
\end{learninggoals}

本章的基本计算使用\cref{b4:thm:bounded-projective-model,b4:thm:derived-ext-identification}及\chapref{b4:ch:18}的导出 Hom、张量积；判断三顶点例子不构成普通 Morita 等价时，还调用第二卷定理15.3.10的不可分解投射模分类。与不同心的比较只涉及\chapref{b4:ch:19}的主题，不调用其中的倾斜构造。Hom 复形、\subsecref{b4:subsec:2002:01}至\subsecref{b4:subsec:2002:02}的有限完美模型，以及三顶点例子的扩张和自同态计算均不依赖\subsecref{b4:subsec:2002:03}；正式判定该例为倾斜对象则使用该小节的完美复形厚性。该小节对紧对象的反向刻画还依赖\secref{b4:sec:1804}的半自由消解。第三卷附录 E 仅用于对照“指定导子”与“DG 微分”两种结构，并非本章前提。

\ProseParagraphBreak
Keller 的《Deriving DG categories》\cite[Sections 1.1--2.2]{Keller1994DerivingDG}给出分次态射、DG 模和同伦范畴的约定；紧对象及导出等价分别见\cite[Section 5.3; Sections 8.2 and 9.1--9.2]{Keller1994DerivingDG}。普通环的完美复形和 Rickard 定理也可查其《Introduction to abelian and derived categories》\cite[Sections 3.1 and 4]{Keller1998AbelianDerived}。本章使用前一论文时同时采用作者的勘误：无界消解中的零调判据须保留圈复形条件，DG 复形的平坦性须按保持无界零调复形来理解。

% !TeX root = ../../Book4.tex
% 纲要标题：### 20.1 标准截断
% 纲要标签：b4:sec:1901
% 纲要位置：计划纲要.md，第 3742 行。

\section{标准截断}
\label{b4:sec:1901}

截去复形的一端会改变截口处的核或像。典范截断修正端点项，既保留指定范围的上同调，又将原对象放入一个截断三角。


% 纲要内容（仅作写作计划，尚非教材正文）：
%
% * 朴素截断与典范截断
% * 标准 t-结构
% * 截断三角
%

\subsection{朴素截断与典范截断}
\label{b4:subsec:1901:01}

把复形截成两段，首先遇到的问题是截口的微分。直接删掉高次项会扩大最后一项的圈；直接删掉低次项则会丢掉原来进入第一项的边界。这两种操作适合描述复形的项，却不能原样保留指定次数的上同调。

\begin{definition}\label{b4:def:standard-truncations}
设 \(C\) 为交换范畴 \(\mathcal A\) 中的上链复形，\(n\in\mathbb Z\)。保留次数 \(i\le n\) 的原项与原微分、把其余项置零，所得复形记为 \(\sigma^{\le n}C\)；相应地定义 \(\sigma^{\ge n}C\)。这些构造称为\term[pusujianduan]{朴素截断}[brutal truncation]。把截口改为
\[
 (\tau^{\le n}C)^i=
 \begin{cases}C^i&i<n,\\ \ker \mathrm{d}^n&i=n,\\0&i>n,\end{cases}
 \qquad
 (\tau^{\ge n}C)^i=
 \begin{cases}0&i<n,\\ \operatorname{coker}\mathrm{d}^{n-1}&i=n,\\C^i&i>n,\end{cases}
\]
并采用原微分诱导的映射，得到\term[dianfanjianduan]{典范截断}[canonical truncation]，也称好截断。它们分别带有自然链映射 \(\tau^{\le n}C\to C\) 和 \(C\to\tau^{\ge n}C\)。
\symidx{\(\tau^{\le n},\tau^{\ge n}\)：典范截断函子}
\end{definition}
\symidx{\(\sigma^{\le n},\sigma^{\ge n}\)：上链复形的朴素截断}

复形映射在截口的核上限制、在余核上诱导映射，因而两种典范截断都是复形范畴上的函子。它们也保持链同伦。具体地，设 \(f,g:C\to D\) 及次数为 \(-1\) 的 \(h\) 满足
\[
f^i-g^i=\mathrm d_D^{i-1}h^i+h^{i+1}\mathrm d_C^i.
\]
记 \(\iota_C^n:Z^n(C)\hookrightarrow C^n\) 为核包含。对 \(\tau^{\le n}\)，低于截口的同伦仍取 \(h^i\)，在截口取
\(\bar h^n=h^n\iota_C^n:Z^n(C)\to D^{n-1}\)，其余取零。因为
\(\mathrm d_C^n\iota_C^n=0\)，原同伦等式在截口变为
\[
(f^n-g^n)\iota_C^n
=\mathrm d_D^{n-1}h^n\iota_C^n.
\]
两边都经 \(Z^n(D)\) 因子化，核包含的单射性便给出截断后的同伦等式；在 \(n-1\) 次，微分经 \(Z^n(C)\) 的因子化也使原等式保持。

\ProseParagraphBreak
对 \(\tau^{\ge n}\)，记 \(\pi_C^n:C^n\twoheadrightarrow\operatorname{coker}\mathrm d_C^{n-1}\)、\(\pi_D^n:D^n\twoheadrightarrow\operatorname{coker}\mathrm d_D^{n-1}\) 为商映射。同伦在 \(i\le n\) 时取零，在 \(n+1\) 次取
\(\bar h^{n+1}=\pi_D^n h^{n+1}\)，在更高次数仍取 \(h^i\)。若 \(\bar f^n,\bar g^n,\bar{\mathrm d}_C^n\) 表示截口的诱导映射，则
\[
\begin{aligned}
(\bar f^n-\bar g^n)\pi_C^n
&=\pi_D^n(f^n-g^n)\\
&=\pi_D^n h^{n+1}\mathrm d_C^n
=\bar h^{n+1}\bar{\mathrm d}_C^n\pi_C^n.
\end{aligned}
\]
这里 \(\pi_D^n\mathrm d_D^{n-1}=0\) 消去了另一项；\(\pi_C^n\) 为满态射，所以截口的同伦等式成立。在 \(n+1\) 次，恒等式
\(\bar{\mathrm d}_D^n\pi_D^n=\mathrm d_D^n\) 恢复原来的同伦等式，其他次数不变。因此两种截断都定义了 \(K(\mathcal A)\) 上的函子。

\ProseParagraphBreak
上截断在截口的核是整个 \(\ker\mathrm{d}_C^n\)，进入它的像仍为 \(\operatorname{im}\mathrm{d}_C^{n-1}\)，所以
\[
H^n(\tau^{\le n}C)
=\frac{\ker\mathrm{d}_C^n}{\operatorname{im}\mathrm{d}_C^{n-1}}
=H^n(C).
\]
下截断在截口使用 \(C^n/\operatorname{im}\mathrm{d}_C^{n-1}\)。其诱导微分
\(\bar{\mathrm{d}}^n:C^n/\operatorname{im}\mathrm{d}_C^{n-1}\rightarrow C^{n+1}\)
的核为 \(\ker\mathrm{d}_C^n/\operatorname{im}\mathrm{d}_C^{n-1}\)，而截口前一项为零，故
\[
H^n(\tau^{\ge n}C)=\ker\bar{\mathrm{d}}^n\cong H^n(C).
\]
其他保留次数的核与像不变，被删去次数的上同调为零。因此上截断保留 \(i\le n\) 的上同调，下截断保留 \(i\ge n\) 的上同调；核与余核的端点修正正是为了保留截口的原上同调。它们由此保持拟同构，将上述同伦范畴上的函子与局部化函子复合，依\cref{b4:def:derived-category}的泛性质得到导出范畴上的截断函子。

\ProseParagraphBreak
例如 \(C=(\mathbb Z\xrightarrow{2}\mathbb Z)\) 位于次数 \(-1,0\)。朴素截断 \(\sigma^{\ge0}C=\mathbb Z[0]\)，而典范截断 \(\tau^{\ge0}C=(\mathbb Z/2\mathbb Z)[0]\)。前者忘记了进入零次项的边界，后者才保留 \(H^0(C)\)。

\subsection{标准 t-结构}
\label{b4:subsec:1901:02}

\begin{definition}\label{b4:def:standard-t-halves}
设 \(\mathcal A\) 为交换范畴，\(D(\mathcal A)\) 为其导出范畴，\(n\in\mathbb Z\)。令
\[
 D^{\le n}(\mathcal A)=\{X:H^i(X)=0\ (i>n)\},\qquad
 D^{\ge n}(\mathcal A)=\{X:H^i(X)=0\ (i<n)\}
\]
表示由相应对象组成的全子范畴。它们也分别由同构于次数不超过 \(n\)、次数不低于 \(n\) 的复形的对象组成。\((D^{\le0},D^{\ge0})\) 称为\term[biaozhuntjiegou]{标准 t-结构}[standard t-structure]；下一节将抽取它满足的公理。
\symidx{\(D^{\le n},D^{\ge n}\)：标准上同调次数条件}
\end{definition}

\begin{proposition}\label{b4:prop:standard-truncation-orthogonality}
设 \(\mathcal A\) 为局部小交换范畴，\(n\in\mathbb Z\)。若 \(X\in D^{\le n}(\mathcal A)\)、\(Y\in D^{\ge n+1}(\mathcal A)\)，则 \(\operatorname{Hom}_{D(\mathcal A)}(X,Y)=0\)。
\end{proposition}
\begin{proof}
以典范截断把 \(X,Y\) 分别换成项集中在 \(i\le n\)、\(i\ge n+1\) 的复形。这只能说明直接的链映射 \(X\to Y\) 为零；导出态射还可能经由屋顶的中间复形，而它未必有相同的项次数界。用屋顶 \(X\xleftarrow{s}Z\xrightarrow{f}Y\) 表示任意导出态射。因为 \(s\) 为拟同构且 \(X\) 的高次上同调为零，\(\tau^{\le n}Z\to Z\) 也是拟同构。用它细化得到
\[
X\xleftarrow{\ \simeq\ }\tau^{\le n}Z\longrightarrow Y.
\]
此时中间复形的项才集中在 \(i\le n\)，向 \(Y\) 的链映射逐项为零：源、靶没有共同的非零次数。因此每个屋顶都表示零态射，正交性成立。
\end{proof}

这个零性要求源在低次、靶在高次。反向一般不为零：把零次对象 \(M\) 移位为 \(M[1]\)，它位于 \(-1\) 次，而 \(\operatorname{Hom}(N,M[1])\) 可以记录一次扩张。这里的次数判断始终使用 \(C[1]^i=C^{i+1}\)。

\subsection{截断三角}
\label{b4:subsec:1901:03}

\begin{proposition}\label{b4:prop:standard-truncation-triangle}
设 \(C\) 为交换范畴 \(\mathcal A\) 中的复形。对每个整数 \(n\)，导出范畴中有自然的好三角
\[
 \tau^{\le n}C\longrightarrow C\longrightarrow
 \tau^{\ge n+1}C\longrightarrow(\tau^{\le n}C)[1].
\]
两端分别属于 \(D^{\le n}\) 和 \(D^{\ge n+1}\)。
\end{proposition}
\begin{proof}
取逐项短正合列 \(0\to\tau^{\le n}C\to C\to Q\to0\)。商复形 \(Q\) 在 \(n\) 次为 \(\operatorname{im}\mathrm{d}^n\)，在较高次数与 \(C\) 相同，较低次数为零。从 \(Q\) 到 \(\tau^{\ge n+1}C\) 的映射在 \(n\) 次为零，在 \(n+1\) 次取商，在更高次数为恒等。它的核是两项恒等复形 \(\operatorname{im}\mathrm{d}^n\to\operatorname{im}\mathrm{d}^n\)，所以是拟同构。由\cref{b4:prop:derived-short-exact-triangle}得到的三角，把 \(Q\) 换成这一同构对象即得结论。所有映射来自核、商及原链映射，故对 \(C\) 自然。
\end{proof}

三角的第三箭头不应省去。若 \(C\) 只在 \(n,n+1\) 两次有上同调，该箭头仍可能阻止它分裂为两个上同调对象的直和。截断提供逐层分析的方法，并未断言上同调完全决定复形的导出同构类型。标准截断与下面的抽象构造可对照 Beilinson、Bernstein、Deligne 的《Faisceaux pervers》\cite[Exemples 1.3.2; Proposition 1.3.3]{BBD1982FaisceauxPervers}。

% !TeX root = ../../Book4.tex
% 纲要标题：### 20.2 抽象 t-结构
% 纲要标签：b4:sec:1902
% 纲要位置：计划纲要.md，第 3748 行。

\section{抽象 t-结构}
\label{b4:sec:1902}

截断的核心数据是两类对象之间的正交性和每个对象的分解三角。将这些条件抽象出来，就能在未指定复形代表的三角范畴中恢复上同调函子。


% 纲要内容（仅作写作计划，尚非教材正文）：
%
% * 两个子范畴及正交、分解公理
% * 有界性与非退化性
% * 由截断恢复上同调函子
%

\subsection{t-结构的正交与分解公理}
\label{b4:subsec:1902:01}

标准截断的存在不必靠复形的每一项描述。上节的两类对象、它们之间的态射零性，以及把任意对象分为两部分的三角，已经足以决定截断函子。

\begin{definition}\label{b4:def:t-structure}
设 \(\mathcal D\) 为三角范畴。取包含零对象且在同构下封闭的两个全子范畴 \(\mathcal D^{\le0},\mathcal D^{\ge0}\)，并记 \(\mathcal D^{\le n}=\mathcal D^{\le0}[-n]\)、\(\mathcal D^{\ge n}=\mathcal D^{\ge0}[-n]\)。若满足
\begin{enumerate}
\item \(\mathcal D^{\le0}[1]\subseteq\mathcal D^{\le0}\)，\(\mathcal D^{\ge0}[-1]\subseteq\mathcal D^{\ge0}\)；
\item \(\operatorname{Hom}(U,V)=0\)，其中 \(U\in\mathcal D^{\le0}\)、\(V\in\mathcal D^{\ge1}\)；
\item 每个 \(X\in\mathcal D\) 都能放入好三角 \(U\to X\to V\to U[1]\)，其中 \(U\in\mathcal D^{\le0}\)、\(V\in\mathcal D^{\ge1}\)，
\end{enumerate}
则称这对子范畴为 \(\mathcal D\) 上的一个\term[tjiegou]{t-结构}[t-structure]。记 \(\mathcal D^{[a,b]}=\mathcal D^{\ge a}\cap\mathcal D^{\le b}\)，其中 \(a,b\in\mathbb Z\)。交 \(\mathcal H=\mathcal D^{\le0}\cap\mathcal D^{\ge0}\) 称为它的\term[xin]{心}[heart]。
\end{definition}
\symidx{\(\mathcal D^{\le n},\mathcal D^{\ge n}\)：给定 t-结构的次数子范畴}
\symidx{\(\mathcal D^{[a,b]}\)：给定 t-结构中次数位于 \([a,b]\) 的子范畴}
\symidx{\(\mathcal H\)：给定 t-结构的心}

移位方向由 \(C[1]^i=C^{i+1}\) 决定。若 \(M\) 是交换范畴中的对象，把它看成标准心中的零次复形，则
\[
 H^{-1}(M[1])=M,\qquad H^0(M[0])=M,\qquad H^1(M[-1])=M,
\]
且各对象的其他上同调均为零。因此 \([1]\) 将这一层移到负一次，\([-1]\) 将它移到一次；定义中的 \([-n]\) 正是把零次条件移到第 \(n\) 次。

\ProseParagraphBreak
例如，设 \(k\) 为域，\(\mathcal D=D^b(\mathbf{Vect}_k^{\mathrm{fd}})\)。标准 t-结构的心由集中在零次的有限维向量空间组成。把两个标准子范畴同时平移为 \((\mathcal D^{\le1},\mathcal D^{\ge1})\)，仍得到 t-结构，其心由 \(V[-1]\) 组成。于是同一个 \(k[-1]\) 在标准心之外，却在平移后的心中。这说明“哪些对象作为零次对象”取决于 t-结构的选择，而三角范畴本身并未改变。

\begin{theorem}\label{b4:thm:t-truncation-adjunction}
设 \(\mathcal D\) 为带 t-结构 \((\mathcal D^{\le0},\mathcal D^{\ge0})\) 的三角范畴，\(n\in\mathbb Z\)，并以 \([-n]\) 平移定义 \(\mathcal D^{\le n},\mathcal D^{\ge n}\)。包含 \(\mathcal D^{\le n}\to\mathcal D\) 有右伴随 \(\tau^{\le n}\)，包含 \(\mathcal D^{\ge n+1}\to\mathcal D\) 有左伴随 \(\tau^{\ge n+1}\)。每个对象 \(X\in\mathcal D\) 有截断三角
\[
 \tau^{\le n}X\longrightarrow X\longrightarrow\tau^{\ge n+1}X
 \longrightarrow(\tau^{\le n}X)[1],
\]
它在保持中间对象及两端次数条件的意义下唯一到唯一同构，并对 \(X\) 的态射函子性地变化。
\end{theorem}
\begin{proof}
低端截断要使每个来自 \(T\in\mathcal D^{\le n}\) 的态射 \(T\to X\) 唯一经过 \(\tau^{\le n}X\to X\)；高端截断要使每个指向 \(W\in\mathcal D^{\ge n+1}\) 的态射 \(X\to W\) 唯一经过 \(X\to\tau^{\ge n+1}X\)。换言之，所需的自然双射分别是
\[
\begin{aligned}
 \operatorname{Hom}(T,\tau^{\le n}X)&\xrightarrow{\sim}\operatorname{Hom}(T,X),\\
 \operatorname{Hom}(\tau^{\ge n+1}X,W)&\xrightarrow{\sim}\operatorname{Hom}(X,W).
\end{aligned}
\]
前一式把截断放在 Hom 的第二变量，给出包含函子的右伴随；后一式把截断放在第一变量，给出包含函子的左伴随。

平移后只须处理 \(n=0\)。取分解公理中的三角 \(U\to X\to V\)。对 \(T\in\mathcal D^{\le0}\)，由\cref{b4:prop:triangle-hom-exact}得到的 Hom 长正合列在 \(\operatorname{Hom}(T,U)\to\operatorname{Hom}(T,X)\) 两旁的项为 \(\operatorname{Hom}(T,V[-1])\) 与 \(\operatorname{Hom}(T,V)\)。它们为零，因为 \(V[-1]\in\mathcal D^{\ge2}\subseteq\mathcal D^{\ge1}\)。所以该映射是同构，给出右伴随的泛性质。对 \(W\in\mathcal D^{\ge1}\)，在 \(\operatorname{Hom}(V,W)\to\operatorname{Hom}(X,W)\) 两旁出现 \(\operatorname{Hom}(U[1],W)\)、\(\operatorname{Hom}(U,W)\)，也都为零；这给出左伴随。

泛性质唯一确定两端对象及其结构映射。给定 \(X\to X'\)，两端态射分别由伴随唯一确定；三角公理 TR3 将前两项的交换方块补成三角态射，而所补的第三项因左伴随的唯一性，必须等于已经确定的态射。

对固定的中间对象，先用这些泛性质给出的唯一同构识别两端，便可把两个分解三角写成
\[
\begin{gathered}
U\xrightarrow{i}X\xrightarrow{p}V\xrightarrow{w}U[1],\\
U\xrightarrow{i}X\xrightarrow{p}V\xrightarrow{w'}U[1].
\end{gathered}
\]
TR3 首先给出三角态射 \((1_U,1_X,c)\)，其中 \(c:V\to V\) 满足 \(cp=p\)。于是 \((c-1_V)p=0\)。对第一三角应用反变 Hom 正合列，得到
\[
\operatorname{Hom}(U[1],V)\longrightarrow
\operatorname{Hom}(V,V)\longrightarrow\operatorname{Hom}(X,V).
\]
因此存在 \(a:U[1]\to V\)，使 \(c-1_V=aw\)。但 \(U[1]\in\mathcal D^{\le0}\)、\(V\in\mathcal D^{\ge1}\)，正交性给出 \(a=0\)，故 \(c=1_V\)。最后一个方块随之给出 \(w'=w\)，而同一论证也说明第三分量唯一。由此得到保持中间对象的唯一三角同构；态射上的恒等与复合则由两端伴随的唯一性保证。
\end{proof}

\begin{proposition}\label{b4:prop:t-orthogonals-extensions}
设 \(\mathcal D\) 为带 t-结构 \((\mathcal D^{\le0},\mathcal D^{\ge0})\) 的三角范畴，\(n\in\mathbb Z\)。则 \(\mathcal D^{\le n}\) 恰是 \(\mathcal D^{\ge n+1}\) 的左正交类，\(\mathcal D^{\ge n+1}\) 恰是 \(\mathcal D^{\le n}\) 的右正交类。二者均对好三角中的扩张、有限直和及直和因子封闭。
\end{proposition}
\begin{proof}
若 \(\operatorname{Hom}(T,X)=0\) 对所有 \(T\in\mathcal D^{\le n}\) 成立，伴随同构使 \(\operatorname{Hom}(\tau^{\le n}X,\tau^{\le n}X)=0\)，故 \(\tau^{\le n}X=0\)，截断三角给出 \(X\simeq\tau^{\ge n+1}X\)。反向是正交公理。左正交刻画对偶可得。对任意好三角 \(X\to Y\to Z\)，若 \(X,Z\) 都与指定的一类对象正交，Hom 长正合列表明 \(Y\) 也正交。有限直和与直和因子的正交性则由 Hom 的加性得到。
\end{proof}

\subsection{有界性与非退化性}
\label{b4:subsec:1902:02}

\begin{definition}\label{b4:def:bounded-nondegenerate-t}
设 \(\mathcal D\) 为带 t-结构 \((\mathcal D^{\le0},\mathcal D^{\ge0})\) 的三角范畴。若每个对象 \(X\in\mathcal D\) 都属于某个 \(\mathcal D^{\ge a}\cap\mathcal D^{\le b}\)，其中整数 \(a,b\) 可以依赖 \(X\)，则称该 t-结构\term[youjietjiegou]{有界}[bounded]。若
\[
 \bigcap_n\mathcal D^{\le n}=0,\qquad
 \bigcap_n\mathcal D^{\ge n}=0,
\]
则称它\term[feituihuatjiegou]{非退化}[nondegenerate]。这里等于零表示交中只有零对象的同构类。
\end{definition}

有界性蕴含非退化性。例如若 \(X\) 属于所有 \(\mathcal D^{\le n}\)，又因有界性属于某个 \(\mathcal D^{\ge a}\)，则 \(X\in\mathcal D^{\le a-1}\cap\mathcal D^{\ge a}\)，正交性迫使 \(1_X=0\)。另一侧相同。标准 t-结构在 \(D^b(\mathcal A)\) 上有界，在一般 \(D(\mathcal A)\) 上未必有界，却仍非退化：所有上同调为零的复形在导出范畴中就是零对象。

\ProseParagraphBreak
非退化性是一项实质条件。对任意非零三角范畴 \(\mathcal D\)，取 \((\mathcal D,0)\) 也满足 t-结构公理，分解三角为 \(X\xrightarrow{1}X\to0\)。它的心为零，不能从心看见原范畴的非零对象。

\subsection{由截断恢复上同调函子}
\label{b4:subsec:1902:03}

\begin{proposition}\label{b4:prop:t-truncation-commute}
设 \(\mathcal D\) 为带 t-结构 \((\mathcal D^{\le0},\mathcal D^{\ge0})\) 的三角范畴，\(a,b\in\mathbb Z\) 且 \(a\le b\)。两种截断函子的复合 \(\tau^{\ge a}\tau^{\le b}\) 与 \(\tau^{\le b}\tau^{\ge a}\) 自然同构，所得对象属于 \(\mathcal D^{\ge a}\cap\mathcal D^{\le b}\)。
\end{proposition}
\begin{proof}
要提取 \(X\) 中次数位于 \([a,b]\) 的部分，可以先取 \(B=\tau^{\le b}X\)，再去掉其中低于 \(a\) 的部分。令 \(A=\tau^{\le a-1}X\)。因为 \(A\in\mathcal D^{\le b}\)，右伴随泛性质唯一给出 \(A\to B\to X\)，其复合为 \(A\) 的截断结构态射。对每个 \(T\in\mathcal D^{\le a-1}\)，它诱导的映射 \(\operatorname{Hom}(T,A)\to\operatorname{Hom}(T,B)\) 经两端伴随同构，成为 \(\operatorname{Hom}(T,X)\) 的恒等映射。因此 \(A\to B\) 也满足 \(B\) 在 \(a-1\) 处的低端截断泛性质。

把 \(A\to B\) 补成好三角 \(A\to B\to Y\to A[1]\)。由上述泛性质，把前两项与 \(B\) 的截断三角识别；(TR3) 补出第三项的态射，\cref{b4:lem:pretriangle-hom}的二取三同构判据给出 \(Y\simeq\tau^{\ge a}B\)。所以 \(Y\in\mathcal D^{\ge a}\)。同时 \(B,A[1]\in\mathcal D^{\le b}\)，对旋转三角 \(B\to Y\to A[1]\) 使用\cref{b4:prop:t-orthogonals-extensions}，得到 \(Y\in\mathcal D^{\le b}\)。这说明 \(Y\) 确实只保留所需次数区间。

对复合 \(A\to B\to X\)，两条已有的截断三角分别以 \(\tau^{\ge a}X\)、\(\tau^{\ge b+1}X\) 为第三项。八面体公理将它们与 \(A\to B\) 的三角联系为
\[
 Y\longrightarrow\tau^{\ge a}X\longrightarrow\tau^{\ge b+1}X\longrightarrow Y[1].
\]
由于 \(Y\in\mathcal D^{\le b}\)，最后一个高端对象属于 \(\mathcal D^{\ge b+1}\)，这又是 \(\tau^{\ge a}X\) 在 \(b\) 处的截断三角，故 \(Y\simeq\tau^{\le b}\tau^{\ge a}X\)。两次识别得到所述同构。八面体中的 \(Y\to\tau^{\ge a}X\) 与 \(B\to X\to\tau^{\ge a}X\) 相容，因而也由 \(B\) 的高端截断泛性质唯一确定；各结构态射都对 \(X\) 自然，故所得同构自然。
\end{proof}

两侧截断 \(\tau^{\ge n}\tau^{\le n}X\) 将对象限制在第 \(n\) 层；它属于 \(\mathcal D^{\ge n}\cap\mathcal D^{\le n}\)，还须平移 \([n]\) 才把这一层移回零次的心。因此在未指定复形各项的情况下定义
\[
 H_t^n(X)\coloneqq\bigl(\tau^{\ge n}\tau^{\le n}X\bigr)[n]
                 =H_t^0(X[n])\in\mathcal H.
\]
这称为关于给定 t-结构的\term[tshangtongdiao]{上同调对象}[cohomology object]。例如在标准结构中，\(M[-2]\) 只有二次上同调，取两侧截断后仍是 \(M[-2]\)；再施加 \([2]\)，得到 \(M[0]\)，所以 \(H_t^2(M[-2])=M\) 是心中的对象。下标 \(t\) 用来区别随后出现的新心；标准情形可省去。上述定义先给出取值于加性心的函子；心的交换性和三角诱导的长正合列将在下一节证明。
\symidx{\(H_t^n\)：给定 t-结构的上同调函子}

\ProseParagraphBreak
相邻截断之间有好三角 \(\tau^{\le n-1}X\to\tau^{\le n}X\to H_t^n(X)[-n]\to\)，这是前一证明取 \(a=b=n\) 的情形。有界对象因此能由有限多个上同调对象逐次扩张得到。每一层的连接映射仍是重建对象所需的数据。抽象截断与交换关系的原始表述见\cite[Propositions 1.3.3, 1.3.5]{BBD1982FaisceauxPervers}。

% !TeX root = ../../Book4.tex
% 纲要标题：### 20.3 心
% 纲要标签：b4:sec:1903
% 纲要位置：计划纲要.md，第 2993 行。

\section{心}
\label{b4:sec:1903}

心中的对象集中在同一上同调位置，态射的锥却占据相邻位置。这两个位置提供核与余核，正交性再保证余像与像一致。


% 纲要内容（仅作写作计划，尚非教材正文）：
%
% * 心是交换范畴
% * 短正合列与三角的关系
% * 标准心与原交换范畴
%

\subsection{心是交换范畴}
\label{b4:subsec:1903:01}

在心中，一切对象都位于同一个上同调次数；但两个对象之间态射的锥通常占据相邻两次。这正好提供核与余核所需的两个对象。困难在于还要证明余像等于像，不能只造出两个候选对象便结束。

\begin{theorem}\label{b4:thm:heart-abelian}
设 \(\mathcal D\) 为带 t-结构 \((\mathcal D^{\le0},\mathcal D^{\ge0})\) 的三角范畴，\(\mathcal H=\mathcal D^{\le0}\cap\mathcal D^{\ge0}\) 为其心。则 \(\mathcal H\) 是交换范畴。若 \(f:X\to Y\) 是心中态射，完成为好三角
\[
 X\xrightarrow{f}Y\longrightarrow C\longrightarrow X[1],
\]
则心中的核、余核分别为
\[
 \ker_{\mathcal H}f=H_t^{-1}(C),\qquad
 \operatorname{coker}_{\mathcal H}f=H_t^0(C).
\]
这两个对象及其核、余核映射由泛性质唯一确定，不依赖锥的选择。
\end{theorem}
\begin{proof}
心是满的加性子范畴，因为\cref{b4:prop:t-orthogonals-extensions}保证有限直和仍在心中。旋转三角 \(Y\to C\to X[1]\) 表明 \(C\in\mathcal D^{\ge-1}\cap\mathcal D^{\le0}\)。令
\[
 K=(\tau^{\le-1}C)[-1],\qquad Q=\tau^{\ge0}C.
\]
两者都在心中，有截断三角 \(K[1]\to C\to Q\to K[2]\)。其与原三角给出 \(K\to X\) 及 \(Y\to Q\)。

对心中任意 \(T\)，在原三角上取 \(\operatorname{Hom}(T,-)\)，得到
\[
 0\longrightarrow\operatorname{Hom}(T,C[-1])
 \longrightarrow\operatorname{Hom}(T,X)
 \longrightarrow\operatorname{Hom}(T,Y),
\]
首项前的 \(\operatorname{Hom}(T,Y[-1])\) 为零。截断三角及 \(Q[-1],Q[-2]\) 的正次数条件又给出 \(\operatorname{Hom}(T,K)\simeq\operatorname{Hom}(T,C[-1])\)。所以 \(K\to X\) 是核。对偶地，\(\operatorname{Hom}(Q,T)\simeq\operatorname{Hom}(C,T)\)，而 \(\operatorname{Hom}(X[1],T)=0\)，因此 \(Y\to Q\) 是余核。

还须核对余像与像。把要比较的复合写成 \(Y\xrightarrow{u}C\xrightarrow{v}Q\)。原三角与截断三角分别给出
\[
 \operatorname{Cone}(u)=X[1],\qquad
 \operatorname{Cone}(v)=K[2].
\]
取 \(\operatorname{Cone}(vu)=I[1]\)。八面体公理的第四个三角便是
\[
 X[1]\longrightarrow I[1]\longrightarrow K[2]\longrightarrow X[2].
\]
将它平移 \([-1]\) 后逆旋转，再把复合的三角 \(Y\to Q\to I[1]\to Y[1]\) 逆旋转，得到相容的两个好三角
\[
 K\longrightarrow X\longrightarrow I\longrightarrow K[1],
 \qquad
 I\longrightarrow Y\longrightarrow Q\longrightarrow I[1],
\]
旋转中出现的负号按三角的约定吸收到相应箭头中，八面体的交换关系保证 \(X\to I\to Y\) 的复合为 \(f\)。第一三角旋转为 \(X\to I\to K[1]\)，两端都在 \(\mathcal D^{\le0}\)，所以扩张封闭性给出 \(I\in\mathcal D^{\le0}\)。第二三角逆旋转为 \(Q[-1]\to I\to Y\)，两端都在 \(\mathcal D^{\ge0}\)，同理得到 \(I\in\mathcal D^{\ge0}\)。故 \(I\) 在心中。在第一三角上向任意心对象取 Hom，\(\operatorname{Hom}(K[1],T)=0\) 表明 \(X\to I\) 是 \(K\to X\) 的余核。在第二三角上从心对象取 Hom，\(\operatorname{Hom}(T,Q[-1])=0\) 表明 \(I\to Y\) 是 \(Y\to Q\) 的核。于是 \(I\) 同时给出余像与像，且它们间的典范态射是同构。这验证了交换范畴的全部条件。
\end{proof}

上述八面体只用于比较两个实际构造出的三角：一边先去掉 \(f\) 的核，另一边先去掉它的余核。它所产生的共同对象 \(I\) 才使两种做法一致。原始证明可对照\cite[Proposition 1.2.4; Théorème 1.3.6]{BBD1982FaisceauxPervers}。

\begin{theorem}\label{b4:thm:t-cohomology-exact}
设 \(\mathcal D\) 为带 t-结构 \((\mathcal D^{\le0},\mathcal D^{\ge0})\) 的三角范畴，心为 \(\mathcal H=\mathcal D^{\le0}\cap\mathcal D^{\ge0}\)，\(H_t^n(X)=(\tau^{\ge n}\tau^{\le n}X)[n]\)。函子 \(H_t^0:\mathcal D\to\mathcal H\) 是上同调函子：每个好三角 \(X\to Y\to Z\xrightarrow{w}X[1]\) 诱导自然长正合列
\[
 \cdots\longrightarrow H_t^n(X)\longrightarrow H_t^n(Y)
 \longrightarrow H_t^n(Z)\xrightarrow{H_t^n(w)}H_t^{n+1}(X)
 \longrightarrow\cdots.
\]
\end{theorem}
\begin{proof}
只须证明 \(H_t^0(X)\to H_t^0(Y)\to H_t^0(Z)\) 正合，再平移并旋转三角。

先设三项都在 \(\mathcal D^{\le0}\)。对 \(T\in\mathcal H\)，伴随给出 \(\operatorname{Hom}(U,T)=\operatorname{Hom}(H_t^0(U),T)\)，其中 \(U=X,Y,Z\)。Hom 长正合列因 \(\operatorname{Hom}(X[1],T)=0\) 而给出
\[
 0\longrightarrow\operatorname{Hom}(H_t^0Z,T)
 \longrightarrow\operatorname{Hom}(H_t^0Y,T)
 \longrightarrow\operatorname{Hom}(H_t^0X,T).
\]
这正是余核泛性质，所以 \(H_t^0X\to H_t^0Y\to H_t^0Z\to0\) 正合。

再设仅有 \(X\in\mathcal D^{\le0}\)。对 \(T\in\mathcal D^{\ge1}\)，正交性使 \(\operatorname{Hom}(Z,T)\to\operatorname{Hom}(Y,T)\) 同构。因此 \(\tau^{\ge1}Y\to\tau^{\ge1}Z\) 是同构。把 \(Y\to Z\to\tau^{\ge1}Z\) 代入八面体公理，其余两锥分别是截断的低次部分的移位，得到好三角
\[
 X\longrightarrow\tau^{\le0}Y\longrightarrow\tau^{\le0}Z\longrightarrow X[1].
\]
应用前一情形得到 \(H_t^0X\to H_t^0Y\to H_t^0Z\to0\) 正合。

再证另一端的正合性。若三项都在 \(\mathcal D^{\ge0}\)，对 \(T\in\mathcal H\)，截断伴随给出
\(\operatorname{Hom}(T,U)\simeq\operatorname{Hom}(T,H_t^0U)\)，其中 \(U=X,Y,Z\)。又 \(Z[-1]\in\mathcal D^{\ge1}\)，故 \(\operatorname{Hom}(T,Z[-1])=0\)。Hom 长正合列于是给出
\[
0\longrightarrow\operatorname{Hom}(T,H_t^0X)
\longrightarrow\operatorname{Hom}(T,H_t^0Y)
\longrightarrow\operatorname{Hom}(T,H_t^0Z).
\]
这正是核泛性质，所以 \(0\to H_t^0X\to H_t^0Y\to H_t^0Z\) 正合。

若仅有 \(Z\in\mathcal D^{\ge0}\)，则对 \(T\in\mathcal D^{\le-1}\)，Hom 长正合列与正交性给出
\(\operatorname{Hom}(T,X)\xrightarrow{\sim}\operatorname{Hom}(T,Y)\)：相邻两项是 \(\operatorname{Hom}(T,Z[-1])\) 和 \(\operatorname{Hom}(T,Z)\)，均为零。因此 \(\tau^{\le-1}X\to\tau^{\le-1}Y\) 是同构。令 \(L=\tau^{\le-1}X\)，把复合 \(L\to X\to Y\) 代入八面体公理。前两个锥分别为 \(\tau^{\ge0}X\) 和 \(Z\)，复合的锥由上述低次截断同构识别为 \(\tau^{\ge0}Y\)，所以得到好三角
\[
\tau^{\ge0}X\longrightarrow\tau^{\ge0}Y
\longrightarrow Z\longrightarrow(\tau^{\ge0}X)[1].
\]
其三项都在 \(\mathcal D^{\ge0}\)，且截掉负次部分不改变零次上同调。前一情形便给出 \(0\to H_t^0X\to H_t^0Y\to H_t^0Z\) 正合。

一般情形下，令 \(A=\tau^{\le0}X\)、\(B=\tau^{\ge1}X\)。对 \(A\to X\to Y\) 应用八面体，得到三角 \(A\to Y\to U\) 与 \(U\to Z\to B[1]\)。第一三角给出 \(H_t^0X\to H_t^0Y\to H_t^0U\to0\) 正合；因 \(B[1]\in\mathcal D^{\ge0}\)，第二三角给出单态射 \(H_t^0U\to H_t^0Z\)。复合后正是所需的三项正合性。连接映射及其自然性来自原三角的态射，因而各次正合列可接成所述长列。
\end{proof}

\subsection{短正合列与三角的关系}
\label{b4:subsec:1903:02}

\begin{proposition}\label{b4:prop:heart-exact-triangles}
设 \(\mathcal H\) 为三角范畴 \(\mathcal D\) 中一个 t-结构的心。心中的列 \(0\to A\xrightarrow{i}E\xrightarrow{p}B\to0\) 短正合，当且仅当存在 \(\delta:B\to A[1]\) 使 \(A\xrightarrow{i}E\xrightarrow{p}B\xrightarrow{\delta}A[1]\) 为好三角。
\end{proposition}
\begin{proof}
若三角存在，\cref{b4:thm:t-cohomology-exact}的长正合列因三项都在心中而成为所述短列。反过来，将 \(i\) 完成为三角 \(A\to E\to C\)。\cref{b4:thm:heart-abelian}给出 \(H_t^{-1}(C)=\ker i=0\) 和 \(H_t^0(C)=\operatorname{coker}i=B\)。又 \(C\in\mathcal D^{[-1,0]}\)，故其截断三角直接给出 \(C\simeq B\)，并把原来的余核映射识别为 \(p\)。
\end{proof}

\begin{proposition}\label{b4:prop:heart-ext1}
设 \(\mathcal D\) 为局部小三角范畴，\(\mathcal H\) 为 \(\mathcal D\) 上一个 t-结构的心，\(A,B\in\mathcal H\)。心中的一次 Yoneda 扩张与环境中的移位态射自然对应：
\[
 \operatorname{YExt}_{\mathcal H}^1(B,A)
       \simeq\operatorname{Hom}_{\mathcal D}(B,A[1]).
\]
左边采用短正合列的等价类与 Baer 和；若心有足够内射或投射对象，它就是通常的一次 Ext。
\end{proposition}
\begin{proof}
短正合列对应前一命题中的连接映射。先证连接映射由短列唯一确定。若同一短列有两个好三角结构，其连接映射为 \(\delta_1,\delta_2:B\to A[1]\)，则对前两项取 \(1_A,1_E\)，(TR3) 给出第三项 \(c:B\to B\)，满足
\[
cp=p,\qquad \delta_2c=\delta_1.
\]
\(c\) 是心中的态射，而 \(p\) 是心中的满态射，所以第一式迫使 \(c=1_B\)，第二式便给出 \(\delta_1=\delta_2\)。若两个短列等价，对 \(A\) 取恒等、对中间项取该等价同构，仍用 (TR3) 和源短列的满态射 \(p\) 得到第三项为 \(1_B\)。因此连接映射在扩张等价类上良定义。

反之，给定 \(\delta:B\to A[1]\)，将它完成为好三角并逆旋转两次，得到 \(A\to E\to B\xrightarrow{\delta}A[1]\)。心对扩张封闭，故 \(E\in\mathcal H\)，前一命题给出短正合列。若两个三角的连接映射同为 \(\delta\)，将它们各旋转两次，写成
\[
B\xrightarrow{\delta}A[1]\xrightarrow{-i_j[1]}E_j[1]
\xrightarrow{-p_j[1]}B[1],\qquad j=1,2.
\]
对前两项取 \(1_B,1_{A[1]}\)，(TR3) 补出 \(h:E_1[1]\to E_2[1]\)。由\cref{b4:lem:pretriangle-hom}，三角态射的两个分量同构就使第三个分量同构；该结论正是对 Hom 长正合列应用五引理所得。因此 \(h\) 为同构。把三角态射旋转回来，得到 \(h[-1]:E_1\to E_2\)，满足
\[
h[-1]i_1=i_2,\qquad p_2h[-1]=p_1.
\]
它在两端 \(A,B\) 上保持恒等，故两个扩张等价。两种操作由此互逆。

对端点的态射取推出或拉回，三角态射的最后方块表明连接映射相应地作后复合或前复合。直和扩张对应连接映射的直和。Baer 和先沿 \(B\to B\oplus B\) 拉回，再沿 \(A\oplus A\to A\) 推出，因此对应于两个连接映射相加。这也证明群结构及对两端的自然性。
\end{proof}

这里把扩张送到三角的连接箭头 \(\delta\)，与\secref{b4:sec:1704}经内射消解识别一次 Ext 的约定一致。用其他符号约定书写锥时，仍须同时核对连接箭头与消解中的余圈代表。

\ProseParagraphBreak
高阶扩张可以把一次扩张拼接成链，再把连接态射移位后复合，从而得到到 \(\operatorname{Hom}_{\mathcal D}(B,A[n])\) 的比较映射；它在扩张态射下不变，因为每个方块的连接态射自然。一次的双射不能直接推广为每个 \(n\) 的双射。\secref{b4:sec:1905}将给出二次比较失效的具体例子。

\subsection{标准心与原交换范畴}
\label{b4:subsec:1903:03}

\begin{proposition}\label{b4:prop:standard-heart-equivalence}
设 \(\mathcal A\) 为交换范畴。把对象视为集中于零次的复形，得到 \(\mathcal A\) 到 \(D(\mathcal A)\) 标准心的正合等价。
\end{proposition}
\begin{proof}
若 \(X\) 只有 \(H^0\) 可能非零，截断给出拟同构屋顶 \(X\leftarrow\tau^{\le0}X\to H^0(X)[0]\)，所以该函子本质满。对 \(A,B\in\mathcal A\)，用屋顶 \(A\leftarrow Z\to B\) 表示导出态射；先以 \(\tau^{\le0}Z\) 细化，它到 \(A,B\) 的映射都经 \(H^0(Z)[0]\) 分解，且 \(H^0(Z)\to A\) 是同构。因此态射唯一来自 \(A\to B\)。唯一性也可直接对该屋顶取 \(H^0\) 看出。这证明全忠实。最后，原范畴的短正合列给出好三角，再由\cref{b4:prop:heart-exact-triangles}成为心中短正合列，所以等价正合。
\end{proof}

\begin{proposition}\label{b4:prop:nondegenerate-cohomology-conservative}
设 \(\mathcal D\) 为带非退化 t-结构的三角范畴，\(H_t^n\) 为该 t-结构的上同调函子。若对象 \(X\in\mathcal D\) 满足 \(H_t^n(X)=0\) 对所有整数 \(n\) 成立，则 \(X=0\)；若态射 \(f:X\to Y\) 在所有 \(H_t^n\) 上诱导同构，则 \(f\) 是同构。
\end{proposition}
\begin{proof}
先设 \(X\in\mathcal D^{\le0}\)。相邻截断三角和各次上同调的消失依次使 \(X\in\mathcal D^{\le-1},\mathcal D^{\le-2},\ldots\)，故非退化性给出 \(X=0\)。对 \(X\in\mathcal D^{\ge1}\) 对偶成立。一般 \(X\) 的 \(\tau^{\le0}X\) 与 \(\tau^{\ge1}X\) 仍有所有上同调为零，分别用前两种情形即可。若 \(f\) 诱导全部上同调同构，其锥由长正合列知全部上同调为零，因而锥为零，\(f\) 为同构。
\end{proof}

“检测同构”不等于“用所有上同调恢复每个态射”。例如非零的 \(\operatorname{Ext}_{\mathbb Z}^1(\mathbb Z/2,\mathbb Z)\) 给出非零导出态射 \(\mathbb Z/2\to\mathbb Z[1]\)，但它在每一个相同次数的普通上同调之间都诱导零映射。

% !TeX root = ../../Book4.tex
% 纲要标题：### 20.4 扭对与 Happel–Reiten–Smalø 倾斜
% 纲要标签：b4:sec:1904
% 纲要位置：计划纲要.md，第 2999 行。

\section{扭对与 Happel–Reiten–Smalø 倾斜}
\label{b4:sec:1904}

把心中的一类对象移到相邻次数，需要与每个对象的分解相容。扭对提供这种分解，使倾斜后的对象仍组成具有核、余核和扩张的交换范畴。


% 纲要内容（仅作写作计划，尚非教材正文）：
%
% * 在交换范畴 \(\mathcal A\) 中定义扭对 \((\mathcal T,\mathcal F)\)：\(\operatorname{Hom}(\mathcal T,\mathcal F)=0\)，每个对象有扭部分与无扭商组成的短正合列
% * 在 \(D^b(\mathcal A)\) 中由标准截断和扭分解构造倾斜后的两个子范畴，逐项证明正交性与截断三角公理
% * 固定本书方向：新心由仅在 \(-1,0\) 次有上同调、且 \(H^{-1}\in\mathcal F\)、\(H^0\in\mathcal T\) 的对象组成，即 \(\langle\mathcal F[1],\mathcal T\rangle\)
% * 利用第20.3节已证的心的交换范畴结构解释新心的核、余核和扩张；区分扭对倾斜与第21章倾斜对象的不同数据
% * 新的心不自动给出与原交换范畴的导出等价；须分别核对导出等价所需条件
%

\subsection{扭对与扭分解}
\label{b4:subsec:1904:01}

标准心把每个模或交换范畴对象放在零次。若某一类对象适合放在相邻次数，问题便成为：怎样移动这一类而仍保留核、余核和短正合列？扭分解提供了与任意对象都相容的分类方法。

\begin{definition}\label{b4:def:torsion-pair}
设 \(\mathcal A\) 为交换范畴，\(\mathcal T,\mathcal F\) 为在同构下封闭的全子范畴。若 \(\operatorname{Hom}_{\mathcal A}(T,F)=0\) 对 \(T\in\mathcal T,F\in\mathcal F\) 成立，且每个 \(M\in\mathcal A\) 都有短正合列
\[
 0\longrightarrow tM\longrightarrow M\longrightarrow fM\longrightarrow0,
 \qquad tM\in\mathcal T,\quad fM\in\mathcal F,
\]
则称 \((\mathcal T,\mathcal F)\) 为 \(\mathcal A\) 中的\term[niudui]{扭对}[torsion pair]，此列称为扭分解。这里的扭与无扭是对这对子范畴的命名，不默认它们由某个环的非零元是否零化元素来定义。
\end{definition}

\begin{proposition}\label{b4:prop:torsion-pair-properties}
设 \(\mathcal A\) 为交换范畴，\((\mathcal T,\mathcal F)\) 为 \(\mathcal A\) 中的扭对。每个对象 \(M\in\mathcal A\) 的扭分解唯一到保持 \(M\) 的唯一同构，并且对 \(M\) 函子性变化；\(\mathcal T={}^\perp\mathcal F\)、\(\mathcal F=\mathcal T^\perp\)，正交均在 \(\mathcal A\) 中取 Hom。此外，\(\mathcal T\) 对商及扩张封闭，\(\mathcal F\) 对子对象及扩张封闭。
\end{proposition}
\begin{proof}
对任意 \(T\in\mathcal T\)，映射 \(T\to M\to fM\) 为零，所以唯一经 \(tM\) 分解；这使 \(tM\to M\) 成为所有此类映射的泛对象。对 \(M\to F\)、\(F\in\mathcal F\)，映射在 \(tM\) 上为零，所以唯一经 \(fM\) 分解。这证明唯一性与函子性。若 \(M\) 与所有 \(F\) 正交，则其满射 \(M\to fM\) 为零，故 \(fM=0\) 且 \(M\in\mathcal T\)；另一侧对偶。商封闭性由 \(\operatorname{Hom}(Q,F)\to\operatorname{Hom}(T,F)\) 单射得到，子对象封闭性对偶。若两个端点都属于同一正交类，短正合列的 Hom 正合性使中间项也属于该类，故有扩张封闭性。
\end{proof}

交换群中的挠群与无挠群给出一个扭对：\(tM\) 是所有有限阶元素组成的子群，\(M/tM\) 无挠，因为 \(n\bar x=0\) 意味着 \(nx\) 有限阶，进而 \(x\) 有限阶。下一节的扭对则由线性映射的两个端点定义，与整数挠元没有关系。

\subsection{扭对倾斜的构造与公理}
\label{b4:subsec:1904:02}

\begin{theorem}\label{b4:thm:hrs-tilt}
设 \(\mathcal A\) 为交换范畴，\((\mathcal T,\mathcal F)\) 是其中的扭对。在 \(D^b(\mathcal A)\) 中令
\begin{align*}
 \mathcal D_t^{\le0}&=\{X:H^i(X)=0\ (i>0),\ H^0(X)\in\mathcal T\},\\
 \mathcal D_t^{\ge0}&=\{X:H^i(X)=0\ (i<-1),\ H^{-1}(X)\in\mathcal F\}.
\end{align*}
则 \((\mathcal D_t^{\le0},\mathcal D_t^{\ge0})\) 是有界 t-结构。这个构造称为\term[hrsqingxie]{Happel--Reiten--Smalø 倾斜}[Happel--Reiten--Smalø tilt]，简称 HRS 倾斜\cite[Chapter I, \S2]{HappelReitenSmalo1996Tilting}。\footnote{哈佩尔（Dieter Happel，1953-2012），德国数学家，研究代数表示论与导出范畴；赖滕（Idun Reiten，1942-2025）和斯莫勒（Sverre Olaf Smalø，1951-）为挪威数学家，研究有限维代数的表示与倾斜理论。}
\end{theorem}
\begin{proof}
先检查移位条件。若 \(X\in\mathcal D_t^{\le0}\)，则 \(X[1]\) 的正次上同调与零次上同调均为零，所以仍在此类；若 \(Y\in\mathcal D_t^{\ge0}\)，则 \(Y[-1]\) 在低于零次的上同调为零，故仍在此类。

新结构的 \(\mathcal D_t^{\ge1}=\mathcal D_t^{\ge0}[-1]\) 由满足 \(H^i(Y)=0\ (i<0)\)、\(H^0(Y)\in\mathcal F\) 的对象组成。若 \(X\in\mathcal D_t^{\le0}\)、\(Y\in\mathcal D_t^{\ge1}\)，标准截断的伴随首先把任意 \(X\to Y\) 唯一分解为 \(X\to H^0(X)\to Y\)，再把后一个映射唯一分解为 \(H^0(X)\to H^0(Y)\to Y\)。中间映射从 \(\mathcal T\) 到 \(\mathcal F\)，所以为零。这证明新正交性。

为构造新截断，给定 \(X\)，记 \(A=\tau^{\le0}X\)、\(T=tH^0(X)\)、\(F=fH^0(X)\)。取复合 \(A\to H^0(X)\to F\)，完成为三角
\[
 U\longrightarrow A\longrightarrow F\longrightarrow U[1].
\]
普通上同调长正合列表明 \(H^i(U)=H^i(X)\) 当 \(i<0\)，\(H^0(U)=T\)，\(H^i(U)=0\) 当 \(i>0\)；这里用到了 \(H^0(X)\to F\) 是满射。故 \(U\in\mathcal D_t^{\le0}\)。把 \(U\to A\to X\) 的复合完成为三角 \(U\to X\to V\)，八面体公理另给出三角
\[
 F\longrightarrow V\longrightarrow\tau^{\ge1}X\longrightarrow F[1].
\]
由它得 \(H^i(V)=0\) 当 \(i<0\)，且 \(H^0(V)=F\)，所以 \(V\in\mathcal D_t^{\ge1}\)。分解公理因此成立。

最后有包含关系 \(D^{\le-1}\subseteq\mathcal D_t^{\le0}\subseteq D^{\le0}\) 及 \(D^{\ge0}\subseteq\mathcal D_t^{\ge0}\subseteq D^{\ge-1}\)。标准 t-结构在有界导出范畴中有界，故新结构也有界。
\end{proof}

\subsection{倾斜方向与新心的上同调条件}
\label{b4:subsec:1904:03}

\begin{corollary}\label{b4:cor:hrs-heart-description}
设 \(\mathcal A\) 为交换范畴，\((\mathcal T,\mathcal F)\) 为扭对，\(\mathcal B\subseteq D^b(\mathcal A)\) 为\cref{b4:thm:hrs-tilt}给出的新心。则 \(\mathcal B\) 由恰好满足
\[
 H^i(E)=0\quad(i\ne-1,0),\qquad
 H^{-1}(E)\in\mathcal F,\quad H^0(E)\in\mathcal T
\]
的对象组成。每个 \(E\in\mathcal B\) 都有新心中的短正合列
\[
 0\longrightarrow H^{-1}(E)[1]\longrightarrow E
 \longrightarrow H^0(E)\longrightarrow0.
\]
因此常记 \(\mathcal B=\langle\mathcal F[1],\mathcal T\rangle\)，尖括号在这里表示按上述次序取扩张。
\end{corollary}
\begin{proof}
两个新子范畴的交给出列出的上同调条件。标准截断三角为 \(H^{-1}(E)[1]\to E\to H^0(E)\to H^{-1}(E)[2]\)，三项都在新心中，由\cref{b4:prop:heart-exact-triangles}即成短正合列。反之，以 \(F[1]\)、\(T\) 为两端的任何三角扩张都在新心中，因为心对扩张封闭。
\end{proof}

这里移动的是 \(\mathcal F\)：它的对象由零次移到 \(-1\) 次。文献也常采用反向约定，得到 \(\langle\mathcal F,\mathcal T[-1]\rangle=\mathcal B[-1]\)。比较结论时须先换算方向，不能只看“倾斜心”这一名称。

\subsection{新心的核、余核与扩张}
\label{b4:subsec:1904:04}

新心已经是交换范畴；它的核与余核由\cref{b4:thm:heart-abelian}计算，但那里要采用新的 \(H_t^i\)。普通上同调 \(H^i\) 只用来说明新心的对象是什么，并不直接给出新心中每个态射的核与余核。

\begin{example}\label{b4:exm:hrs-integer-multiplication}
取 \(\mathcal A=\mathbb Z\text{-}\mathrm{Mod}\)，扭对为挠群与无挠群，\(p\) 为素数。对象 \(\mathbb Z[1]\) 位于新心。普通模上的单射 \(p:\mathbb Z\to\mathbb Z\)，在新心中成为满态射
\[
 p:\mathbb Z[1]\longrightarrow\mathbb Z[1],\qquad
 \ker_{\mathcal B}p=\mathbb Z/p,
\]
并给出短正合列
\[
 0\longrightarrow\mathbb Z/p\longrightarrow\mathbb Z[1]
 \xrightarrow{p}\mathbb Z[1]\longrightarrow0.
\]
事实上，标准短正合列 \(0\to\mathbb Z\xrightarrow{p}\mathbb Z\to\mathbb Z/p\to0\) 的三角旋转后，前三项就是这三个新心对象；必要时把一个结构映射改乘 \(-1\)，即与旋转符号一致。由\cref{b4:prop:heart-exact-triangles}得到所述短列。
\end{example}

\((\mathcal F[1],\mathcal T)\) 本身还是新心中的扭对。其分解由\cref{b4:cor:hrs-heart-description}给出，正交性来自 \(\operatorname{Hom}(F[1],T)=\operatorname{Hom}(F,T[-1])=0\)。这一事实说明新心并非没有原扭分解的痕迹；变化的是两类对象所在的次数，以及它们之间的正合列。

\subsection{倾斜心为什么不自动给出导出等价}
\label{b4:subsec:1904:05}

HRS 倾斜的数据是一整个交换范畴中的两类对象及其扭分解。\chapref{b4:ch:20}讨论的倾斜对象则是一项对象及其自扩张消失、生成等条件。两种构造有关联，却不能仅凭名称相同便认作同一个定义。

\ProseParagraphBreak
尤其是，在 \(D^b(\mathcal A)\) 中得到新心 \(\mathcal B\)，并不已经得到 \(D^b(\mathcal B)\simeq D^b(\mathcal A)\)。若存在与心的包含相容的三角等价，它对 \(X,Y\in\mathcal B\) 必须给出每个次数的同构 \(\operatorname{Hom}_{D^b(\mathcal B)}(X,Y[n])\simeq\operatorname{Hom}_{D^b(\mathcal A)}(X,Y[n])\)。当新心有足够投射或内射对象时，左边可用其中的消解计算为 Ext。一次扩张与环境态射的对应已由\cref{b4:prop:heart-ext1}保证，高阶的一致性却需要另作证明。

\ProseParagraphBreak
下一节直接计算一个两顶点范畴：新心为半单范畴，故心内二次 Ext 为零，而环境的某个二次移位 Hom 为 \(k\)。这将排除任何与该包含相容的导出等价。一般倾斜等价还需要额外条件。


\begin{chaptersummary}
DG 范畴把复形之间的态射组织成复形，零次余圈是复形映射，零次余边是零伦映射。微分与复合的带符号相容性使 \(H^0\) 成为普通范畴；要得到三角结构，还须有相容的移位与锥。增强记录的是这些更细的数据，而不只是为原来的三角范畴另换名称。

\ProseParagraphBreak
完美复形具有有界有限生成投射模型，既能通过有限次构造从正则模得到，又能用与任意直和交换的紧性刻画。倾斜对象进一步要求全部非零移位的自身态射消失。三顶点计算把这两个条件落实为一个短消解及一张自同态矩阵；导出 Morita 定理则把这样的计算转化为导出等价。紧对象刻画用望远镜与有限因子分解提取有限模型，Morita 等价经 DG 自同态代数和双模替换构造。锥映射中的同伦选择、导出张量中的额外同调，说明进一步研究这些等价时为何需要保留增强及其相容数据。
\end{chaptersummary}

\BookExercises

\begin{optionalexercise}\label{b4:ex:20-koszul-dependent-relations}
在 \(\mathbb Z\) 上取 DG 外代数 \(A=\bigwedge(\varepsilon_1,\varepsilon_2)\)，两生成元次数均为 \(-1\)，微分为 \(\mathrm{d}\varepsilon_1=2\)、\(\mathrm{d}\varepsilon_2=4\)。计算全部上同调，并解释为何第二个关系没有产生一个正则序列的消解。
\end{optionalexercise}
\begin{solution}
底层复形为 \(\mathbb Z\to\mathbb Z^2\to\mathbb Z\)，次数为 \(-2,-1,0\)，微分分别为 \(t\mapsto(-4t,2t)\) 和 \((a,b)\mapsto2a+4b\)。第一映射单射；第二映射的核由 \((-2,1)\) 生成，而第一映射的像由其两倍生成。故 \(H^{-2}=0\)、\(H^{-1}=\mathbb Z/2\mathbb Z\)、\(H^0=\mathbb Z/2\mathbb Z\)，其他次数为零。模去第一个关系 \(2\) 后，第二个元素 \(4\) 为零，不是非零商环上的非零因子，所以 Koszul 复形留下负次同调。
\end{solution}

\begin{optionalexercise}\label{b4:ex:20-unit-boundary-contractibility}
设 \(A\) 为交换环 \(k\) 上的 DG 代数，存在 \(h\in A^{-1}\) 使 \(\mathrm{d}h=1\)。设左 DG \(A\)-模 \(M\) 的作用满足 \(\mathrm{d}(am)=\mathrm{d}a\,m+(-1)^{|a|}a\,\mathrm{d}m\)。证明 \(M\) 的底层 \(k\)-复形可缩。这个缩同伦一定是 \(A\)-线性映射吗？
\end{optionalexercise}
\begin{solution}
令 \(s(m)=hm\)，它是次数 \(-1\) 的 \(k\)-线性映射。由作用的微分公式，\(\mathrm{d}(hm)+h\,\mathrm{d}m=(\mathrm{d}h)m=m\)，所以 \(\mathrm{d}s+s\mathrm{d}=1\)。一般不能声称它为齐次 \(A\)-线性映射：次数 \(-1\) 的左模态射应满足 \(s(am)=(-1)^{|a|}a s(m)\)，这要求 \(ha=(-1)^{|a|}ah\)，而一般 DG 代数并不分次交换。因此底层可缩不等于已经给出 DG 模意义的缩同伦。
\end{solution}

\begin{optionalexercise}\label{b4:ex:20-endomorphism-two-term}
设整数 \(m\ge2\)，\(P=(\mathbb Z\xrightarrow{m}\mathbb Z)\) 位于次数 \(-1,0\)。计算 DG 代数 \(E=\underline{\operatorname{End}}_{\mathbb Z}(P)\) 的全部上同调与正次数乘法，说明为何完美对象未必满足倾斜的消失条件。
\end{optionalexercise}
\begin{solution}
只有 \(E^{-1}=\mathbb Z\)、\(E^0=\mathbb Z^2\)、\(E^1=\mathbb Z\) 非零。把零次元素写成分别作用在 \(P^{-1},P^0\) 上的 \((a,b)\)，微分为 \(h\mapsto(mh,mh)\)、\((a,b)\mapsto m(a-b)\)。故 \(H^{-1}=0\)、\(H^0=\mathbb Z/m\mathbb Z\)、\(H^1=\mathbb Z/m\mathbb Z\)，其余为零。零次类按通常标量作用于一次类，一次类乘积因二次项为零而为零。虽然 \(P\) 有界有限自由，其到一次移位的导出态射仍非零。
\end{solution}

\begin{optionalexercise}\label{b4:ex:20-hzero-not-quasi-equivalence}
设 \(k\) 为域，把 \(k\) 与 \(k[u]\)、\(|u|=2\)、\(\mathrm{d}=0\) 分别视为只有一个对象的 DG 范畴。证明包含映射在 \(H^0\) 上给出等价，但不是拟等价。说明这个例子是否已经是两个三角增强之间的反例。
\end{optionalexercise}
\begin{solution}
两者的零次上同调代数均为 \(k\)，对象集合也都是单点，所以 \(H^0\) 上为等价。但第二个态射复形的二次上同调含有非零类 \(u\)，第一个没有，因此 Hom 复形映射不是拟同构。两个单对象 DG 范畴没有包含所需的移位和锥，不是这里的预三角增强；此例只说明任意 DG 范畴的 \(H^0\) 丢失了资料，不能据此推出增强间的更强结论。
\end{solution}

\begin{optionalexercise}\label{b4:ex:20-perfect-nonprojective-summand}
设 \(R=\mathbb Z\)，\(P=(R\xrightarrow{6}R)\) 位于次数 \(-1,0\)。利用中国剩余定理，在导出范畴中把 \(P\) 分成两个非零完美直和项，并为每一项写出有限自由模型。
\end{optionalexercise}
\begin{solution}
\(P\) 拟同构于零次的 \(\mathbb Z/6\mathbb Z\)，中国剩余定理给出 \(\mathbb Z/6\mathbb Z\cong\mathbb Z/2\mathbb Z\oplus\mathbb Z/3\mathbb Z\)。两项的模型分别为 \((\mathbb Z\xrightarrow{2}\mathbb Z)\) 和 \((\mathbb Z\xrightarrow{3}\mathbb Z)\)，均置于次数 \(-1,0\)。由这些拟同构在导出范畴中可逆，得到所需分解。不能把这个分解理解成原来每项 \(\mathbb Z\) 都分为两个非零直和项。
\end{solution}

\begin{optionalexercise}\label{b4:ex:20-field-perfect}
设 \(k\) 为域，\(X\) 为任意 \(k\)-向量空间上链复形。证明 \(X\) 同伦等价于零微分复形 \(\bigoplus_n H^n(X)[-n]\)。再证明 \(X\) 完美当且仅当 \(\sum_n\dim_kH^n(X)<\infty\)。
\end{optionalexercise}
\begin{solution}
在每次取补空间 \(Z^n=B^n\oplus H_n\)、\(X^n=Z^n\oplus L_n\)。微分限制给出同构 \(L_n\to B^{n+1}\)：它的核为零，像为全部边界。于是 \(X\) 分为零微分的 \(H_n\) 与各两项恒等复形的直和；后一部分用上述同构的逆得到缩同伦。每个次数至多涉及相邻的两个可缩项，所以这个定义适用于无界复形。若总同调维数有限，则零微分部分为有界有限维复形，因而完美。反之，有界有限维模型的同调只能在有限个次数非零，且每次有限维，给出有限总维数。
\end{solution}

\begin{optionalexercise}\label{b4:ex:20-thick-product-ring}
设 \(k\) 为域，\(R=k\times k\)，\(e=(1,0)\)。在左 \(R\)-模导出范畴中，准确描述 \(\operatorname{thick}(Re)\) 中的对象，并证明 \(Re\) 不是 \(\operatorname{Perf}(R)\) 的倾斜生成元。
\end{optionalexercise}
\begin{solution}
每个 \(R\)-模及复形都由两个幂等元分解为一对 \(k\)-向量空间模及复形；映射、锥和同调均按分量计算。\(Re\) 对应 \((k,0)\)。因此它的厚子范畴所有对象第二分量为零，第一分量完美。反过来第一分量的每个完美 \(k\)-复形由有限直和及移位的 \(k\) 构成，故这些对象全在其中。\(R(1-e)=(0,k)\) 不属于该子范畴，因此生成条件失败，尽管 \(Re\) 自身的非零移位态射全为零。
\end{solution}

\begin{optionalexercise}\label{b4:ex:20-tilting-shifts-over-field}
设 \(k\) 为域，\(T\in\operatorname{Perf}(k)\) 非零。证明 \(T\) 是倾斜复形当且仅当存在整数 \(a\) 和正整数 \(r\)，使 \(T\cong k^r[a]\)。计算其反自同态代数。
\end{optionalexercise}
\begin{solution}
由向量空间复形的分裂，\(T\) 是各次同调的有限直和。若两个不同次数 \(p<q\) 的同调非零，则从 \(H^p(T)[-p]\) 到 \(H^q(T)[-q][q-p]\) 存在非零线性映射，得到 \(T\to T[q-p]\) 的非零态射，违背消失条件。故同调只能集中在一次，此时为 \(k^r[a]\)。反之这种对象没有非零移位的自身态射，而 \(k[a]\) 是其直和项，故厚生成所有完美复形。自同态环为 \(M_r(k)\)，取反后仍借矩阵转置同构于 \(M_r(k)\)。
\end{solution}

\begin{optionalexercise}\label{b4:ex:20-a3-generation-sequence}
沿用\cref{b4:exm:a3-tilting-object}的记号，构造一个符合\cref{b4:prop:tilting-module-criterion}的正合列 \(0\to R\to T_0\to T_1\to0\)。
\end{optionalexercise}
\begin{solution}
取 \(T_0=P_1\oplus P_2\oplus P_2\)、\(T_1=S_2\)。从 \(R=P_1\oplus P_2\oplus P_3\) 到 \(T_0\) 的映射为前两分量恒等、第三分量为 \(i:P_3\to P_2\)；从 \(T_0\) 到 \(S_2\) 取第三分量的商映射 \(q\)。核、像关系由\eqref{b4:eq:a3-simple-resolution}得到。\(T_0,T_1\) 均是有限个 \(T\) 的直和项，因此这把正文的厚生成计算写成该判据要求的正合列。
\end{solution}

\begin{optionalexercise}\label{b4:ex:20-a3-replace-wrong-simple}
在 \(R=k(1\to2\to3)\) 上令 \(S_1=(k\to0\to0)\)，考虑 \(U=P_1\oplus P_2\oplus S_1\)。证明 \(U\) 不满足倾斜消失条件，指出一个非零 Ext 群。
\end{optionalexercise}
\begin{solution}
有 \(0\to P_2\to P_1\to S_1\to0\)。施加 \(\operatorname{Hom}_R(-,P_2)\)，得到 \(\operatorname{Ext}^1_R(S_1,P_2)=\operatorname{coker}(\operatorname{Hom}(P_1,P_2)\to\operatorname{End}(P_2))\)。前一空间为 \((P_2)_1=0\)，后一空间为 \(k\)，所以该 Ext 群为 \(k\)。它是 \(\operatorname{Ext}^1_R(U,U)\) 的一个直和项，因此消失条件失败。
\end{solution}

\begin{optionalexercise}\label{b4:ex:20-formality-is-zigzag}
设 \(k\) 为域，\(E\) 为 DG 代数，\(H^n(E)=0\) 对 \(n\ne0\) 成立。说明“为每个零次上同调类任取一个代表”为何不自动构成代数映射 \(H^0(E)\to E\)。构造一个保持单位的 \(k\)-线性代表选择，要求它仍然破坏乘法。
\end{optionalexercise}
\begin{solution}
代表的选择首先只是一条到零次余圈的集合截面；即使取保持单位的 \(k\)-线性截面，也没有保证 \(s(ab)=s(a)s(b)\)。取交换 DG 代数 \(E=k[z,t]\otimes_k\bigwedge(\varepsilon)\)，\(|z|=|t|=0\)、\(|\varepsilon|=-1\)、\(\mathrm{d}\varepsilon=t\)、\(\mathrm{d}z=\mathrm{d}t=0\)。乘 \(t\) 在 \(k[z,t]\) 上单射，所以其上同调仅为零次的 \(k[z]\)。令 \(c_2(f)\) 为多项式 \(f(z)\) 的二次项系数，取 \(s(f)=f(z)+c_2(f)t\)。它保持单位且 \(k\)-线性，模去 \(t\) 后为恒等，确为代表选择；但 \(s(z)^2=z^2\ne z^2+t=s(z^2)\)。正文的截断拟同构避免了依赖这种不相容的选择。
\end{solution}

\begin{optionalexercise}\label{b4:ex:20-cone-homotopy-change}
沿用\eqref{b4:eq:cone-map-with-homotopy}的记号。若两种方块同伦 \(h,h'\) 满足 \(h-h'=\mathrm{d}_{D'}s-s\mathrm{d}_C\)，其中 \(s:C\to D'\) 的次数为 \(-2\)，证明它们给出的两个锥映射同伦。
\end{optionalexercise}
\begin{solution}
在锥上取次数 \(-1\) 的映射 \(S(y,x)=(sx,0)\)。锥微分给出
\[
 \mathrm{d}_{\operatorname{Cone}(f')}S+S \mathrm{d}_{\operatorname{Cone}(f)}
 =\begin{pmatrix}0&\mathrm{d}_{D'}s-s\mathrm{d}_C\\0&0\end{pmatrix}.
\]
这恰为两个锥映射的差。因此二次同伦控制一次同伦选择所导致的锥映射差别。
\end{solution}

\begin{optionalexercise}\label{b4:ex:20-derived-intersection-transverse}
设 \(k\) 为域，\(A=k[x,y]\)，比较 \(A/(x)\otimes_A^{\mathbf L}A/(y)\) 与 \(A/(x)\otimes_A^{\mathbf L}A/(x)\) 的上同调。说明额外同调由哪个乘法映射的核产生。
\end{optionalexercise}
\begin{solution}
用 \((A\xrightarrow{x}A)\) 消解 \(A/(x)\)。与 \(A/(y)=k[x]\) 张量后，微分仍为单射乘 \(x\)，故仅零次同调 \(k\) 非零。与 \(A/(x)=k[y]\) 张量后，微分为零，故零次与负一次同调均为 \(k[y]\)。额外负一次同调来自乘 \(x\) 在另一个商环中出现的核：前一商环中核为零，后一商环中为整个商环。
\end{solution}

% !TeX root = ../../Book4.tex
% 纲要标题：## 第21章*　DG 范畴与导出等价
% 纲要标签：b4:ch:20
% 纲要位置：计划纲要.md，第 3024 行。

\OptionalChapter{DG 范畴与导出等价}{b4:ch:20}
\BookChapterNumber{b4:ch:20}

\begin{chapterabstract}
本章用 Hom 复形建立 DG 范畴及其同伦信息，研究完美复形与倾斜对象，并以三顶点路径代数计算一个导出等价而非普通 Morita 等价的例子。
\end{chapterabstract}

\begin{learninggoals}
\begin{itemize}
\item 正确展开 DG 乘法、Hom 微分与复合中的符号，区分 DG 微分与普通指定导子。
\item 说明 \(H^0\) 忘掉了哪些资料，并在有界投射复形中识别一个具体增强。
\item 区分完美性、同调有界、每项有限生成和紧性，核对厚生成的有限构造。
\item 在三个顶点的表示上计算倾斜条件、自同态代数以及缺失投射模的恢复。
\item 说明紧对象刻画中怎样提取有限模型，以及导出 Morita 等价中双模替换的作用。
\end{itemize}
\end{learninggoals}

两个复形之间的链映射可以逐项相加，也可以作复合；但若只记录它们的同伦类，就忘记了产生同伦的具体映射。把各次齐次映射一起放入 Hom 复形，微分的零次闭元正是链映射，零次边界则给出零同伦的链映射。取 \(H^0\) 便回到熟悉的同伦范畴；保留整个复形，就还能比较高次 Ext 所在的次数。DG 范畴要保存的正是这份比普通态射集合更丰富的数据。它如何帮助选择导出模型、怎样产生导出等价，本章将从完美复形与具体路径代数算起。

\ProseParagraphBreak
本章的基本计算使用\cref{b4:thm:bounded-projective-model,b4:thm:derived-ext-identification}及\chapref{b4:ch:18}的导出 Hom、张量积；判断三顶点例子不构成普通 Morita 等价时，还调用第二卷定理15.3.10的不可分解投射模分类。不同心的比较参照\chapref{b4:ch:19}。三顶点例子的倾斜判定使用\subsecref{b4:subsec:2002:03}证明的完美复形厚性。该小节对紧对象的反向刻画还依赖\secref{b4:sec:1804}的半自由消解。

\ProseParagraphBreak
Keller 的《Deriving DG categories》\cite[Sections 1.1--2.2]{Keller1994DerivingDG}给出分次态射、DG 模和同伦范畴的约定；紧对象及导出等价分别见\cite[Section 5.3; Sections 8.2 and 9.1--9.2]{Keller1994DerivingDG}。普通环的完美复形和 Rickard 定理也可查其《Introduction to abelian and derived categories》\cite[Sections 3.1 and 4]{Keller1998AbelianDerived}。本章使用前一论文时同时采用作者的勘误：无界消解中的零调判据须保留圈复形条件，DG 复形的平坦性须按保持无界零调复形来理解。

% !TeX root = ../../Book4.tex
% 纲要标题：### 21.1 DG 代数与 DG 范畴
% 纲要标签：b4:sec:2001
% 纲要位置：计划纲要.md，第 3026 行。

\section{DG 代数与 DG 范畴}
\label{b4:sec:2001}

在适当的投射或内射模型上，导出态射由复形映射的同伦类表示。完整 Hom 复形还保留具体同伦及其复合；将这些运算放在一起，乘法与微分的符号便须满足分次 Leibniz 法则。


% 纲要内容（仅作写作计划，尚非教材正文）：
%
% * 区别第三卷附录 E 的指定导子：这里采用分次微分、带符号 Leibniz 法则及微分平方为零
% * 带微分的分次态射空间
% * 复合与符号
% * 同伦范畴及增强
%

\subsection{分次微分与带符号 Leibniz 法则}
\label{b4:subsec:2001:01}

先看一个可以完全算出的 Hom 复形。设 \(k\) 为域，\(C=(k\xrightarrow{1}k)\) 位于零次和一次。它可缩，所以在同伦范畴中为零对象。不过，次数 \(-1\) 的映射 \(h:C^1\to C^0\)、\(h=1\)，仍记录了收缩同伦。其分次自同态空间 \(E=\underline{\operatorname{End}}_k(C)\) 连同微分为
\[
 E^{-1}=k\xrightarrow{\ c\mapsto(c,c)\ }E^0=k\oplus k
 \xrightarrow{\ (a,b)\mapsto a-b\ }E^1=k.
\]
这里微分由 \(\partial f=\mathrm{d}_Cf-(-1)^{|f|}f\mathrm{d}_C\) 给出，尤其 \(\partial h=1_C\)。三处上同调都为零，但 \(E\) 还保留了使恒等映射零伦的具体 \(h\)，以及所有映射的复合。要同时保存这些数据，就须把分次乘法与微分放在同一结构中。

\begin{definition}\label{b4:def:dg-algebra}
设 \(k\) 为交换环，\(A=\bigoplus_{n\in\mathbb Z}A^n\) 为含幺的分次结合 \(k\)-代数，单位位于 \(A^0\)。设 \(\mathrm{d}:A^n\to A^{n+1}\) 为 \(k\)-线性映射，满足 \(\mathrm{d}^2=0\)，并且对 \(a\in A^p,b\in A^q\) 有
\[
 \mathrm{d}(ab)=\mathrm{d}(a)b+(-1)^p a\,\mathrm{d}(b).
\]
称 \((A,\mathrm{d})\) 为\term[weifenfencidaishu]{微分分次代数}[differential graded algebra]，简称 DG 代数。其态射保持分次、乘法、单位与微分。
\end{definition}

由 \(\mathrm{d}(1)=\mathrm{d}(1\cdot1)=2\mathrm{d}(1)\) 得 \(\mathrm{d}(1)=0\)，这里不需要除以 \(2\)。符号取决于左因子的次数：微分越过一个 \(p\) 次元素时出现 \((-1)^p\)。普通 \(k\)-代数放在零次、微分取零，便成为 DG 代数；DG 代数则不要求乘法分次交换。

\ProseParagraphBreak
第三卷附录 E 的微分代数从一个满足普通 Leibniz 法则的指定导子出发，导子不一定平方为零。本节的微分升高次数，既有符号条件，又必须平方为零。例如 \(k[x]\) 上的 \(\partial/\partial x\) 并不因名称中含有“微分”就给出这里的结构；在特征零时，它作用于 \(x^2\) 两次的结果为 \(2\)。

\begin{example}\label{b4:exm:dg-koszul-one-element}
设 \(R\) 为交换环，\(a\in R\)。令 \(A^0=R\)、\(A^{-1}=R\varepsilon\)，其余次数为零，规定 \(\varepsilon^2=0\)、\(\mathrm{d}\varepsilon=a\)、\(\mathrm{d}R=0\)。这给出一个 DG \(R\)-代数。唯一需要特别检查的关系是
\[
 \mathrm{d}(\varepsilon^2)=a\varepsilon-\varepsilon a=0.
\]
其上同调为 \(H^0(A)=R/aR\)、\(H^{-1}(A)=\operatorname{ann}_R(a)\)。当 \(a\) 不是零因子时，增广 \(A\to R/aR\) 是拟同构；当 \(a\) 为单位时，底层复形可缩。这里重现了单元素 Koszul 复形，但同时记录了外代数乘法。
\end{example}

\begin{proposition}\label{b4:prop:dg-cohomology-product}
设 \(A\) 为交换环 \(k\) 上的 DG 代数。公式 \([a][b]=[ab]\) 在 \(H^*(A)=\bigoplus_n H^n(A)\) 上给出分次结合代数结构。
\end{proposition}
\begin{proof}
若 \(\mathrm{d}a=\mathrm{d}b=0\)，则 Leibniz 法则给出 \(\mathrm{d}(ab)=0\)。若把 \(a\) 改为 \(a+\mathrm{d}u\)，其中 \(u\) 为 \(p-1\) 次元素，则乘积之差为 \((\mathrm{d}u)b=\mathrm{d}(ub)\)；若把 \(b\) 改为 \(b+\mathrm{d}v\)，则 \(a(\mathrm{d}v)=(-1)^p \mathrm{d}(av)\)。因此乘积不依赖余圈代表。结合律和单位由 \(A\) 中的对应等式下降；若 \([1]=0\)，所得上同调代数为零代数。
\end{proof}

\subsection{带微分的分次态射空间}
\label{b4:subsec:2001:02}

设 \(R\) 为含幺环，\(C,D\) 为左 \(R\)-模上链复形。一个次数为 \(r\) 的齐次映射，是一族不必与微分交换的 \(R\)-线性映射 \(f^n:C^n\to D^{n+r}\)。全部这类映射组成
\[
 \underline{\operatorname{Hom}}_R(C,D)^r
 =\prod_{n\in\mathbb Z}\operatorname{Hom}_R(C^n,D^{n+r}).
\]
积号不能随意改为直和：一个齐次映射可以在无限多个次数上非零。令
\begin{equation}\label{b4:eq:dg-hom-differential}
 (\partial f)^n=\mathrm{d}_D^{n+r}f^n-(-1)^r f^{n+1}\mathrm{d}_C^n.
\end{equation}
这正是前面总 Hom 复形的微分；此处把同一公式用于组织所有复形之间的映射。

\begin{proposition}\label{b4:prop:dg-hom-complex}
设 \(R\) 为含幺环，\(C,D\) 为左 \(R\)-模上链复形。在 \(\underline{\operatorname{Hom}}_R(C,D)\) 上，对次数 \(r\) 的齐次映射 \(f\) 定义 \(\partial f=\mathrm d_Df-(-1)^rf\mathrm d_C\)，则 \(\partial^2=0\)。其零次余圈是复形映射，零次余边是零伦映射，因而
\[
 H^0\underline{\operatorname{Hom}}_R(C,D)
 \cong\operatorname{Hom}_{K(R\text{-}\mathrm{Mod})}(C,D).
\]
\end{proposition}
\begin{proof}
省略上标而保留次数，展开两次微分得到
\[
 \partial^2 f=\mathrm{d}_D^2f-(-1)^r \mathrm{d}_Df\mathrm{d}_C
             -(-1)^{r+1}\mathrm{d}_Df\mathrm{d}_C-f\mathrm{d}_C^2=0.
\]
在零次，\(\partial f=0\) 就是 \(\mathrm{d}_Df=f\mathrm{d}_C\)。若 \(h\) 的次数为 \(-1\)，则 \(\partial h=\mathrm{d}_Dh+h\mathrm{d}_C\)，恰是零伦公式。取商即得到所述同伦范畴中的态射群。
\end{proof}

较高次数也有明确意义：\(\partial f=0\) 等价于把 \(f^n\) 看作 \(C^n\to D[r]^n\) 后得到复形映射，因为 \(\mathrm{d}_{D[r]}=(-1)^r \mathrm{d}_D\)。相应余边与同伦只差一个整体符号，故 \(H^r\underline{\operatorname{Hom}}_R(C,D)=\operatorname{Hom}_{K(R\text{-}\mathrm{Mod})}(C,D[r])\)。若 \(C\) 是有界投射复形，\cref{b4:thm:bounded-projective-model}又把这里的同伦态射识别为导出态射。任意复形则不能省去这个条件。

\subsection{复合与符号}
\label{b4:subsec:2001:03}

齐次映射仍按通常次序复合。若 \(f:C\to D\) 的次数为 \(p\)，\(g:D\to E\) 的次数为 \(q\)，则 \((gf)^n=g^{n+p}f^n\)，次数为 \(p+q\)。直接展开可得
\begin{equation}\label{b4:eq:dg-composition-leibniz}
 \partial(gf)=(\partial g)f+(-1)^q g(\partial f).
\end{equation}
右边的中间两项为 \(-(-1)^q g\mathrm{d}_Df\) 与 \(( -1)^q g\mathrm{d}_Df\)，正好抵消。因此符号由写在左边、后执行的 \(g\) 的次数决定，而不是由先执行的 \(f\) 决定。

\begin{definition}\label{b4:def:dg-category}
设 \(k\) 为交换环。给定一组数据 \(\mathcal E\)：一族对象；对每两个对象 \(X,Y\)，给定 \(k\)-模上链复形 \(\underline{\operatorname{Hom}}_{\mathcal E}(X,Y)\)，称其为态射复形；给定双线性、保持总次数的复合，以及零次余圈 \(1_X\in\underline{\operatorname{Hom}}_{\mathcal E}(X,X)^0\)。若复合结合，\(1_X\) 是单位，并且对次数 \(p,q\) 的齐次态射 \(f:X\to Y\)、\(g:Y\to Z\) 满足 \(\mathrm d(gf)=(\mathrm dg)f+(-1)^qg(\mathrm df)\)，则称这些数据构成一个\term[weifenfencifanchou]{微分分次范畴}[differential graded category]，简称 DG 范畴。

两个 \(k\)-线性 DG 范畴之间，若对象对应与各 Hom 复形之间的复形映射保持单位和复合，称其为\term[DGhanzi]{DG 函子}[DG functor]。
\end{definition}

\(R\)-模复形及上述 Hom 复形组成 DG 范畴 \(\mathrm{Ch}_{\mathrm{dg}}(R)\)。只有一个对象的 DG 范畴等价于一个 DG 代数；特别是 \(\underline{\operatorname{End}}_R(C)=\underline{\operatorname{Hom}}_R(C,C)\) 以复合作为乘法，是 DG 代数。其零次上同调才是 \(C\) 的同伦自同态环，二者一般不是同一个对象。

\ProseParagraphBreak
张量积中出现的符号有相同来源。设 \(A,B\) 为交换环 \(k\) 上的 DG 代数，对齐次元素规定
\[
 (a\otimes b)(a'\otimes b')=(-1)^{|b|\cdot|a'|}aa'\otimes bb',
 \qquad
 \mathrm{d}(a\otimes b)=\mathrm{d}a\otimes b+(-1)^{|a|}a\otimes \mathrm{d}b.
\]
在展开乘积的微分时，把 \(b\) 越过 \(a'\) 所需的符号恰好使 Leibniz 法则成立。与此相配，DG 反代数的乘法为 \(a\mathbin{\cdot_{\mathrm{op}}}b=(-1)^{|a|\cdot|b|}ba\)，而不只是反转次序。这些约定在 DG 双模和张量函子中不可省略；集中在零次时，它们还原为普通约定。

\subsection{同伦范畴及增强}
\label{b4:subsec:2001:04}

对 DG 范畴 \(\mathcal E\)，保留对象而把态射复形换成零次上同调，得到范畴 \(H^0(\mathcal E)\)，其中
\[
 \operatorname{Hom}_{H^0(\mathcal E)}(X,Y)
 =H^0\underline{\operatorname{Hom}}_{\mathcal E}(X,Y).
\]
复合良定义，是因为余圈复合仍是余圈，而且一个零次余边与零次余圈复合后仍是余边。\cref{b4:prop:dg-hom-complex}说明
\[
 H^0\bigl(\mathrm{Ch}_{\mathrm{dg}}(R)\bigr)=K(R\text{-}\mathrm{Mod}).
\]
因此取 \(H^0\) 已经忘掉了同伦的具体选择。它还没有自动把所有拟同构倒置。

\begin{definition}\label{b4:def:dg-quasi-equivalence}
设 \(F:\mathcal E\to\mathcal F\) 为交换环 \(k\) 上的 DG 函子。若每个 \(\underline{\operatorname{Hom}}_{\mathcal E}(X,Y)\to\underline{\operatorname{Hom}}_{\mathcal F}(FX,FY)\) 是拟同构，而且 \(H^0(F)\) 本质满，称 \(F\) 为\term[nidengjia]{拟等价}[quasi-equivalence]。
\end{definition}

拟等价诱导 \(H^0\) 上的等价，因为零次上同调上的映射已经双射。反过来，仅知道零次上同调范畴等价，并没有提供高次映射或与微分、复合相容的提升。

\ProseParagraphBreak
说一个三角范畴具有 DG 描述，还必须解释它的移位与好三角来自哪里。为此可以用 DG 模表达闭性条件。一个右 DG \(\mathcal E\)-模是 DG 函子 \(\mathcal E^{\mathrm{op}}\to\mathrm{Ch}_{\mathrm{dg}}(k)\)，其中反范畴的复合采用上一小节的符号。

\ProseParagraphBreak
先把这个定义翻译到只有一个对象的情形，此时 \(\mathcal E\) 就是 DG 代数 \(A\)。通常的右 DG \(A\)-模写作右作用 \(ma\)，满足
\[
 \mathrm d(ma)=(\mathrm dm)a+(-1)^{|m|}m\,\mathrm da.
\]
它对应的 DG 函子在齐次元素上采用
\[
 M(a^{\mathrm{op}})(m)=(-1)^{|a||m|}ma.
\]
这个符号与 \(a^{\mathrm{op}}b^{\mathrm{op}}=(-1)^{|a||b|}(ba)^{\mathrm{op}}\) 相配。例如 \(|a|=|b|=|m|=1\) 时，连续作用给出 \(M(a^{\mathrm{op}})M(b^{\mathrm{op}})(m)=-mba\)，与反代数中乘积的作用相同。因而右作用的通常写法与反范畴函子的写法之间，需要这一次符号转换。

\ProseParagraphBreak
对两个右 DG \(\mathcal E\)-模 \(M,N\)，次数 \(r\) 的态射是一族次数 \(r\) 的 \(k\)-线性映射 \(\eta_Y:M(Y)\to N(Y)\)，满足 \(N(f)\eta_Y=(-1)^{r|f|}\eta_{Y'}M(f)\)，其中 \(f:Y'\to Y\) 为 \(\mathcal E\) 中的任意齐次态射。在单对象情形，把上面的转换式代入，两边的次数符号抵消，便得到通常的右线性条件 \(\eta(ma)=\eta(m)a\)。态射复形的微分仍为\eqref{b4:eq:dg-hom-differential}；模的移位、映射锥逐对象构造。对齐次态射 \(f\)，移位作用为 \(M[1](f)=(-1)^{|f|}M(f)\)；对闭的零次模态射 \(u:M\to N\)，锥上的作用为 \(\operatorname{diag}(N(f),(-1)^{|f|}M(f))\)。这些符号使作用与移位及锥的微分相容。

\ProseParagraphBreak
可表模 \(h_X\) 把 \(Y\) 送到 \(\underline{\operatorname{Hom}}_{\mathcal E}(Y,X)\)，齐次 \(f:Y'\to Y\) 对 \(a:Y\to X\) 的作用为 \((-1)^{|f|\cdot|a|}af\)。在单对象情形，\(h_*=A\)；若 \(a,f\) 都为一次，函子作用是 \(h_*(f^{\mathrm{op}})(a)=-af\)。再按上面的转换式还原通常右作用，得到的正是右正则模的乘法 \(af\)。这说明 Yoneda 公式中的负号来自反范畴约定，而不是把右乘法本身改成负的乘法。

\begin{definition}\label{b4:def:dg-enhancement}
设 \(\mathcal E\) 为交换环 \(k\) 上的 DG 范畴。若其可表右 DG 模及与它们同构的对象在 DG 模的同伦范畴中组成一个包含零对象、对正负移位及零次闭态射的映射锥封闭的满子范畴，称 \(\mathcal E\) 为\term[yu-sanjiao-DGfanchou]{预三角 DG 范畴}[pretriangulated DG category]。设 \(\mathcal T\) 为 \(k\)-线性三角范畴；一对数据 \((\mathcal E,\Phi)\)，其中 \(\mathcal E\) 预三角而 \(\Phi:H^0(\mathcal E)\to\mathcal T\) 为三角等价，称为 \(\mathcal T\) 的\term[DGzengqiang]{DG 增强}[DG enhancement]。
\end{definition}

先在全部右 DG 模的同伦范畴中构造三角结构。移位与锥逐对象形成复形；上述作用符号保证它们仍是 DG 模。\cref{b4:thm:homotopy-triangulated}证明中的锥映射及同伦矩阵也与模作用相容，因此同一验证给出这里的三角公理。DG Yoneda 函子则在态射复形上给出同构
\[
\underline{\operatorname{Hom}}_{\mathrm{DGMod}(\mathcal E)}(h_X,h_Y)
\cong\underline{\operatorname{Hom}}_{\mathcal E}(X,Y).
\]
其中 \(\mathrm{DGMod}(\mathcal E)\) 表示右 DG 模的 DG 范畴；该同构把模态射送到它在 \(1_X\) 处的值，逆映射把齐次 \(g:X\to Y\) 送到逐分量的后复合 \(a\mapsto ga\)。带符号的自然性与 Leibniz 法则保证这两个对应与微分相容。取零次上同调，便得到 \(H^0(\mathcal E)\) 到 DG 模同伦范畴的全忠实函子。

\ProseParagraphBreak
若 \(\mathcal E\) 预三角，这个函子的本质像包含零对象，并对正负移位及锥封闭。它也对有限双积封闭，因为零态射 \(h_X[-1]\to h_Y\) 的锥就是 \(h_Y\oplus h_X\)。把 DG 模同伦范畴的好三角限制到该满子范畴，便得到三角结构：三角的补全、旋转及八面体中所需的对象都是这些移位和锥，因而仍在本质像中，三角态射的各分量则由满性保留。DG Yoneda 的全忠实性给出 \(H^0(\mathcal E)\) 与这个本质像的等价，由此将三角结构传回 \(H^0(\mathcal E)\)。这才是定义中 DG 增强所用的三角结构。DG 模的这些构造也可查 Keller 的《Deriving DG categories》\cite[Sections 1.2 and 2.2]{Keller1994DerivingDG}。一般 DG 范畴可能既没有移位对象，也没有锥对象，不能仅凭 DG 模同伦范畴具有三角结构，就断言其 \(H^0\) 是三角范畴。

\begin{example}\label{b4:exm:bounded-projective-enhancement}
设 \(R\) 为含幺环。取所有有界、每项有限生成投射的左 \(R\)-模复形以及它们的整个 Hom 复形，得到 \(\mathrm{Ch}^b_{\mathrm{dg}}(R\text{-}\mathrm{proj})\)。移位仍在其中；映射锥的项为两个有限生成投射模的直和，也仍在其中。因此它是预三角的，\(H^0\) 是 \(K^b(R\text{-}\mathrm{proj})\)。下一节把这个范畴识别为完美复形的导出范畴模型。这是一个直接可写出的增强，而非从抽象三角公理中反求出来的结构。
\end{example}

% !TeX root = ../../Book4.tex
% 纲要标题：### 21.2 完美复形
% 纲要标签：b4:sec:2002
% 纲要位置：计划纲要.md，第 3794 行。

\section{完美复形}
\label{b4:sec:2002}

真正有限的导出对象需要有界有限生成投射模型。厚生成描述它的有限构造，紧性则用它到任意直和的态射来识别同一类对象。前两小节只涉及有界模型与有限构造；第三小节的任意余积与紧对象反向刻画处在无界导出范畴，依赖\secref{b4:sec:1804}建立的半自由替换及\cref{b4:thm:semifree-resolution,b4:prop:k-resolution-comparison}。


% 纲要内容（仅作写作计划，尚非教材正文）：
%
% * 有限投射复形
% * 厚子范畴
% * 紧性概念的适用背景
%

\subsection{有限投射复形}
\label{b4:subsec:2002:01}

一个复形的同调可能仅在有限个次数非零，但计算这些同调所需的投射消解却可能无限延伸。要在导出范畴中表达真正有限的投射数据，应把有限性加在一个可选取的投射模型上。

\begin{definition}\label{b4:def:perfect-complex}
设 \(R\) 为含幺环，\(X\in D(R\text{-}\mathrm{Mod})\)。若 \(X\) 在该导出范畴中同构于某个有界复形 \(P\)，且每个 \(P^n\) 都是有限生成投射左 \(R\)-模，称 \(X\) 为\term[wanmeifuxing]{完美复形}[perfect complex]。全部完美复形组成的全子范畴记为 \(\operatorname{Perf}(R)\)。
\end{definition}
\symidx{\(\operatorname{Perf}(R)\)：环上的完美复形范畴}

\begin{proposition}\label{b4:prop:perfect-homotopy-model}
设 \(R\) 为含幺环，以 \(R\text{-}\mathrm{proj}\) 表示有限生成投射左 \(R\)-模的满子范畴。从同伦范畴到导出范畴的自然函子限制为三角等价
\[
 K^b(R\text{-}\mathrm{proj})\ \simeq\ \operatorname{Perf}(R).
\]
完美复形对有限直和、移位和导出范畴中的映射锥封闭。
\end{proposition}
\begin{proof}
由\cref{b4:prop:resolution-computes-derived-hom}，有界投射复形之间的导出态射就是复形映射的同伦类，所以所给函子全忠实；本质满则是完美的定义。有限直和和移位保持项的有限生成投射性。对完美对象之间的态射，先取有界有限生成投射模型，再用全忠实性把态射表示成复形映射。其映射锥依然有界，每项是两个有限生成投射模的直和，因此仍然完美。好三角来自这样的锥，故等价保持三角结构。
\end{proof}

例如，\(\mathbb Z/m\mathbb Z\) 置于零次是 \(\mathbb Z\) 上的完美复形，因为它有两项有限自由消解 \(\mathbb Z\xrightarrow{m}\mathbb Z\)。这并不要求零次同调本身投射。

\begin{example}\label{b4:exm:bounded-not-perfect}
设 \(k\) 为域，\(R=k[\varepsilon]/(\varepsilon^2)\)，\(S=R/(\varepsilon)\)。模 \(S\) 置于零次虽然有界且有限生成，却不是完美复形。它的自由消解为
\[
 \cdots\xrightarrow{\varepsilon}R\xrightarrow{\varepsilon}R
 \xrightarrow{\varepsilon}R\longrightarrow S\longrightarrow0,
\]
因为每个乘 \(\varepsilon\) 映射的核、像都是 \(\varepsilon R\)。施加 \(\operatorname{Hom}_R(-,S)\) 后微分全为零，所以 \(\operatorname{Ext}^n_R(S,S)=k\) 对所有 \(n\ge0\) 成立。若 \(S\) 有有界投射模型 \(P\)，则 \(\underline{\operatorname{Hom}}_R(P,S)\) 有界，其高次上同调终将消失；结合\cref{b4:thm:derived-ext-identification}即得矛盾。
\end{example}

\subsection{厚子范畴}
\label{b4:subsec:2002:02}

在有界有限生成投射复形中，有限直和、移位和锥并不能单独表达“取一个投射直和项”的操作。例如，从自由模 \(R^n\) 取出有限生成投射直和项，是完美对象形成过程的一部分，不能只允许取自由项。

\begin{definition}\label{b4:def:thick-subcategory}
设 \(\mathcal T\) 为三角范畴。一个含零对象、对同构、移位、逆移位和映射锥封闭的全子范畴 \(\mathcal S\)，若还满足：当 \(X\in\mathcal S\) 且 \(X\cong Y\oplus Z\) 在 \(\mathcal T\) 中成立时，\(Y,Z\in\mathcal S\)，便称为\term[houzifanchou]{厚子范畴}[thick subcategory]。包含对象族 \(\mathcal U\) 的最小厚子范畴记为 \(\operatorname{thick}(\mathcal U)\)。
\end{definition}
\symidx{\(\operatorname{thick}(T)\)：由对象生成的厚子范畴}

厚生成只允许有限次构造以及有限直和；没有规定可以取无限直和。“某对象属于 \(\operatorname{thick}(T)\)”因而是一条有限构造声明。

\begin{proposition}\label{b4:prop:perfect-built-from-ring}
设 \(R\) 为含幺环，\(P\) 为有界有限生成投射左 \(R\)-模复形。在 \(D(R\text{-}\mathrm{Mod})\) 中有 \(P\in\operatorname{thick}(R)\)，这里 \(R\) 作为集中在零次的左正则模。
\end{proposition}
\begin{proof}
每个 \(P^n\) 都是某个有限自由模 \(R^{r_n}\) 的直和项，故 \(P^n[-n]\in\operatorname{thick}(R)\)。按非零项的数目归纳。若最低非零次数为 \(a\)，取从 \(a+1\) 次开始保留原复形的子复形 \(P^{\ge a+1}_{\mathrm{br}}\)。有逐次分裂短正合列
\[
 0\longrightarrow P^{\ge a+1}_{\mathrm{br}}
 \longrightarrow P\longrightarrow P^a[-a]\longrightarrow0.
\]
它在同伦范畴、因而也在导出范畴中给出好三角。归纳假设处理左项，右项已属于所需厚子范畴，于是中项也属于其中。
\end{proof}

反向包含有一个实质细节：导出范畴中的直和分解，最初只给出同伦意义的幂等态射，未必已经是所选模型逐项分裂的幂等矩阵。下一小节的证明概要说明，怎样从紧性提取有限构造，再把这个直和项写成有界有限生成投射复形。

\subsection{紧性概念的适用背景}
\label{b4:subsec:2002:03}

紧性比较一个对象到任意直和的态射与它到各直和项的态射。因此必须先有任意集合指标的余积。这里采用 \(D(R\text{-}\mathrm{Mod})\)，不把讨论含混地放在可能没有无限直和的有界有限生成模范畴里。

\begin{definition}\label{b4:def:compact-object}
设 \(\mathcal T\) 为具有任意集合指标余积的局部小三角范畴，\(P\in\mathcal T\)。若对每族对象 \((X_i)_{i\in I}\)，自然映射
\[
 \bigoplus_{i\in I}\operatorname{Hom}_{\mathcal T}(P,X_i)
 \longrightarrow\operatorname{Hom}_{\mathcal T}\left(P,\bigoplus_{i\in I}X_i\right)
\]
为同构，则称 \(P\) 为\term[jinduixiang]{紧对象}[compact object]。若对象 \(G\in\mathcal T\) 满足：对任意 \(X\in\mathcal T\)，\(\operatorname{Hom}_{\mathcal T}(G,X[n])=0\) 对所有 \(n\in\mathbb Z\) 成立时必有 \(X=0\)，则称 \(G\) 为检测对象是否为零的生成元。
\end{definition}

模范畴的直和正合，因而复形的逐次直和保持拟同构，并给出模导出范畴的余积。后一断言还须检验态射的泛性质：对一族 \(X_i\)，用\cref{b4:thm:semifree-resolution}取半自由替换 \(P_i\to X_i\)。直和 \(P=\bigoplus_iP_i\) 仍为半自由复形，且 \(P\to\bigoplus_iX_i\) 拟同构。由\cref{b4:prop:k-resolution-comparison}，从 \(P\) 出发的导出态射可以在同伦范畴计算；而从直和出发的复形映射及其同伦都是逐分量给定的，故
\[
 \operatorname{Hom}_D\left(\bigoplus_iX_i,Y\right)
 \cong\prod_i\operatorname{Hom}_K(P_i,Y)
 \cong\prod_i\operatorname{Hom}_D(X_i,Y).
\]
这就是余积泛性质。DG 模的相应构造可对照\cite[Sections 3.1 and 4.1]{Keller1994DerivingDG}；两种论证都不能以任意逐项投射复形替代所需的无界替换。

\begin{proposition}\label{b4:prop:perfect-compact-ring-generator}
设 \(R\) 为含幺环。在 \(D(R\text{-}\mathrm{Mod})\) 中，每个完美复形都是紧对象。左正则模 \(R\) 还是检测对象是否为零的生成元。
\end{proposition}
\begin{proof}
取有界有限生成投射模型 \(P\)，令 \(S=\{\,n\in\mathbb Z\mid P^n\ne0\,\}\)，这是有限集合。固定 Hom 次数 \(r\)，有
\[
\begin{aligned}
 \underline{\operatorname{Hom}}_R\left(P,\bigoplus_iX_i\right)^r
 &=\prod_{n\in S}\operatorname{Hom}_R
       \left(P^n,\bigoplus_iX_i^{n+r}\right)\\
 &\cong\prod_{n\in S}\bigoplus_i\operatorname{Hom}_R(P^n,X_i^{n+r})\\
 &\cong\bigoplus_i\prod_{n\in S}\operatorname{Hom}_R(P^n,X_i^{n+r}).
\end{aligned}
\]
第一处同构使用 \(P^n\) 有限生成：有限个生成元的像总共只涉及有限个直和项。第二处使用 \(S\) 有限，使各次数所涉及的有限指标集可以合成一个有限集。这正是有界性在 Hom 总化中的作用；若直积指标无限，不能如此交换直积与直和。上述同构与 Hom 微分相容，因而自然映射
\[
 \bigoplus_i\underline{\operatorname{Hom}}_R(P,X_i)
 \longrightarrow
 \underline{\operatorname{Hom}}_R\left(P,\bigoplus_iX_i\right)
\]
是复形同构，对它取 \(H^0\) 仍得到同构。左边逐分量取圈与边界，便得到各 Hom 复形零次上同调的直和；有界投射模型的全忠实性把上述同构识别为紧性的映射。另一方面，\(\operatorname{Hom}_{D(R\text{-}\mathrm{Mod})}(R,X[n])=H^n(X)\)。若全部为零，\(X\) 零调，便在导出范畴中同构于零。
\end{proof}

\begin{theorem}[完美复形的内在刻画]\label{b4:thm:perfect-compact-thick}
设 \(R\) 为含幺环。在 \(D(R\text{-}\mathrm{Mod})\) 中，下列三类对象相同：完美复形、紧对象以及 \(\operatorname{thick}(R)\) 中的对象。特别地，完美复形组成厚子范畴。
\end{theorem}

\begin{proofsketch}
完美推出紧、完美属于 \(\operatorname{thick}(R)\) 已分别由\cref{b4:prop:perfect-compact-ring-generator,b4:prop:perfect-built-from-ring}证明。紧对象对移位、有限直和、锥和直和项封闭：对锥使用 Hom 长正合列和五引理，对直和项限制紧性的自然同构即可。因此 \(\operatorname{thick}(R)\) 中的对象都是紧对象。只剩紧推出完美。

设 \(C\) 紧，用\cref{b4:thm:semifree-resolution}取半自由消解 \(P\to C\)，其滤过 \(F_0P\subseteq F_1P\subseteq\cdots\) 的各次商都是若干移位 \(R[n]\) 的直和。下列望远镜三角表示 \(P\)：
\[
 \bigoplus_{j\ge0}F_jP\xrightarrow{\,1-s\,}
 \bigoplus_{j\ge0}F_jP\longrightarrow P
 \longrightarrow\left(\bigoplus_{j\ge0}F_jP\right)[1],
\]
其中 \(s\) 把第 \(j\) 项经滤过包含送到第 \(j+1\) 项。在每个复形次数上，\((1-s)(x)=0\) 的零号坐标给出 \(x_0=0\)，随后递推得到 \(x_1=x_2=\cdots=0\)，所以 \(1-s\) 单射。其余核正是按滤过包含取的余极限，即 \(P\)，因而这一逐次短正合列给出所写好三角。对它取 \(\operatorname{Hom}(C,-)\)，紧性把两端的 Hom 变成直和；相同的坐标递推说明，后一移位的 \(1-s\) 仍单射。Hom 长正合列遂得
\[
 \operatorname*{colim}_j\operatorname{Hom}_D(C,F_jP)
 \cong\operatorname{Hom}_D(C,P).
\]
令 \(u:C\xrightarrow{\sim}P\) 为消解拟同构在导出范畴中的逆。上述余极限公式给出某个有限阶段及 \(v:C\to F_jP\)，使自然映射 \(a:F_jP\to P\) 满足 \(av=u\)。于是
\[
 (u^{-1}a)v=1_C.
\]
这说明 \(C\) 是 \(F_jP\) 的导出直和项。这里仅有滤过步数有限，每一层的自由生成元仍可能有无限多个。

还需把各层的无限直和缩为有限直和，并保留连接态射。两层时可以把这一步具体写出。设
\[
 U\xrightarrow{i}E\xrightarrow{p}V\xrightarrow{\delta}U[1]
\]
为好三角，其中 \(U,V\) 均为移位 \(R[n]\) 的任意直和，给定 \(f:C\to E\)。紧性使 \(pf\) 经某个有限子直和 \(V_0\) 分解为 \(pf=j_Vg\)，其中 \(j_V:V_0\to V\) 为包含。\(V_0\) 本身完美、因而紧，所以 \(\delta j_V\) 又经 \(U[1]\) 中某个有限子直和 \(U_0[1]\) 分解：
\[
 \delta j_V=j_U[1]\delta_0,
 \qquad \delta_0:V_0\longrightarrow U_0[1].
\]
这里 \(j_U:U_0\to U\) 是自由直和的分裂包含。由 \(\delta pf=0\) 得 \(j_U[1]\delta_0g=0\)，再用其左逆才得到 \(\delta_0g=0\)。把 \(\delta_0\) 补成好三角
\[
 U_0\longrightarrow E_0\xrightarrow{p_0}V_0
 \xrightarrow{\delta_0}U_0[1].
\]
\(E_0\) 由有限个 \(R[n]\) 用一次锥构造得到。态射公理给出 \(q:E_0\to E\)，与两端的包含相容；\cref{b4:prop:triangle-hom-exact}给出 \(v_0:C\to E_0\)，使 \(p_0v_0=g\)。因此 \(p(f-qv_0)=0\)，再次用 Hom 正合性得到 \(t:C\to U\)，使 \(f-qv_0=it\)。最后用 \(C\) 紧，将 \(t\) 写成 \(t=j_1t_1\)，其中 \(j_1:U_1\to U\) 是另一个有限子直和的包含。这样
\[
 f=\begin{pmatrix}q&ij_1\end{pmatrix}
   \begin{pmatrix}v_0\\t_1\end{pmatrix}
 :C\longrightarrow E_0\oplus U_1\longrightarrow E
\]
就是所需的有限因子分解。除了压缩两端，还须控制提升的障碍 \(\delta_0g\)，以及提升以后与原映射之差；这两处不能只靠“取有限子直和”略去。

一般有限层数使用 Keller \cite[Section 5.3, Lemma]{Keller1994DerivingDG}所给的有限因子分解引理及其八面体归纳。该引理同时处理态射与锥：若 \(c:Z'\to Z\) 的锥由自由模移位的任意直和经有限次扩张得到，则从紧对象出发的 \(C\to Z\) 可补成从 \(b:C'\to C\) 到 \(c\) 的交换方块，使 \(b\) 的锥只由有限个自由模移位经有限次扩张得到。应用于 \(0\to F_jP\) 及上述 \(v\)，交换性给出 \(vb=0\)；再由 \(v\) 有左逆得 \(b=0\)。\cref{b4:prop:split-triangle-criterion}遂使 \(C\) 成为 \(\operatorname{Cone}(b)\) 的直和项，故 \(C\in\operatorname{thick}(R)\)。至此，已得到紧对象恰为 \(\operatorname{thick}(R)\) 中的对象；还没有把这些导出直和项写成有界有限生成投射模型。

要得到完美模型，还需处理同伦意义的幂等态射。若 \(e:R^m\to R^m\) 已是满足 \(e^2=e\) 的模同态，则 \(R^m=\operatorname{im}e\oplus\ker e\)，其像直接是有限生成投射模。可是在有界投射复形 \(Q\) 的模型上，导出直和项只给出 \([e]^2=[e]\)，即 \(e^2-e\) 零同伦；这不保证每个次数分量 \(e^n:Q^n\to Q^n\) 满足 \((e^n)^2=e^n\)，不能直接逐次取其像作为所需复形。

这里引用 Keller 对负 DG 范畴的加强结论在普通环上的版本：紧对象本身可以写成有限次扩张，且各扩张因子是有限个自由模移位之和的导出直和项\cite[Section 5.3, Remark b]{Keller1994DerivingDG}。原处说明此加强结论由 Rickard 的证明改编得到；普通环的紧对象与有界有限生成投射模型的对应也见\cite[Section 3.1]{Keller1998AbelianDerived}。普通环集中在零次，满足正次数态射复形为零的负 DG 条件。这里先将对象重新组织为有限扩张，再对每个扩张因子取直和项；不能从 \(C\) 是某个有限扩张的直和项，直接推断总对象的幂等态射保持原来的扩张滤过。

对上述任一因子 \(G\)，把它写成 \(V=\bigoplus_{j=1}^m R[n_j]\) 的导出直和项。不同移位之间没有非零态射，即
\[
\operatorname{Hom}_D(R[n],R[n'])=0\qquad(n\ne n').
\]
因此 \(V\) 上对应的幂等态射按相同移位分块。在每一块上，它由有限自由左 \(R\)-模的真正幂等自同态表示，其像是有限生成投射模。因此 \(G\) 同构于有限个有限生成投射模移位的直和。最后，有界投射模型的全忠实性使各次扩张的连接态射都能由复形映射表示，逐次取锥便得到 \(C\) 的有界有限生成投射模型。

以上概要省略了有限因子分解引理的一般八面体归纳，并引用了处理导出直和项的有限扩张结论。前者及紧对象论证见\cite[Sections 5.1--5.3, especially the lemma in 5.3]{Keller1994DerivingDG}；后者用来完成最后一项 \(\operatorname{thick}(R)\subseteq\operatorname{Perf}(R)\)，并非已经由逐项的投射模分解证明。
\end{proofsketch}

两个例子说明条件各有作用。若 \(R\ne0\)，无限自由模 \(F=\bigoplus_{j\ge0}R\) 不是紧对象：恒等映射 \(F\to\bigoplus_{j\ge0}R\) 不经过任何有限子直和。另一方面，\cref{b4:exm:bounded-not-perfect}中的单模有限生成且集中在一个次数，但其无限投射消解妨碍紧性。只检查“每项有限”或者“同调有界”，都不能替代完美性。

% !TeX root = ../../Book4.tex
% 纲要标题：### 21.3 倾斜对象
% 纲要标签：b4:sec:2003
% 纲要位置：计划纲要.md，第 3039 行。

\section{倾斜对象}
\label{b4:sec:2003}

用复形代替正则模时，既要能够生成全部完美对象，又要消去非零移位的自身态射。三顶点例子的自同态计算将这两个条件落实，并产生新的导出等价。


% 纲要内容（仅作写作计划，尚非教材正文）：
%
% * 消失条件与生成条件
% * 经典有限维代数例子
% * 导出 Morita 理论的基本版本
%

\subsection{消失条件与生成条件}
\label{b4:subsec:2003:01}

第二卷定理16.3.3已经证明，普通 Morita 理论用有限生成投射生成元代替正则模。导出范畴允许把替代对象选成复形。假如希望一个三角等价把环 \(B\) 的正则模送到 \(R\) 上的完美复形 \(T\)，并在完美对象上给出等价，那么正则模的两项性质也须由 \(T\) 保留：\(B\) 与自身非零移位之间没有态射，而每个完美 \(B\)-复形都能由 \(B\) 经过有限次移位、锥和取直和项构造。因此需要
\[
 \operatorname{Hom}(T,T[n])=0\quad(n\ne0),\qquad
 \operatorname{thick}(T)=\operatorname{Perf}(R).
\]
前者使自同态信息集中在零次，后者保证用 \(T\) 能恢复全部完美对象。如果高次自身态射仍然存在，就须保留完整的 DG 自同态代数。

\ProseParagraphBreak
两个简单例子说明这两项要求各有什么作用。若 \(R=k\times k\)，直和项 \(T=k\times0\) 没有非零高次自身态射，但由它构造的对象都被 \((0,1)\) 零化，因而无法得到 \(0\times k\)。若 \(R=k\)，\(T=k\oplus k[1]\) 包含 \(k\) 作为直和项，所以生成全部完美复形；但从 \(T\) 的 \(k[1]\) 分量到 \(T[1]\) 的 \(k[1]\) 分量可取恒等映射，故 \(\operatorname{Hom}(T,T[1])\ne0\)。

\begin{definition}\label{b4:def:tilting-complex}
设 \(k\) 为域，\(R\) 为含幺结合 \(k\)-代数。若 \(T\in\operatorname{Perf}(R)\) 满足
\[
 \operatorname{Hom}_{D(R\text{-}\mathrm{Mod})}(T,T[n])=0
 \quad(n\ne0),\qquad
 \operatorname{thick}(T)=\operatorname{Perf}(R),
\]
称 \(T\) 为 \(R\) 上的\term[qingxiefuxing]{倾斜复形}[tilting complex]，也称所给三角范畴中的倾斜对象。这里使用的是包括正负所有非零次数的消失条件。
\end{definition}

\(R\) 本身是倾斜复形：其自身高次态射为零，厚生成条件由\cref{b4:thm:perfect-compact-thick}给出。倾斜复形推广的正是正则模的这两项性质。

\ProseParagraphBreak
这里的倾斜对象与\chapref{b4:ch:19}的扭对倾斜使用不同数据。后者从一个心及其中的扭对构造另一个心；本节则要求一个完美对象具有自身态射消失和厚生成性质。构造出新心并不自动产生这样的对象，更不自动给出两个模导出范畴之间的等价。

\begin{proposition}\label{b4:prop:tilting-module-criterion}
设 \(k\) 为域，\(R\) 为 \(k\)-代数，\(T\) 为左 \(R\)-模。假设 \(T\) 有有限长有限生成投射消解，\(\operatorname{Ext}^n_R(T,T)=0\) 对所有 \(n>0\) 成立，并且存在正合列
\[
 0\longrightarrow R\longrightarrow T_0\longrightarrow T_1
 \longrightarrow\cdots\longrightarrow T_m\longrightarrow0,
\]
其中每个 \(T_i\) 是某个有限直和 \(T^{r_i}\) 的直和项。则 \(T\) 集中在零次是倾斜复形。
\end{proposition}
\begin{proof}
有限投射消解给出完美性。对 \(n>0\)，自身导出态射由\cref{b4:thm:derived-ext-identification}识别为所给 Ext 群；对 \(n<0\)，取位于非正次数的投射消解 \(P\to T\)，则 \(\underline{\operatorname{Hom}}_R(P,T)\) 没有负次项，故负次态射为零。

把所给长正合列按像分成短正合列。从末项起逆向使用相应好三角，可知 \(R\in\operatorname{thick}(T)\)。由\cref{b4:thm:perfect-compact-thick}，\(\operatorname{Perf}(R)=\operatorname{thick}(R)\subseteq\operatorname{thick}(T)\)。反向包含由完美子范畴的厚性和 \(T\) 完美得到。
\end{proof}

\subsection{经典有限维代数例子}
\label{b4:subsec:2003:02}

下面用三个向量空间及两个线性映射计算一个倾斜对象，暂不需要第五卷的一般箭图表示论。设 \(k\) 为域，\(R\) 是箭图 \(1\to2\to3\) 的路径代数：以三条长度零的路径、两条箭头及它们的复合为 \(k\)-基，乘法按可接续路径复合，不可接续时取零。具体记箭头为 \(\alpha:1\to2\)、\(\beta:2\to3\)，约定 \(\alpha=e_2\alpha e_1\)、\(\beta=e_3\beta e_2\)，所以 \(\beta\alpha\) 是从一到三的路径。有限维左 \(R\)-模等价于图式 \(V_1\to V_2\to V_3\)；三个顶点幂等元给出三个分量，箭头作用给出两个线性映射。模同态是使这两个方块交换的三元线性映射。

\ProseParagraphBreak
写出三个投射模以及中间顶点的简单模：
\[
\begin{aligned}
 P_1&=(k\xrightarrow{1}k\xrightarrow{1}k),&
 P_2&=(0\longrightarrow k\xrightarrow{1}k),\\
 P_3&=(0\longrightarrow0\longrightarrow k),&
 S_2&=(0\longrightarrow k\longrightarrow0).
\end{aligned}
\]
这里 \(P_i=Re_i\)，并且 \(\operatorname{Hom}_R(P_i,V)\cong V_i\)：态射由顶点 \(i\) 处的向量决定，其余分量由箭头推出。因此它们确为投射模，且 \(R=P_1\oplus P_2\oplus P_3\)。

\begin{example}\label{b4:exm:a3-tilting-object}
在上述路径代数 \(R\) 上，令 \(T=P_1\oplus P_2\oplus S_2\)。则 \(T\) 是倾斜复形，而
\[
 B=\operatorname{End}_R(T)^{\mathrm{op}}
\]
同构于路径代数 \(k(1\to2\leftarrow3)\)。
\end{example}
\begin{proof}
有短正合列
\begin{equation}\label{b4:eq:a3-simple-resolution}
 0\longrightarrow P_3\xrightarrow{i}P_2
 \xrightarrow{q}S_2\longrightarrow0.
\end{equation}
所以 \(T\) 有长度至多一的有限投射消解。对任意表示 \(V\)，这份消解给出
\[
 \operatorname{Ext}^1_R(S_2,V)
 =\operatorname{coker}(V_2\longrightarrow V_3),
 \qquad \operatorname{Ext}^{n}_R(S_2,V)=0\quad(n\ge2).
\]
对 \(V=P_1,P_2,S_2\)，\(V_2\to V_3\) 都满；因此 \(\operatorname{Ext}^n_R(T,T)=0\) 对所有 \(n>0\) 成立。负次态射为零同样由非正次投射消解计算。\eqref{b4:eq:a3-simple-resolution}的三角把 \(P_3\) 从 \(P_2,S_2\) 构造出来，故 \(R\in\operatorname{thick}(T)\)，由完美性即得厚生成。

自同态的计算也可以逐项完成。三个直和项的自同态环均为 \(k\)。不同项之间仅有
\[
 \operatorname{Hom}_R(P_2,P_1)=k j,
 \qquad \operatorname{Hom}_R(P_2,S_2)=k q
\]
非零：从 \(P_1\) 出发的态射由目标顶点一的分量决定，故另外两项为零；从 \(S_2\) 到 \(P_1\) 或 \(P_2\) 的态射则由第二条箭头的交换条件迫为零。相对于 \(T\) 的三个直和项，自同态代数因此为
\[
 E=\left\{\left.
 \begin{pmatrix}a&b&0\\0&c&0\\0&d&e\end{pmatrix}
 \ \right|\ a,b,c,d,e\in k\right\}.
\]
其中两条非恒等箭头从第二个直和项发出，彼此不能复合。取反代数后，两条箭头都指向第二个顶点，正好得到 \(B=k(1\to2\leftarrow3)\)。
\end{proof}

这个例子不是普通 Morita 等价换了一种写法。\(R\) 的三个不可分解投射模的合成长度为 \(3,2,1\)；\(B\) 的三个不可分解投射模分别支持在 \(1\to2\)、顶点二、\(3\to2\)，合成长度为 \(2,1,2\)。这些类型已经穷尽：两代数模去箭头生成的幂零理想均为 \(k^3\)，三个顶点给出全部简单模，第二卷定理15.3.10的简单模与不可分解投射模对应给出所述清单。模范畴等价保持简单对象、投射性、不可分解性和合成列，因而保持这组长度；两组数不同，所以两者的模范畴不等价。下一小节的定理却会给出其导出范畴的等价。

\subsection{导出 Morita 理论的基本版本}
\label{b4:subsec:2003:03}

令 \(T\) 为倾斜复形，选择有界有限生成投射模型 \(P\)，并令
\[
 E=\underline{\operatorname{End}}_R(P),\qquad
 B=\operatorname{End}_{D(R\text{-}\mathrm{Mod})}(T)^{\mathrm{op}}.
\]
倾斜条件恰好说 \(H^n(E)=0\) 对 \(n\ne0\) 成立，而 \(H^0(E)=B^{\mathrm{op}}\)。这一步把一个涉及所有移位的条件变为自同态 DG 代数的上同调消失。

\begin{proposition}\label{b4:prop:dg-degree-zero-formality}
设 \(k\) 为交换环，\(E\) 为 \(k\) 上的 DG 代数，且 \(H^n(E)=0\) 对所有 \(n\ne0\) 成立。令
\[
 (\tau_{\le0}E)^n=
 \begin{cases}
 E^n,&n<0,\\
 \ker(\mathrm{d}:E^0\to E^1),&n=0,\\
 0,&n>0.
 \end{cases}
\]
原微分和乘法使 \(\tau_{\le0}E\) 成为 DG 子代数，而且包含映射与零次取商给出 DG 代数拟同构
\[
 E\ \longleftarrow\ \tau_{\le0}E\ \longrightarrow\ H^0(E).
\]
\end{proposition}
\begin{proof}
非正次数在乘法下封闭；两零次余圈之积仍为余圈。来自 \(E^{-1}\) 的微分落在零次余圈中，故这是子复形。到 \(H^0(E)\) 的映射在负次取零，在零次把余圈取模余边。负次元素的乘积不会产生正次或零次项；两个零次余圈的乘法与取商相容，所以它是代数映射。包含映射在非正次上同调上为同构，而 \(E\) 正次上同调为零；取商映射在零次为同构，两边其他次数上同调均为零。因此两映射都是拟同构。
\end{proof}

注意这里得到的是通过第三个 DG 代数连接的两个拟同构，未必有直接代数映射 \(H^0(E)\to E\)。任选每个同伦类的复形映射代表，通常不能同时保持乘法。正是这一点阻止了把普通 Morita 理论中的论证逐字照搬。

\begin{theorem}[导出 Morita 定理的倾斜版本]\label{b4:thm:derived-morita-tilting}
设 \(k\) 为域，\(R\) 为含幺结合 \(k\)-代数，\(T\in\operatorname{Perf}(R)\) 为倾斜复形，并令
\[
 B=\operatorname{End}_{D(R\text{-}\mathrm{Mod})}(T)^{\mathrm{op}}.
\]
则存在 \(R\)-\(B\) 双模复形 \(X\)，作为左 \(R\)-模复形时在导出范畴中同构于 \(T\)，使
\[
 X\otimes_B^{\mathbf L}-:
 D(B\text{-}\mathrm{Mod})\ \longrightarrow\ D(R\text{-}\mathrm{Mod})
\]
成为 \(k\)-线性三角等价，且把 \(B\) 送到 \(T\)。此等价限制为 \(\operatorname{Perf}(B)\simeq\operatorname{Perf}(R)\)。
\end{theorem}

这一定理是\term[daochuMoritalilun]{导出 Morita 定理}[derived Morita theorem]的倾斜版本。下面通过 DG 自同态代数说明双模及等价的构造；所用文献主要采用右模，本节使用左模，因此自同态代数取反。

\begin{proofsketch}
取 \(T\) 的有界有限生成投射模型 \(P\)，记
\(E=\underline{\operatorname{End}}_R(P)\)、\(C=E^{\mathrm{op}}\)。这里取 DG 反代数：对次数分别为 \(a,b\) 的齐次自同态 \(e,f\)，
\[
e^{\mathrm{op}}f^{\mathrm{op}}=(-1)^{ab}(fe)^{\mathrm{op}},
\qquad \mathrm d_C(e^{\mathrm{op}})=(\mathrm d_Ee)^{\mathrm{op}}.
\]
在 \(P\) 上定义右作用
\(p\cdot e^{\mathrm{op}}=(-1)^{a|p|}e(p)\)。反乘法中的符号保证右作用的结合律；Hom 微分公式则给出
\[
\mathrm d_P(p\cdot e^{\mathrm{op}})
=(\mathrm d_Pp)\cdot e^{\mathrm{op}}
+(-1)^{|p|}p\cdot\mathrm d_C(e^{\mathrm{op}}).
\]
左 \(R\) 作用与这项右作用交换，所以 \(P\) 是严格的左 \(R\)、右 DG \(C\) 双模。

左 DG \(C\) 模是带分次左 \(C\) 作用的复形 \(N\)，其微分满足 \(\mathrm d(cn)=(\mathrm d_Cc)n+(-1)^{|c|}c\,\mathrm dn\)。以下 \(D(C)=D(C\text{-}\mathrm{DGMod})\) 表示左 DG \(C\) 模同伦范畴在拟同构处局部化所得的导出范畴。对左 \(R\)-模复形 \(Y\)，次数 \(b\) 的齐次映射 \(\varphi:P\to Y\) 上的左 \(C\) 作用为
\[
e^{\mathrm{op}}\cdot\varphi=(-1)^{ab}\varphi e.
\]
带符号的反乘法给出作用的结合律，而
\(\mathrm d(\varphi e)=(\mathrm d\varphi)e+(-1)^b\varphi(\mathrm d_Ee)\) 给出
\[
\mathrm d(e^{\mathrm{op}}\cdot\varphi)
=\mathrm d_C(e^{\mathrm{op}})\cdot\varphi
+(-1)^a e^{\mathrm{op}}\cdot\mathrm d\varphi.
\]
因此 Hom 复形确实是左 DG \(C\) 模。对这些 DG 模使用半自由消解，得到伴随
\[
 F=P\otimes_C^{\mathbf L}-:D(C)\rightleftarrows
 D(R\text{-}\mathrm{Mod}):U=\mathbf R\operatorname{Hom}_R(P,-).
\]
因为 \(P\) 有界且逐项有限生成投射，\(U\) 可用普通 Hom 复形计算，并保持任意余积；左伴随 \(F\) 也保持任意余积。复形级伴随的单位在齐次元上为
\(\eta_N(n)(p)=(-1)^{|n||p|}p\otimes n\)。在识别 \(P\otimes_CC\cong P\) 后，它在左正则 DG 模 \(C\) 上的公式是
\[
\eta_C:C\longrightarrow\underline{\operatorname{End}}_R(P),
\qquad
\eta_C(e^{\mathrm{op}})(p)
=(-1)^{a|p|}p\cdot e^{\mathrm{op}}=e(p).
\]
它与微分相容，并且
\[
\eta_C(e^{\mathrm{op}}f^{\mathrm{op}})
=(-1)^{ab}fe
=e^{\mathrm{op}}\cdot\eta_C(f^{\mathrm{op}}).
\]
故这是左 \(C\) 模复形的同构，目标使用刚才的左作用。它并不是把反代数与原自同态代数当作同一 DG 代数。

单位为同构的对象组成一个在同构下封闭的满子范畴。两函子均为三角函子，单位与平移相容，三角的 Hom 正合列给出该类对平移及好三角封闭；两函子保持任意余积，故该类也对任意余积封闭；加性与自然性又使单位在直和分解下分块，故该类对直和项封闭。它包含 \(C\)，而半自由消解把每个 DG 模写成自由移位的逐层锥及望远镜，所以它等于整个 \(D(C)\)。这证明 \(F\) 全忠实。

由伴随三角恒等式，余单位在 \(F(C)\cong P\) 上也是同构。余单位为同构的对象类同样在同构下封闭，并由上述三角、余积和分块论证分别对好三角、任意余积和直和项封闭。倾斜条件使
\(R\in\operatorname{thick}(P)\)，故该类包含 \(R\)。以
\(\operatorname{Loc}(R)\) 表示包含 \(R\)、对平移、好三角、任意余积及直和项封闭的最小满子范畴。普通环的半自由消解与望远镜构造给出
\[
\operatorname{Loc}(R)=D(R\text{-}\mathrm{Mod}).
\]
因此余单位处处为同构，\(F\) 是等价。这里先用有限构造得到 \(R\)，再用任意余积与望远镜恢复无界对象。

自身态射消失使 \(H^{\ne0}(C)=0\)、\(H^0(C)=B\)。由\cref{b4:prop:dg-degree-zero-formality}，令 \(A=\tau_{\le0}C\)，便有 DG 代数拟同构
\[
 C\longleftarrow A\longrightarrow B.
\]
每条拟同构都诱导导出模范畴的等价：在自由模上，延标量与限制标量的伴随单位是这条拟同构；再用半自由消解及余积、三角封闭性推广到全部对象。限制标量反映拟同构，伴随三角恒等式便使余单位也处处为同构。合成上述等价，便得到 \(D(B)\simeq D(R)\)。

严格双模也由这一步得到。把 \(P\) 限制为 \(R\)-\(A\) 双模，取它在 DG 代数 \(R\otimes_kA^{\mathrm{op}}\) 上的半自由替换 \(Q\to P\)，并令
\[
 X=Q\otimes_A B.
\]
设 \(\mathcal B_R\) 是 \(R\) 的 \(k\)-向量空间基。\(Q\) 的半自由滤过逐层在分次模意义下分裂，各滤过商是移位正则模
\(R\otimes_kA^{\mathrm{op}}\) 的直和。由于 \(R\) 集中在零次，限制为左 \(A^{\mathrm{op}}\) 模后有复形同构
\[
R\otimes_kA^{\mathrm{op}}
\cong\bigoplus_{r\in\mathcal B_R}A^{\mathrm{op}}.
\]
因此限制后的滤过商仍是移位自由 DG \(A^{\mathrm{op}}\) 模的直和；换成右 \(A\) 模即说明 \(Q\) 仍半自由，因而 K-平坦。将 \(Q\otimes_A-\) 作用于拟同构 \(A\to B\)，得到左 \(R\) 模复形的拟同构
\(Q\cong Q\otimes_AA\to Q\otimes_AB=X\)。于是
\(X\simeq P\simeq T\) 作为左 \(R\) 模成立。导出张量的结合律把合成等价识别为 \(X\otimes_B^{\mathbf L}-\)。它把 \(B\) 送到 \(T\)，从而把 \(\operatorname{thick}(B)\) 送到 \(\operatorname{thick}(T)\)；\cref{b4:thm:perfect-compact-thick}给出完美子范畴上的限制。

本概要省略了 DG 半自由消解的逐次添加生成元、双模替换的比较及导出伴随的链级相容性。消解与双模张量分别见\cite[Sections 3.1 and 6.1--6.3]{Keller1994DerivingDG}，紧生成下的等价判据及倾斜实现见\cite[Sections 4.2, 8.2 and 9.1--9.2]{Keller1994DerivingDG}；普通代数的同一定理见\cite[Section 4, Theorem (Rickard)]{Keller1998AbelianDerived}。这些构造把同伦意义的自同态作用换成严格双模作用，不能用逐项任取代表代替。
\end{proofsketch}

把定理用于\cref{b4:exm:a3-tilting-object}，便得到 \(k(1\to2\to3)\) 与 \(k(1\to2\leftarrow3)\) 的导出等价。虽然它们的模范畴不同，允许使用复形后，原来的投射生成元可以由 \(P_1\oplus P_2\oplus S_2\) 代替；\eqref{b4:eq:a3-simple-resolution}正是恢复缺少的投射模 \(P_3\) 的计算。

% !TeX root = ../../Book4.tex
% 纲要标题：### 21.4 DG 增强、高阶范畴与几何应用
% 纲要标签：b4:sec:2004
% 纲要位置：计划纲要.md，第 3806 行。

\section{DG 增强、高阶范畴与几何应用}
\label{b4:sec:2004}

同伦交换方块中的同伦会进入锥映射，导出张量也会保留普通张量看不见的同调。这些具体计算说明增强与更高相容结构所需记录的资料。


% 纲要内容（仅作写作计划，尚非教材正文）：
%
% * 增强的用途
% * 稳定高阶范畴的动机
% * 几何表示论与导出几何所需的附加语言
%
% ---
%

\subsection{增强的用途}
\label{b4:subsec:2004:01}

增强的用途可以从一个很小的计算看出。设 \(R\) 为含幺环，\(f:C\to D\)、\(f':C'\to D'\)、\(u:C\to C'\)、\(v:D\to D'\) 为复形映射。若方块严格交换，便有锥之间的对角映射 \((y,x)\mapsto(vy,ux)\)。若方块只在同伦范畴中交换，实际给出的条件则是存在次数 \(-1\) 的 \(h:C\to D'\)，使
\[
 vf-f'u=\mathrm{d}_{D'}h+h\mathrm{d}_C.
\]
此时锥之间的映射为
\begin{equation}\label{b4:eq:cone-map-with-homotopy}
 \operatorname{Cone}(f)\longrightarrow\operatorname{Cone}(f'),
 \qquad (y,x)\longmapsto(vy+hx,ux).
\end{equation}
把两边微分写成 \(\left(\begin{smallmatrix}\mathrm{d}_D&f\\0&-\mathrm{d}_C\end{smallmatrix}\right)\)，矩阵相乘后，右上角的交换条件恰好就是所给同伦等式。因此锥映射需要 \(h\) 这项具体资料；只知道 \([vf]=[f'u]\) 并没有指定它。

\ProseParagraphBreak
这种依赖确实可能改变所得映射。在域 \(k\) 上，取 \(C=C'=k[-1]\)、\(D=D'=k\)、\(f=f'=0\)、\(u=v=0\)。次数 \(-1\) 的 \(h:C\to D'\) 可以是任意标量。此时两个锥都是集中在零次的 \(k\oplus k\)，\eqref{b4:eq:cone-map-with-homotopy}为 \(\left(\begin{smallmatrix}0&h\\0&0\end{smallmatrix}\right)\)。不同标量给出不同同伦类，因为这些锥没有非零微分，也没有非零同伦。一个在同伦范畴里固定的零方块，仍然允许不同锥映射。

\ProseParagraphBreak
DG 增强保留整个态射复形，从而保留同伦所在的负一次以及同伦之间更高次的比较。它不会神奇地消除选择，而是把选择及其差别放在可以继续计算的对象中。对已有的三角等价，是否存在相容的 DG 提升是额外问题；本章只构造了具体模型中的增强，并使用特定双模复形给出的等价，没有断言任意三角等价都能如此提升。

\subsection{稳定高阶范畴的动机}
\label{b4:subsec:2004:02}

上一小节中，\(h\) 是两个映射之间的同伦；若两个同伦的差又是 \(\partial s\)，其中 \(s\) 的次数为 \(-2\)，便有同伦之间的比较。\cref{b4:prop:cone-higher-homotopy}已证明，这种比较给出相应锥映射之间的同伦。继续下去，会出现一串彼此相容的高阶资料。普通范畴只记录对象和态射；三角范畴额外记录移位与好三角，但并不逐层保留这些资料。

\ProseParagraphBreak
\term[wendinggaojiefanchou]{稳定高阶范畴}[stable infinity-category]把态射集合换为包含路径及高阶同伦的映射空间。\secref{b4:sec:F04}用单纯集和内角填充给出无穷范畴的一个精确模型，并定义映射空间。稳定性进一步要求零对象、有限极限与有限余极限存在，且拉回方块与推出方块相同；这些泛性质比较映射空间的等价，纤维与余纤维因而通过移位相联系。

\ProseParagraphBreak
对复形而言，映射纤维可写成 \(\operatorname{Cone}(f)[-1]\)，映射余纤维可写成 \(\operatorname{Cone}(f)\)。\chapref{b4:ch:17}证明的三角关系已经显示了它们与移位的联系；高阶结构还保存构造所需的相容同伦。\secref{b4:sec:B04}以复形的映射空间为例，比较模型结构与稳定无穷范畴所保留的数据。

\subsection{几何表示论与导出几何简介}
\label{b4:subsec:2004:03}

几何中的一个基本现象，是两个子空间的交会出现额外的同调。设 \(k\) 为域，\(A=k[x]\)、\(M=A/(x)\)。普通张量积只给出 \(M\otimes_A M=k\)。若先用自由消解
\[
 P=(A\xrightarrow{x}A),\qquad P^{-1}=P^0=A,
\]
再张量 \(M\)，得到
\[
 M\otimes_A^{\mathbf L}M\simeq(k\xrightarrow{0}k),
 \qquad H^{-1}=k,\quad H^0=k.
\]
这项负次同调还可直接由方程理想识别。记 \(I=(x)\)，对短正合列 \(0\rightarrow I\rightarrow A\rightarrow M\rightarrow0\) 与 \(M\) 作张量。由\cref{b4:thm:tor-long-exact}，\(A\) 自由使其一次 Tor 为零，故有正合列
\[
 0\longrightarrow\operatorname{Tor}_1^A(M,M)
 \longrightarrow I\otimes_A M\longrightarrow A\otimes_A M.
\]
最后一张映射为零，因为 \(I\) 的元素在 \(M=A/I\) 上作用为零。又 \(I\otimes_A(A/I)\cong I/I^2\)：映射 \(i\otimes\bar a\longmapsto ia+I^2\) 的逆为 \(i+I^2\longmapsto i\otimes\bar1\)，两张映射均由张量的平衡关系良定义。因此
\[
 H^{-1}(M\otimes_A^{\mathbf L}M)
 \cong\operatorname{Tor}_1^A(M,M)
 \cong I/I^2=(x)/(x^2)\cong k,
\]
末一同构把 \(x+(x^2)\) 送到 \(1\)。模 \(I/I^2\) 称为闭嵌入 \(\operatorname{Spec}(A/I)\rightarrow\operatorname{Spec}A\) 的\term[yufa-mo]{余法模}[conormal module]。它保留方程模掉方程的乘积后的类；本例的非零类 \(x+(x^2)\) 因而记录了一阶方程信息，普通交会的零次环 \(M\otimes_A M=k\) 看不见这项负次同调。

\ProseParagraphBreak
若保留\cref{b4:exm:dg-koszul-one-element}的乘法，这个计算还得到 \(k[\varepsilon]/(\varepsilon^2)\)、\(|\varepsilon|=-1\)、\(\mathrm{d}=0\) 的 DG 代数。导出几何把类似资料用于研究交、变形与模问题，而不只保留零次商环。

\ProseParagraphBreak
这个算例并不已经定义一般的导出概形。要从仿射计算推广到几何对象，还需要层与下降、局部模型之间的粘合，并说明采用哪一种增强；在正特征中，也不能未经比较就把交换 DG 代数当成所有导出交换环的通用模型。第三卷已经介绍普通层与概形，但一般导出粘合、余切复形及其形变意义不在本章建立。

\ProseParagraphBreak
几何表示论则在空间或概形上研究带群作用的层、它们的上同调以及卷积运算，再从这些结构构造和比较表示。这里不仅需要本卷的导出范畴，还需要具体的层理论、等变性与相应的几何方法。第五卷先承担有限群、箭图和 Lie 代数表示的代数基础；本节只是说明本卷的同调构造为何会在这些方向再次出现，不把这些方向列作本卷已经完成的理论。


\begin{chaptersummary}
DG 范畴把复形之间的态射组织成复形，零次余圈是复形映射，零次余边是零伦映射。微分与复合的带符号相容性使 \(H^0\) 成为普通范畴；要得到三角结构，还须有相容的移位与锥。增强记录的是这些更细的数据，而不只是为原来的三角范畴另换名称。

\ProseParagraphBreak
完美复形具有有界有限生成投射模型，既能通过有限次构造从正则模得到，又能用与任意直和交换的紧性刻画。倾斜对象进一步要求全部非零移位的自身态射消失。三顶点计算把这两个条件落实为一个短消解及一张自同态矩阵；导出 Morita 定理则把这样的计算转化为导出等价。紧对象刻画用望远镜与有限因子分解提取有限模型，Morita 等价经 DG 自同态代数和双模替换构造。锥映射中的同伦选择、导出张量中的额外同调，说明进一步研究这些等价时为何需要保留增强及其相容数据。
\end{chaptersummary}

\BookExercises

\begin{optionalexercise}\label{b4:ex:20-koszul-dependent-relations}
在 \(\mathbb Z\) 上取 DG 外代数 \(A=\bigwedge(\varepsilon_1,\varepsilon_2)\)，两生成元次数均为 \(-1\)，微分为 \(\mathrm{d}\varepsilon_1=2\)、\(\mathrm{d}\varepsilon_2=4\)。计算全部上同调，并解释为何第二个关系没有产生一个正则序列的消解。
\end{optionalexercise}
\begin{solution}
底层复形为 \(\mathbb Z\to\mathbb Z^2\to\mathbb Z\)，次数为 \(-2,-1,0\)，微分分别为 \(t\mapsto(-4t,2t)\) 和 \((a,b)\mapsto2a+4b\)。第一映射单射；第二映射的核由 \((-2,1)\) 生成，而第一映射的像由其两倍生成。故 \(H^{-2}=0\)、\(H^{-1}=\mathbb Z/2\mathbb Z\)、\(H^0=\mathbb Z/2\mathbb Z\)，其他次数为零。模去第一个关系 \(2\) 后，第二个元素 \(4\) 为零，不是非零商环上的非零因子，所以 Koszul 复形留下负次同调。
\end{solution}

\begin{optionalexercise}\label{b4:ex:20-unit-boundary-contractibility}
设 \(A\) 为交换环 \(k\) 上的 DG 代数，存在 \(h\in A^{-1}\) 使 \(\mathrm{d}h=1\)。设左 DG \(A\)-模 \(M\) 的作用满足 \(\mathrm{d}(am)=\mathrm{d}a\,m+(-1)^{|a|}a\,\mathrm{d}m\)。证明 \(M\) 的底层 \(k\)-复形可缩。这个缩同伦一定是 \(A\)-线性映射吗？
\end{optionalexercise}
\begin{solution}
令 \(s(m)=hm\)，它是次数 \(-1\) 的 \(k\)-线性映射。由作用的微分公式，\(\mathrm{d}(hm)+h\,\mathrm{d}m=(\mathrm{d}h)m=m\)，所以 \(\mathrm{d}s+s\mathrm{d}=1\)。一般不能声称它为齐次 \(A\)-线性映射：次数 \(-1\) 的左模态射应满足 \(s(am)=(-1)^{|a|}a s(m)\)，这要求 \(ha=(-1)^{|a|}ah\)，而一般 DG 代数并不分次交换。因此底层可缩不等于已经给出 DG 模意义的缩同伦。
\end{solution}

\begin{optionalexercise}\label{b4:ex:20-endomorphism-two-term}
设整数 \(m\ge2\)，\(P=(\mathbb Z\xrightarrow{m}\mathbb Z)\) 位于次数 \(-1,0\)。计算 DG 代数 \(E=\underline{\operatorname{End}}_{\mathbb Z}(P)\) 的全部上同调与正次数乘法，说明为何完美对象未必满足倾斜的消失条件。
\end{optionalexercise}
\begin{solution}
只有 \(E^{-1}=\mathbb Z\)、\(E^0=\mathbb Z^2\)、\(E^1=\mathbb Z\) 非零。把零次元素写成分别作用在 \(P^{-1},P^0\) 上的 \((a,b)\)，微分为 \(h\mapsto(mh,mh)\)、\((a,b)\mapsto m(a-b)\)。故 \(H^{-1}=0\)、\(H^0=\mathbb Z/m\mathbb Z\)、\(H^1=\mathbb Z/m\mathbb Z\)，其余为零。零次类按通常标量作用于一次类，一次类乘积因二次项为零而为零。虽然 \(P\) 有界有限自由，其到一次移位的导出态射仍非零。
\end{solution}

\begin{optionalexercise}\label{b4:ex:20-hzero-not-quasi-equivalence}
设 \(k\) 为域，把 \(k\) 与 \(k[u]\)、\(|u|=2\)、\(\mathrm{d}=0\) 分别视为只有一个对象的 DG 范畴。证明包含映射在 \(H^0\) 上给出等价，但不是拟等价。说明这个例子是否已经是两个三角增强之间的反例。
\end{optionalexercise}
\begin{solution}
两者的零次上同调代数均为 \(k\)，对象集合也都是单点，所以 \(H^0\) 上为等价。但第二个态射复形的二次上同调含有非零类 \(u\)，第一个没有，因此 Hom 复形映射不是拟同构。两个单对象 DG 范畴没有包含所需的移位和锥，不是这里的预三角增强；此例只说明任意 DG 范畴的 \(H^0\) 丢失了资料，不能据此推出增强间的更强结论。
\end{solution}

\begin{optionalexercise}\label{b4:ex:20-perfect-nonprojective-summand}
设 \(R=\mathbb Z\)，\(P=(R\xrightarrow{6}R)\) 位于次数 \(-1,0\)。利用中国剩余定理，在导出范畴中把 \(P\) 分成两个非零完美直和项，并为每一项写出有限自由模型。
\end{optionalexercise}
\begin{solution}
\(P\) 拟同构于零次的 \(\mathbb Z/6\mathbb Z\)，中国剩余定理给出 \(\mathbb Z/6\mathbb Z\cong\mathbb Z/2\mathbb Z\oplus\mathbb Z/3\mathbb Z\)。两项的模型分别为 \((\mathbb Z\xrightarrow{2}\mathbb Z)\) 和 \((\mathbb Z\xrightarrow{3}\mathbb Z)\)，均置于次数 \(-1,0\)。由这些拟同构在导出范畴中可逆，得到所需分解。不能把这个分解理解成原来每项 \(\mathbb Z\) 都分为两个非零直和项。
\end{solution}

\begin{optionalexercise}\label{b4:ex:20-field-perfect}
设 \(k\) 为域，\(X\) 为任意 \(k\)-向量空间上链复形。证明 \(X\) 同伦等价于零微分复形 \(\bigoplus_n H^n(X)[-n]\)。再证明 \(X\) 完美当且仅当 \(\sum_n\dim_kH^n(X)<\infty\)。
\end{optionalexercise}
\begin{solution}
在每次取补空间 \(Z^n=B^n\oplus H_n\)、\(X^n=Z^n\oplus L_n\)。微分限制给出同构 \(L_n\to B^{n+1}\)：它的核为零，像为全部边界。于是 \(X\) 分为零微分的 \(H_n\) 与各两项恒等复形的直和；后一部分用上述同构的逆得到缩同伦。每个次数至多涉及相邻的两个可缩项，所以这个定义适用于无界复形。若总同调维数有限，则零微分部分为有界有限维复形，因而完美。反之，有界有限维模型的同调只能在有限个次数非零，且每次有限维，给出有限总维数。
\end{solution}

\begin{optionalexercise}\label{b4:ex:20-thick-product-ring}
设 \(k\) 为域，\(R=k\times k\)，\(e=(1,0)\)。在左 \(R\)-模导出范畴中，准确描述 \(\operatorname{thick}(Re)\) 中的对象，并证明 \(Re\) 不是 \(\operatorname{Perf}(R)\) 的倾斜生成元。
\end{optionalexercise}
\begin{solution}
每个 \(R\)-模及复形都由两个幂等元分解为一对 \(k\)-向量空间模及复形；映射、锥和同调均按分量计算。\(Re\) 对应 \((k,0)\)。因此它的厚子范畴所有对象第二分量为零，第一分量完美。反过来第一分量的每个完美 \(k\)-复形由有限直和及移位的 \(k\) 构成，故这些对象全在其中。\(R(1-e)=(0,k)\) 不属于该子范畴，因此生成条件失败，尽管 \(Re\) 自身的非零移位态射全为零。
\end{solution}

\begin{optionalexercise}\label{b4:ex:20-tilting-shifts-over-field}
设 \(k\) 为域，\(T\in\operatorname{Perf}(k)\) 非零。证明 \(T\) 是倾斜复形当且仅当存在整数 \(a\) 和正整数 \(r\)，使 \(T\cong k^r[a]\)。计算其反自同态代数。
\end{optionalexercise}
\begin{solution}
由向量空间复形的分裂，\(T\) 是各次同调的有限直和。若两个不同次数 \(p<q\) 的同调非零，则从 \(H^p(T)[-p]\) 到 \(H^q(T)[-q][q-p]\) 存在非零线性映射，得到 \(T\to T[q-p]\) 的非零态射，违背消失条件。故同调只能集中在一次，此时为 \(k^r[a]\)。反之这种对象没有非零移位的自身态射，而 \(k[a]\) 是其直和项，故厚生成所有完美复形。自同态环为 \(M_r(k)\)，取反后仍借矩阵转置同构于 \(M_r(k)\)。
\end{solution}

\begin{optionalexercise}\label{b4:ex:20-a3-generation-sequence}
沿用\cref{b4:exm:a3-tilting-object}的记号，构造一个符合\cref{b4:prop:tilting-module-criterion}的正合列 \(0\to R\to T_0\to T_1\to0\)。
\end{optionalexercise}
\begin{solution}
取 \(T_0=P_1\oplus P_2\oplus P_2\)、\(T_1=S_2\)。从 \(R=P_1\oplus P_2\oplus P_3\) 到 \(T_0\) 的映射为前两分量恒等、第三分量为 \(i:P_3\to P_2\)；从 \(T_0\) 到 \(S_2\) 取第三分量的商映射 \(q\)。核、像关系由\eqref{b4:eq:a3-simple-resolution}得到。\(T_0,T_1\) 均是有限个 \(T\) 的直和项，因此这把正文的厚生成计算写成该判据要求的正合列。
\end{solution}

\begin{optionalexercise}\label{b4:ex:20-a3-replace-wrong-simple}
在 \(R=k(1\to2\to3)\) 上令 \(S_1=(k\to0\to0)\)，考虑 \(U=P_1\oplus P_2\oplus S_1\)。证明 \(U\) 不满足倾斜消失条件，指出一个非零 Ext 群。
\end{optionalexercise}
\begin{solution}
有 \(0\to P_2\to P_1\to S_1\to0\)。施加 \(\operatorname{Hom}_R(-,P_2)\)，得到 \(\operatorname{Ext}^1_R(S_1,P_2)=\operatorname{coker}(\operatorname{Hom}(P_1,P_2)\to\operatorname{End}(P_2))\)。前一空间为 \((P_2)_1=0\)，后一空间为 \(k\)，所以该 Ext 群为 \(k\)。它是 \(\operatorname{Ext}^1_R(U,U)\) 的一个直和项，因此消失条件失败。
\end{solution}

\begin{optionalexercise}\label{b4:ex:20-formality-is-zigzag}
设 \(k\) 为域，\(E\) 为 DG 代数，\(H^n(E)=0\) 对 \(n\ne0\) 成立。说明“为每个零次上同调类任取一个代表”为何不自动构成代数映射 \(H^0(E)\to E\)。构造一个保持单位的 \(k\)-线性代表选择，要求它仍然破坏乘法。
\end{optionalexercise}
\begin{solution}
代表的选择首先只是一条到零次余圈的集合截面；即使取保持单位的 \(k\)-线性截面，也没有保证 \(s(ab)=s(a)s(b)\)。取交换 DG 代数 \(E=k[z,t]\otimes_k\bigwedge(\varepsilon)\)，\(|z|=|t|=0\)、\(|\varepsilon|=-1\)、\(\mathrm{d}\varepsilon=t\)、\(\mathrm{d}z=\mathrm{d}t=0\)。乘 \(t\) 在 \(k[z,t]\) 上单射，所以其上同调仅为零次的 \(k[z]\)。令 \(c_2(f)\) 为多项式 \(f(z)\) 的二次项系数，取 \(s(f)=f(z)+c_2(f)t\)。它保持单位且 \(k\)-线性，模去 \(t\) 后为恒等，确为代表选择；但 \(s(z)^2=z^2\ne z^2+t=s(z^2)\)。正文的截断拟同构避免了依赖这种不相容的选择。
\end{solution}

\begin{optionalexercise}\label{b4:ex:20-cone-homotopy-change}
沿用\eqref{b4:eq:cone-map-with-homotopy}的记号。若两种方块同伦 \(h,h'\) 满足 \(h-h'=\mathrm{d}_{D'}s-s\mathrm{d}_C\)，其中 \(s:C\to D'\) 的次数为 \(-2\)，证明它们给出的两个锥映射同伦。
\end{optionalexercise}
\begin{solution}
在锥上取次数 \(-1\) 的映射 \(S(y,x)=(sx,0)\)。锥微分给出
\[
 \mathrm{d}_{\operatorname{Cone}(f')}S+S \mathrm{d}_{\operatorname{Cone}(f)}
 =\begin{pmatrix}0&\mathrm{d}_{D'}s-s\mathrm{d}_C\\0&0\end{pmatrix}.
\]
这恰为两个锥映射的差。因此二次同伦控制一次同伦选择所导致的锥映射差别。
\end{solution}

\begin{optionalexercise}\label{b4:ex:20-derived-intersection-transverse}
设 \(k\) 为域，\(A=k[x,y]\)，比较 \(A/(x)\otimes_A^{\mathbf L}A/(y)\) 与 \(A/(x)\otimes_A^{\mathbf L}A/(x)\) 的上同调。说明额外同调由哪个乘法映射的核产生。
\end{optionalexercise}
\begin{solution}
用 \((A\xrightarrow{x}A)\) 消解 \(A/(x)\)。与 \(A/(y)=k[x]\) 张量后，微分仍为单射乘 \(x\)，故仅零次同调 \(k\) 非零。与 \(A/(x)=k[y]\) 张量后，微分为零，故零次与负一次同调均为 \(k[y]\)。额外负一次同调来自乘 \(x\) 在另一个商环中出现的核：前一商环中核为零，后一商环中为整个商环。
\end{solution}

